\documentclass[index,biblatex]{MUMAv-UPM}

%%%%%%%%%%%%%%%%%%%%%%%%%%%%%%%%%%%%%%%%%%%%%%%%%%%%%%%%%%%%%%%%%%%%%%%%
% PAQUETES ADICIONALES DEL TFM
%%%%%%%%%%%%%%%%%%%%%%%%%%%%%%%%%%%%%%%%%%%%%%%%%%%%%%%%%%%%%%%%%%%%%%%%
\usepackage{xcolor}
\usepackage[hidelinks]{hyperref}

\hypersetup{
  pdftitle={Análisis Espectral y Matricial del Operador de Multiplicación en Espacios de Sobolev Discretos},
  pdfauthor={Gabriel Suárez},
  pdfstartview={FitH}
}

%%%%%%%%%%%%%%%%%%%%%%%%%%%%%%%%%%%%%%%%%%%%%%%%%%%%%%%%%%%%%%%%%%%%%%%%
% COMANDOS PERSONALIZADOS DEL TFM
%%%%%%%%%%%%%%%%%%%%%%%%%%%%%%%%%%%%%%%%%%%%%%%%%%%%%%%%%%%%%%%%%%%%%%%%
\newcommand{\C}{\mathbb{C}}
\newcommand{\Poly}{\mathbb{P}[z]}
\newcommand{\norm}[1]{\left\| #1 \right\|_S}
\newcommand{\abs}[1]{\left| #1 \right|}
\newcommand{\ip}[2]{\langle #1, #2 \rangle_S}
\newcommand{\D}{\mathcal{D}}

%%%%%%%%%%%%%%%%%%%%%%%%%%%%%%%%%%%%%%%%%%%%%%%%%%%%%%%%%%%%%%%%%%%%%%%%
% ENTORNOS MATEMÁTICOS COMPATIBLES CON TFM.TEX
% (La plantilla usa thm, prop, lem, defn, exm, rmk; TFM.tex usa nombres en inglés)
%%%%%%%%%%%%%%%%%%%%%%%%%%%%%%%%%%%%%%%%%%%%%%%%%%%%%%%%%%%%%%%%%%%%%%%%
\theoremstyle{theoremMUMAv}
\newtheorem{theorem}[thm]{Teorema}
\newtheorem{corollary}[thm]{Corolario}
\newtheorem{proposition}[thm]{Proposición}
\newtheorem{lemma}[thm]{Lema}
\theoremstyle{definitionMUMAv}
\newtheorem{definition}[thm]{Definición}
\newtheorem{example}[thm]{Ejemplo}
\theoremstyle{remarkMUMAv}
\newtheorem{remark}[thm]{Observación}

%%%%%%%%%%%%%%%%%%%%%%%%%%%%%%%%%%%%%%%%%%%%%%%%%%%%%%%%%%%%%%%%%%%%%%%%
% TITULO, AUTORES Y DIRECTORES
%
\title{Análisis Espectral y Matricial del Operador de Multiplicación en Espacios de Sobolev Discretos}
\author{Gabriel Suárez}
\genderauthor{M}% Male or Female?
\bibauthor{Suárez, Gabriel}
\tutor{Tutor 1\\
Tutor 2}
% \tutor{Sofía Kobaleskáya}
%\bibtutor{Kobaleskáya, Sofía \and Hilbert, David}
\fecha{10 de junio de 2026}
%\fecha{10th June 2025}

%%%%%%%%%%%%%%%%%%%%%%%%%%%%%%%%%%%%%%%%%%%%%%%%%%%%%%%%%%%%%%%%%%%%%%%%
% RESUMEN Y ABSTRACT
%
\abstract{spanish}{
    Se estudia el operador de multiplicación por la variable compleja en espacios de polinomios dotados de un producto interno de Sobolev discreto, cuya componente continua es la medida de Lebesgue normalizada sobre una circunferencia de radio $R$. Se introduce la noción de punto $\delta$-singular ---un punto $c$ donde la evaluación puntual $\delta_c$ restringida al núcleo de la evaluación de la derivada no es acotada respecto a la norma de Sobolev--- y se demuestra que el número de puntos $\delta$-singulares determina exactamente la codimensión algebraica del subespacio donde el operador se estabiliza. Se obtienen cotas precisas para la localización de los ceros mediante la cadena $Q_m \to$ valores singulares $\to$ desigualdad de Weyl-Horn $\to$ atracción asintótica, y se establece que $Q_m(\D)$ desempeña en el régimen no acotado el papel que $\|\D\|$ juega en el caso clásico acotado. Para la medida de Lebesgue en circunferencia, se demuestra la equivalencia $\delta$-singular $\iff |c|\ge R \iff$ no-BPE.
}
\keywords{spanish}{
    Polinomios ortogonales de Sobolev;
    Operador de multiplicación;
    Análisis espectral;
    Valores singulares;
    Localización de ceros;
    Puntos $\delta$-singulares
}

\abstract{english}{
    \noindent This work studies the multiplication operator by the complex variable in spaces of polynomials endowed with a discrete Sobolev inner product, whose continuous component is the normalized Lebesgue measure on a circle of radius $R$. The notion of $\delta$-singular point is introduced ---a point $c$ where the point evaluation $\delta_c$ restricted to the kernel of the derivative evaluation is not bounded with respect to the Sobolev norm--- and it is proved that the number of $\delta$-singular points determines exactly the algebraic codimension of the subspace where the operator stabilizes. Precise bounds for the localization of zeros are obtained via the chain $Q_m \to$ singular values $\to$ Weyl-Horn inequality $\to$ asymptotic attraction, and it is established that $Q_m(\D)$ plays in the unbounded regime the same role that $\|\D\|$ plays in the classical bounded case. For the Lebesgue measure on the circle, the equivalence $\delta$-singular $\iff |c|\ge R \iff$ non-BPE is proved.
}
\keywords{english}{
    Sobolev orthogonal polynomials;
    Multiplication operator;
    Spectral analysis;
    Singular values;
    Zero location;
    $\delta$-singular points
}

\acknowledgements{
     Aquí van los agradecimientos.
}

%%%%%%%%%%%%%%%%%%%%%%%%%%%%%%%%%%%%%%%%%%%%%%%%%%%%%%%%%%%%%%%%%%%%%%%%
% DOCUMENTO
%

\begin{document}

\chapter{Introducción y Preliminares}\label{sec:intro}

\subsection{Contexto y Motivación}

La teoría de Polinomios Ortogonales (PO) es una rama clásica del análisis matemático con profundas aplicaciones en teoría de la aproximación, física matemática y métodos numéricos. Clásicamente, la teoría estándar consolidada por autores como Szegő estudia medidas de ortogonalidad soportadas en la recta real o en la circunferencia unidad. En estos contextos, el operador de multiplicación por la variable independiente (o la matriz de Jacobi asociada) suele ser un operador acotado definidos sobre un espacio de Hilbert estándar $L^2(\mu)$.

Sin embargo, en las últimas décadas ha surgido un gran interés en el estudio de los llamados \textbf{Polinomios Ortogonales de Sobolev}. A diferencia del caso estándar, estos polinomios son ortogonales respecto a un producto interno que involucra no solo evaluaciones de las funciones, sino también evaluaciones de sus derivadas. Formalmente, si $\mu$ es una medida positiva, el producto interno de Sobolev toma la forma general:
\begin{equation}
	\langle p, q \rangle_S = \langle p, q \rangle_{L^2(\mu)} + \sum_{k=1}^N \lambda_k p'(c_k) \overline{q'(c_k)},
\end{equation}
donde $\{c_k\}$ son puntos discretos en el plano complejo y $\lambda_k > 0$ son pesos.

La introducción de derivadas en la métrica cambia radicalmente la geometría del espacio de Hilbert subyacente. Una de las consecuencias más notables, y que constituye el núcleo de este trabajo, es el comportamiento del \textbf{operador multiplicación}. Mientras que en la teoría clásica este operador es acotado y, según el contexto, autoadjunto o unitario, en los espacios de Sobolev discretos su comportamiento depende críticamente de la ubicación de los puntos $c_k$ respecto al soporte de la medida $\mu$.

El objetivo principal es analizar la estructura espectral del operador multiplicación en el caso \textit{no acotado}. Específicamente, estudiamos la situación inestable donde los puntos $c_k$ se encuentran fuera del soporte esencial de la medida (el caso que denominaremos \textit{no-BPE} o de evaluaciones puntuales no acotadas). Para ello, se emplea la teoría de perturbaciones de rango finito y la teoría de los números-$s$ (valores singulares), con el fin de cuantificar la no-acotación del operador y su impacto en la asintótica de los ceros.

Este enfoque no solo proporciona un marco riguroso para los fenómenos observados anteriormente, sino que conecta el problema de los polinomios ortogonales con problemas modernos de análisis espectral y teoría de la aproximación en espacios de codimensión topológica finita.

\subsection{Medidas de Borel, Regularidad y Medida de Lebesgue}

Antes de introducir la familia concreta de productos internos de Sobolev que estudiaremos, conviene fijar el marco general de medidas en el que estos aparecen de manera natural. En la formulación clásica, un producto interno de Sobolev se construye a partir de un vector finito de medidas de Borel positivas y finitas, que ponderan separadamente las distintas derivadas de los polinomios.

\begin{definition}[Medida]
	Sea $X$ un conjunto y sea $\mathcal{A}$ una $\sigma$-álgebra de subconjuntos de $X$. Una aplicación
	\[
		\mu : \mathcal{A} \to [0,+\infty]
	\]
	se llama \emph{medida} si satisface las siguientes propiedades \cite{folland1999real,rudin1987real}:
	\begin{itemize}
		\item $\mu(\varnothing)=0$.
		\item Si $\{E_n\}_{n\geq 1} \subset \mathcal{A}$ es una familia numerable de conjuntos dos a dos disjuntos, entonces
		      \[
			      \mu\Bigl(\bigcup_{n=1}^{\infty} E_n\Bigr)=\sum_{n=1}^{\infty}\mu(E_n).
		      \]
	\end{itemize}
	Diremos que $\mu$ es \emph{finita} si $\mu(X)<\infty$, y \emph{positiva} si toma únicamente valores no negativos, como ocurrirá en todo este trabajo.
\end{definition}

\begin{definition}[Medida de Borel]
	Sea $X$ un espacio topológico. La $\sigma$-álgebra generada por los abiertos de $X$ se denomina $\sigma$-álgebra de Borel y se denota por $\mathcal{B}(X)$. Una \emph{medida de Borel} en $X$ es una medida definida sobre $\mathcal{B}(X)$ \cite{folland1999real,bogachev2007measure}. En particular, cuando $X=\mathbb{R}$ o $X=\mathbb{C}\cong\mathbb{R}^2$, las medidas de Borel permiten integrar funciones continuas, polinomios y sus derivadas sobre conjuntos geométricamente naturales como intervalos, discos o circunferencias.
\end{definition}

\begin{definition}[Soporte de una medida]
	Sea $\mu$ una medida de Borel en un espacio topológico $X$. El \emph{soporte} de $\mu$, denotado por $\operatorname{supp}(\mu)$, es el menor cerrado $F \subset X$ tal que $\mu(X\setminus F)=0$. Equivalentemente, $\operatorname{supp}(\mu)$ es el conjunto de puntos $x \in X$ tales que toda vecindad abierta de $x$ tiene medida estrictamente positiva \cite{folland1999real,bogachev2007measure}.
\end{definition}

Las propiedades del soporte son esenciales en este trabajo por dos motivos. Por un lado, la parte continua del producto interno queda determinada por la geometría de $\operatorname{supp}(\mu)$; por otro, la localización de las masas discretas respecto a dicho soporte gobierna la aparición o no de evaluaciones puntuales acotadas.

\begin{definition}[Medida de Lebesgue]
	La \emph{medida de Lebesgue} en $\mathbb{R}$, y más generalmente en $\mathbb{R}^n$, es la medida de Borel que extiende la noción usual de longitud, área o volumen \cite{folland1999real,rudin1987real}. Mediante la identificación $\mathbb{C} \cong \mathbb{R}^2$, hablaremos también de medida de Lebesgue en el plano complejo. Esta medida es completa, invariante por traslaciones y regular; además, los conjuntos suaves de dimensión uno, como las circunferencias, inducen de manera natural medidas de longitud de arco que son absolutamente continuas respecto a la parametrización angular.
\end{definition}

En particular, cuando el soporte continuo del problema es una circunferencia, la medida de referencia adecuada no es el área bidimensional del plano, sino la medida de longitud de arco normalizada sobre dicha curva. Esta elección conserva la estructura analítica unidimensional compatible con la teoría de Hardy y con las fórmulas integrales de Cauchy \cite{duren1970theory,szego1975orthogonal}.

\begin{remark}[Propiedades relevantes para el producto de Sobolev]
	Las medidas que intervienen en este trabajo serán siempre medidas de Borel positivas y finitas, con soporte compacto. Estas hipótesis garantizan que:
	\begin{itemize}
		\item todos los momentos de los polinomios que aparecen en el producto interno están bien definidos;
		\item el producto $L^2(\mu)$ asociado a la componente continua tiene sentido y genera una estructura de pre-Hilbert natural sobre $\Poly$;
		\item la posición relativa entre el soporte continuo y los puntos discretos $c_k$ puede analizarse mediante herramientas geométricas y de teoría del potencial \cite{stahl1992general,szego1975orthogonal}.
	\end{itemize}
\end{remark}

\begin{definition}[Medida de Lebesgue normalizada en $\mathbb{S}_r(a)$]
	Denotamos por $\mathbf{m}_{a;r}$ la medida de Lebesgue normalizada en la circunferencia
	$$\mathbb{S}_r(a) = \{z \in \C : \abs{z-a} = r\},$$
	es decir,
	\begin{equation}
		d\mathbf{m}_{a;r}(z) = \frac{\abs{dz}}{2\pi r}.
	\end{equation}
	En el caso particular de la circunferencia unidad centrada en el origen escribiremos simplemente $\mathbf{m} = \mathbf{m}_{0;1}$.
\end{definition}

\begin{definition}[Descomposición de Lebesgue en la circunferencia]
	Sea $\nu$ una medida de Borel finita sobre $\mathbb{S}_r(a)$. Entonces existe una descomposición única
	\begin{equation}
		d\nu = w\, d\mathbf{m}_{a;r} + d\nu_s,
	\end{equation}
	donde $w \in L^1(\mathbf{m}_{a;r})$ es la densidad de la parte absolutamente continua de $\nu$ respecto a la medida de Lebesgue normalizada, y $\nu_s$ es la parte singular. En particular, las masas discretas o átomos quedan contenidas en la componente singular.
\end{definition}

\begin{definition}[Condición de Szegő]
	Sea $\nu$ una medida de Borel finita en la circunferencia unidad, con descomposición
	\begin{equation}
		d\nu(e^{i\theta}) = w(\theta)\, d\mathbf{m}(e^{i\theta}) + d\nu_s(e^{i\theta}).
	\end{equation}
	Diremos que $\nu$ satisface la condición de Szegő si
	\begin{equation}
		\int_0^{2\pi} \log w(\theta)\, d\theta > -\infty.
	\end{equation}
\end{definition}

\begin{remark}[El caso de la medida de Lebesgue]
	Si $\nu = c\,\mathbf{m}$ con $c>0$, entonces $w(\theta)=c$ casi en todas partes y la condición de Szegő se satisface trivialmente, puesto que $\log(c)$ es integrable \cite{szego1975orthogonal}. Más aún, en este contexto la clausura de los polinomios analíticos en $L^2(\nu)$ se identifica con un espacio de Hardy ponderado, lo que proporciona el entorno funcional en el que aparecen de forma natural las evaluaciones puntuales acotadas en el interior del disco \cite{duren1970theory}. Esta observación explica por qué la medida de Lebesgue constituye la medida base idónea en nuestro problema.
\end{remark}

\subsection{Espacios de Polinomios y Ortogonalidad}

Esta sección recoge las nociones básicas de la teoría clásica de Polinomios Ortogonales que sirven de base para las generalizaciones de Sobolev desarrolladas en las secciones posteriores.

Sea $\mu$ una medida de Borel positiva con soporte infinito en el plano complejo $\mathbb{C}$. Denotamos por $\mathbb{P}$ al espacio vectorial de los polinomios con coeficientes complejos.

\begin{definition}[Producto Interno Estándar] \label{def:producto_estandar}
	Definimos el producto interno estándar en el espacio de Hilbert $L^2(\mu)$ como:
	\begin{equation}
		\langle f, g \rangle_{L^2(\mu)} = \int_{\text{supp}(\mu)} f(z) \overline{g(z)} \, d\mu(z).
	\end{equation}
\end{definition}

A partir de este producto interno, se construye la sucesión de Polinomios Ortogonales (PO) mediante el proceso de Gram-Schmidt aplicado a la base canónica $\{1, z, z^2, \dots\}$ \cite{szego1975orthogonal}.

\begin{definition}[Polinomios Ortonormales] \label{def:polinomios_ortonormales}
	Sea $\{\phi_n\}_{n \geq 0}$ la sucesión de polinomios ortonormales respecto a $\mu$, donde $\deg(\phi_n) = n$ y el coeficiente principal $\kappa_n > 0$. Estos satisfacen la condición de ortogonalidad:
	\begin{equation}
		\langle \phi_n, \phi_m \rangle_{L^2(\mu)} = \delta_{nm}, \quad \forall n,m \geq 0.
	\end{equation}
\end{definition}

Una herramienta fundamental para el estudio algebraico de estos polinomios es la matriz que recoge los momentos de la medida.

\begin{definition}[Matriz de Momentos] \label{def:matriz_momentos}
	La matriz de momentos $\mathbf{M}$ asociada a la medida $\mu$ es la matriz de Gram de la base canónica, cuyos elementos están dados por:
	\begin{equation}
		(\mathbf{M})_{i,j} = \langle z^i, z^j \rangle_{L^2(\mu)} = \int z^i \overline{z^j} \, d\mu(z), \quad i,j \geq 0.
	\end{equation}
	La existencia y positividad de los menores principales de Gram de $\mathbf{M}$ es condición necesaria y suficiente para la existencia de la familia de polinomios ortogonales.
\end{definition}

Un concepto geométrico fundamental que conecta el soporte de la medida con los puntos de evaluación acotada es la envoltura convexa polinómica.

\begin{definition}[Envoltura convexa polinómica]\label{def:envoltura_polinomica}
	Para un compacto $K \subset \mathbb{C}$, la \emph{envoltura convexa polinómica} de $K$ se define como
	\[
		P_C(K) = \Bigl\{ z \in \mathbb{C} : |p(z)| \leq \max_{\xi \in K} |p(\xi)|\ \text{ para todo } p \in \Poly \Bigr\}.
	\]
\end{definition}

Todo punto de evaluación acotada de una medida $\mu$ pertenece a $P_C(\operatorname{supp}(\mu))$. Para el caso en que $K = \mathbb{S}_r(a)$, se tiene $P_C(K) = \overline{\mathbb{D}_r(a)}$, lo que conecta la geometría del soporte con la región interior al disco.

\subsection{Espacios de Sobolev y Productos Discretos}

La generalización central de este trabajo consiste en perturbar el producto interno estándar $L^2(\mu)$ mediante la adición de términos que involucran derivadas en un conjunto finito de puntos discretos.

Sea $C = \{c_1, \dots, c_N\} \subset \mathbb{C}$ un conjunto finito de puntos en el plano complejo.

\begin{definition}[Producto Interno de Sobolev Discreto]
	Definimos el producto interno de Sobolev $\langle \cdot, \cdot \rangle_S$ en el espacio de polinomios $\Poly$ como:
	\begin{equation} \label{eq:sobolev_inner_prod}
		\langle p, q \rangle_S = \langle p, q \rangle_{L^2(\mu)} + \sum_{k=1}^N p'(c_k) \overline{q'(c_k)}.
	\end{equation}
\end{definition}

Un aspecto central que distingue este marco del caso clásico es la completitud del espacio. El espacio de Hilbert de Sobolev $\mathcal{H}$ se define como la clausura de los polinomios bajo la métrica inducida por $\langle \cdot, \cdot \rangle_S$.

\begin{definition}[Espacio de Hilbert $\mathcal{H}$]
	\begin{equation}
		\mathcal{H} = \overline{\Poly}^{\|\cdot\|_S}
	\end{equation}
\end{definition}

En el caso concreto estudiado en este trabajo, tomaremos como componente continua la medida de Lebesgue normalizada $\mathbf{m}_{a;r}$ introducida en la subsección anterior. Con esta notación, el producto interno de Sobolev objeto de estudio toma la forma
\begin{equation} \label{eq:sobolev_concreto}
	\langle p, q \rangle_S = \int_{\mathbb{S}_r(a)} p(z)\overline{q(z)} \, d\mathbf{m}_{a;r}(z) + \sum_{k=1}^N p'(c_k)\overline{q'(c_k)}, \qquad p, q \in \Poly.
\end{equation}

\subsection{Evaluaciones Puntuales Acotadas (BPE) y Puntos $\delta$-singulares}

El comportamiento del espacio $\mathcal{H}$ y, en consecuencia, la estabilidad numérica de los ceros de los polinomios ortogonales, depende críticamente de la posición de los puntos $c_k$ respecto al soporte de la medida $\mu$.

Denotamos por $R$ al radio del disco minimal que contiene al soporte de la medida $\mu$ (o el radio esencial en caso de medidas con soporte no compacto).

\begin{definition}[Funcionales de Evaluación Algebraicos]
	Para cualquier $\xi \in \C$, definimos el funcional de evaluación lineal sobre el espacio de polinomios $\delta_\xi : \Poly \to \C$ como $\delta_\xi(p) = p(\xi)$. Análogamente, definimos la evaluación de la derivada como $\delta'_\xi(p) = p'(\xi)$.
\end{definition}

\begin{definition}[Funcionales de evaluación restringidos]\label{def:funcionales_restringidos}
	Sea $c\in\C$. Definimos los funcionales lineales
	\[
		\delta_c:\Poly\to\C,\qquad \delta_c(p)=p(c),
	\]
	\[
		\delta'_c:\Poly\to\C,\qquad \delta'_c(p)=p'(c).
	\]
	Definimos el subespacio \emph{núcleo de la derivada en $c$} como
	\[
		N_c:=\ker(\delta'_c)=\{p\in\Poly:p'(c)=0\}.
	\]
	Notemos que $\operatorname{codim}_{\Poly}(N_c)=1$, ya que $\delta'_c$ no es el funcional nulo
	(basta tomar $p(z)=z-c$, para el cual $p'(c)=1$).
\end{definition}

\begin{definition}[Punto $\delta$-singular respecto de $\|\cdot\|_S$]\label{def:delta_singular}
	Sea $\langle\cdot,\cdot\rangle_S$ un producto interno de Sobolev sobre $\Poly$.
	Diremos que un punto $c\in\C$ es \emph{$\delta$-singular} respecto a
	$\|\cdot\|_S$ si la restricción del funcional de evaluación
	$\delta_c$ al subespacio $N_c$ no es acotada respecto de la norma
	$\|\cdot\|_S$, es decir, si
	\[
		\sup\Bigl\{\frac{|p(c)|}{\|p\|_S}:p\in N_c\setminus\{0\}\Bigr\}
		=+\infty.
	\]
	Equivalentemente, existe una sucesión $(p_n)\subset N_c$ con
	\[
		\|p_n\|_S\to 0,\qquad |p_n(c)|=1\quad\text{para todo }n\ge1.
	\]
\end{definition}

Basándonos en la geometría de los puntos respecto al soporte continuo, clasificamos el conjunto de índices $\mathcal{I} = \{1, \dots, N\}$ en dos categorías:

% [REEMPLAZADO] Definición original de Puntos Interiores/Exteriores basada en |c_k|≥R — ahora generalizada mediante δ-singularidad
\begin{definition}[Puntos BPE-derivada y puntos $\delta$-singulares]\label{def:class_points_new}
	Para un producto de Sobolev discreto $\langle\cdot,\cdot\rangle_S$
	fijado, dividimos el conjunto de masas discretas como
	$C=I_{\mathrm{in}}\cup I_{\mathrm{out}}$, donde:
	\begin{itemize}
		\item $I_{\mathrm{in}}:=\{k:c_k\text{ no es }\delta\text{-singular}\}$
		      (puntos donde $\delta_{c_k}$ y $\delta'_{c_k}$ son ambos acotados
		      respecto a $\|\cdot\|_S$).
		\item $I_{\mathrm{out}}:=\{k:c_k\text{ es }\delta\text{-singular}\}$.
	\end{itemize}
	Sea $m=|I_{\mathrm{out}}|$.
\end{definition}

\begin{remark}[Compatibilidad con la versión anterior]\label{rem:compatibilidad_class_points}
	En el caso modelo (Lebesgue normalizada en $\mathbb{S}_R$), por el
	Teorema~\ref{thm:equivalencia_lebesgue},
	\[
		I_{\mathrm{in}}=\{k:|c_k|<R\},\qquad I_{\mathrm{out}}=\{k:|c_k|\ge R\}.
	\]
	Esto recupera exactamente la clasificación geométrica usada en el TFM
	original. La nueva definición es estrictamente más general.
\end{remark}

Esta distinción entre $I_{\mathrm{in}}$ y $I_{\mathrm{out}}$ constituye la base del análisis espectral y del estudio de la codimensión topológica del espacio desarrollado en la Sección~\ref{sec:espectral}. La noción de \textit{Bounded Point Evaluation} (BPE) formaliza esta separación.

\begin{definition}[Punto de Evaluación Acotada (BPE)]
	Decimos que $a \in \C$ es un punto de evaluación acotada (\textit{bounded point evaluation}, BPE) de una medida $\mu$ \cite{miller1999bounded,aleman2009nontangential} si existe una constante $C > 0$ tal que:
	\begin{equation}
		\abs{p(a)}^2 \leq C \int \abs{p(z)}^2 \, d\mu(z), \qquad p(z) \in \Poly.
	\end{equation}
	Equivalentemente, por el Teorema de Representación de Riesz, esto corresponde a la existencia de un núcleo reproductor $K_a \in \overline{\Poly}^{L^2(\mu)}$ tal que $p(a) = \langle p, K_a \rangle_{L^2(\mu)}$ para todo $p \in \Poly$.
\end{definition}

En el contexto de nuestro problema, la distinción entre puntos $I_{\mathrm{in}}$ y $I_{\mathrm{out}}$ se interpreta como una dicotomía de continuidad de evaluaciones. Para $k\in I_{\mathrm{in}}$, la evaluación puntual y la evaluación de la derivada son continuas (véanse la Proposición siguiente y el Lema~\ref{lema:ABPE_derivada}). En cambio, en el régimen exterior no-BPE la evaluación puntual deja de ser continua respecto a la norma $L^2(\mu)$ de la componente continua. Esta pérdida de continuidad es precisamente la que obliga, más adelante, a reemplazar el enfoque topológico por uno algebraico en el estudio de $\D$.

\begin{proposition}[Acotación y Extensión Continua en el Interior]
	Para los puntos interiores ($k \in I_{\mathrm{in}}$), el funcional de evaluación $\delta_{c_k}$ es continuo respecto a la norma de Sobolev $\norm{\cdot}$. En consecuencia, puede extenderse de forma única a un funcional lineal acotado en todo el espacio de Hilbert $\mathcal{H}$, y su núcleo es un subespacio cerrado.
\end{proposition}

\begin{proof}
	Para cualquier polinomio $p \in \Poly$ y un punto $c_k$ con $\abs{c_k} < R$, podemos aplicar la Fórmula Integral de Cauchy sobre la circunferencia $\mathbb{S}_R$:
	$$ p(c_k) = \frac{1}{2\pi i} \int_{\mathbb{S}_R} \frac{p(z)}{z - c_k} dz. $$
	Tomando valores absolutos y aplicando la desigualdad de Cauchy-Schwarz:
	$$ \abs{p(c_k)}^2 \leq \left( \frac{1}{2\pi} \int_{\mathbb{S}_R} \frac{1}{\abs{z - c_k}^2} \abs{dz} \right) \left( \int_{\mathbb{S}_R} \abs{p(z)}^2 \abs{dz} \right). $$
	Puesto que $\abs{c_k} < R$, la distancia al borde de la circunferencia está estrictamente acotada por abajo, luego la primera integral es una constante finita $M_{c_k} > 0$. Además, la segunda integral está dominada por la parte continua de la norma de Sobolev (salvo una constante de proporcionalidad dependiente de la medida de Lebesgue). Por tanto:
	$$ \abs{p(c_k)}^2 \leq \tilde{M}_{c_k} \|p\|_{L^2(\mu)}^2 \leq \tilde{M}_{c_k} \norm{p}^2. $$
	Esto demuestra que el funcional está acotado en $\Poly$. Por el Teorema de Extensión de Operadores Lineales Acotados en espacios densos \cite{kreyszig1989introductory}, el funcional se extiende de manera continua y única a todo $\mathcal{H}$. Al ser continuo, su núcleo es la preimagen de un conjunto cerrado (el $\{0\}$), por lo que es un subespacio cerrado.
\end{proof}

\begin{lemma}[Continuidad del funcional derivada para puntos ABPE]
	\label{lema:ABPE_derivada}
	Sea $\mu$ la medida de Lebesgue (o, más generalmente, una medida que satisfaga la condición de Szegő) soportada en la circunferencia $\mathbb{T}_R = \{z \in \mathbb{C} : |z| = R\}$. Si $c \in \mathbb{C}$ es un punto de evaluación acotada (BPE) para la norma $L^2(\mu)$ y se encuentra en el interior de la curva ($|c| < R$), entonces $c$ es un punto de Evaluación Analítica Acotada (ABPE). En consecuencia, el funcional de evaluación de la derivada, $\delta'_c(p) = p'(c)$, es un funcional lineal continuo.
\end{lemma}

\begin{proof}
	Al ser $c$ un punto BPE situado en el interior de la región acotada por el soporte de la medida, la teoría clásica de espacios de Hardy y de los espacios $P^t(\mu)$ garantiza que $c$ es, de hecho, un punto ABPE \cite{miller1999bounded,aleman2009nontangential}. Esto implica que existe un entorno abierto $\mathcal{U}$ de $c$ contenido en el disco $|z| < R$ donde la evaluación puntual es uniformemente acotada.

	Sea $r > 0$ suficientemente pequeño tal que la circunferencia cerrada $\gamma_r = \{z \in \mathbb{C} : |z-c| = r\}$ esté contenida en $\mathcal{U}$. Dado que los polinomios son funciones enteras, podemos aplicar la Fórmula Integral de Cauchy para la derivada:
	\[
		p'(c) = \frac{1}{2\pi i} \oint_{\gamma_r} \frac{p(z)}{(z-c)^2} dz
	\]
	Tomando valor absoluto y acotando la integral por el máximo de la función en la curva:
	\[
		|p'(c)| \le \frac{1}{2\pi} \left( 2\pi r \right) \frac{\max_{z \in \gamma_r} |p(z)|}{r^2} = \frac{1}{r} \max_{z \in \gamma_r} |p(z)|
	\]
	Como todos los puntos $z \in \gamma_r$ son BPEs, existe una constante $M > 0$ tal que $|p(z)| \le M \|p\|_{L^2(\mu)}$ para todo $z \in \gamma_r$. Sustituyendo esta cota obtenemos:
	\[
		|p'(c)| \le \frac{M}{r} \|p\|_{L^2(\mu)} \le \frac{M}{r} \|p\|_S
	\]
	Lo que demuestra que existe una constante $C = M/r < \infty$ tal que $|\delta'_c(p)| \le C \|p\|_S$, probando así que la evaluación de la derivada es un funcional lineal continuo.
\end{proof}

\begin{lemma}[No acotación de la evaluación de derivada en frontera y exterior]
	\label{lem:no_acotacion_derivada_exterior_lebesgue}
	Sea $\mu=\mathbf m_{0;R}$ la medida de Lebesgue normalizada en la circunferencia
	$\mathbb{S}_R=\{z\in\C:\,|z|=R\}$. Si $c\in\C$ satisface $|c|\ge R$, entonces el funcional
	\[
		\delta'_c:\Poly\to\C,
		\qquad
		\delta'_c(p)=p'(c),
	\]
	no es acotado respecto de la norma $L^2(\mu)$.
\end{lemma}

\begin{proof}
	Para $n\ge 1$, consideremos
	\[
		p_n(z):=\left(\frac{z}{R}\right)^n.
	\]
	Entonces $\|p_n\|_{L^2(\mu)}=1$, pues $|z|=R$ sobre $\mathbb{S}_R$. Además,
	\[
		p_n'(z)=\frac{n}{R}\left(\frac{z}{R}\right)^{n-1},
		\qquad
		|p_n'(c)|=\frac{n}{R}\left(\frac{|c|}{R}\right)^{n-1}.
	\]
	Si $|c|>R$, se tiene $|p_n'(c)|\to\infty$ exponencialmente. Si $|c|=R$,
	entonces $|p_n'(c)|=n/R\to\infty$. En ambos casos,
	\[
		\sup_{\|p\|_{L^2(\mu)}=1}|p'(c)|=\infty,
	\]
	luego $\delta'_c$ no es acotado.
\end{proof}

% [REEMPLAZADO] Corolario original: Implicación en el caso no-BPE de Lebesgue — ahora subsumido por thm:equivalencia_lebesgue

\begin{proposition}[Caracterización equivalente vía $\psi_c$]\label{prop:caracterizacion_psi_c}
	Sea $c\in\C$ y consideremos el funcional
	\[
		\psi_c:\Poly\to\C,\qquad \psi_c(p)=p(c)+cp'(c)=(\D p)'(c).
	\]
	Entonces son equivalentes:
	\begin{enumerate}
		\item $c$ es $\delta$-singular respecto a $\|\cdot\|_S$;
		\item $\psi_c$ no es acotado respecto a $\|\cdot\|_S$;
		\item $\delta_c$ no es acotado en $\Poly$ respecto a $\|\cdot\|_S$.
	\end{enumerate}
\end{proposition}

\begin{proof}
	\textbf{(1)$\Rightarrow$(3):} es inmediato, pues $N_c\subset\Poly$.

	\textbf{(3)$\Rightarrow$(1):} sea $(p_n)\subset\Poly$ con
	$\|p_n\|_S\le 1$ y $|p_n(c)|\to\infty$. Definimos
	\[
		q_n(z):=p_n(z)-p_n'(c)(z-c).
	\]
	Entonces $q_n'(c)=p_n'(c)-p_n'(c)=0$, por tanto $q_n\in N_c$.
	Por otra parte,
	\[
		q_n(c)=p_n(c).
	\]
	Solo falta acotar $\|q_n\|_S$. Como la proyección sobre $N_c$ respecto de la descomposición
	$\Poly=N_c\oplus\operatorname{span}\{z-c\}$ es acotada (por ser $N_c$ de codimensión
	finita), existe $C>0$ tal que $\|q_n\|_S\le C\|p_n\|_S$ para toda la sucesión
	construida. En particular, la sucesión $(q_n)$ está acotada en $\|\cdot\|_S$
	y cumple $q_n(c)=p_n(c)=1$. Normalizando $\tilde q_n=q_n/\|q_n\|_S$ se obtiene
	$\tilde q_n\in N_c$, $\|\tilde q_n\|_S=1$ y $|\tilde q_n(c)|\to\infty$,
	lo cual prueba que $c$ es $\delta$-singular.

	\textbf{(2)$\iff$(3):} Si $\delta'_c$ es acotado, entonces
	$\psi_c=\delta_c+c\delta'_c$ es no acotado si y solo si $\delta_c$ lo es.
	Si $\delta'_c$ no es acotado, entonces \emph{a fortiori} $\psi_c$ es no
	acotado (pues $\psi_c$ extiende a $\delta_c$ sobre $N_c$ más el término
	$c\delta'_c$). Por tanto, en ambos casos, $\psi_c$ no acotado
	$\Longleftrightarrow$ $\delta_c$ no acotado en $\Poly$.
\end{proof}

\begin{theorem}[Equivalencia $\delta$-singular $\iff$ no-BPE en el caso Lebesgue]\label{thm:equivalencia_lebesgue}
	Sea $\mu=\mathbf m_{0;R}$ la medida de Lebesgue normalizada en
	la circunferencia $\mathbb{S}_R=\{z\in\C:|z|=R\}$, y consideremos el producto
	de Sobolev discreto
	\[
		\langle p,q\rangle_S
		=
		\int_{\mathbb{S}_R}p(z)\overline{q(z)}\,d\mathbf m_{0;R}(z)
		+
		\sum_{k=1}^N p'(c_k)\overline{q'(c_k)}.
	\]
	Para cualquier $c\in\C$, son equivalentes:
	\begin{enumerate}
		\item $c$ es $\delta$-singular respecto a $\|\cdot\|_S$;
		\item $|c|\ge R$;
		\item $c$ es no-BPE para $L^2(\mu)$.
	\end{enumerate}
\end{theorem}

\begin{proof}
	\textbf{(2)$\iff$(3):} es la caracterización clásica de los puntos BPE
	para la medida de Lebesgue en circunferencia: los BPE son exactamente
	los puntos del disco abierto $|c|<R$ (espacio de Hardy del disco)
	\cite{duren1970theory,szego1975orthogonal}.

	\textbf{(2)$\Rightarrow$(1):} suponemos $|c|\ge R$.
	Por la Proposición~\ref{prop:no_acotacion_modelo_exterior},
	el funcional $\psi_c(p)=p(c)+cp'(c)=(\D p)'(c)$ no es acotado
	respecto a $\|\cdot\|_S$. Por la
	Proposición~\ref{prop:caracterizacion_psi_c}, $c$ es $\delta$-singular.

	\textbf{(1)$\Rightarrow$(2):} probamos el contrarrecíproco.
	Suponemos $|c|<R$. Entonces $c$ es BPE para $L^2(\mu)$ (caso (3)
	falso) y, por el Lema~\ref{lema:ABPE_derivada}, $c$ es ABPE: tanto $\delta_c$
	como $\delta'_c$ son funcionales acotados respecto a $\|\cdot\|_{L^2(\mu)}$
	y, en consecuencia, también respecto a $\|\cdot\|_S$ (porque
	$\|\cdot\|_{L^2(\mu)}\le\|\cdot\|_S$). Como $\delta_c$ es acotado en
	$\Poly$, en particular es acotado al restringirlo a $N_c\subset\Poly$,
	luego $c$ no es $\delta$-singular.
\end{proof}

\begin{remark}[Hipótesis primaria: $\delta$-singularidad]\label{rem:hipotesis_delta_singular}
	El Teorema~\ref{thm:equivalencia_lebesgue} es un \emph{hecho geométrico}
	propio del caso Lebesgue: la frontera entre la región BPE y la región
	$\delta$-singular coincide exactamente con la circunferencia.
	A partir de este punto, la condición de $\delta$-singularidad se utilizará
	como hipótesis técnica primaria en todo el trabajo. La noción clásica
	de punto no-BPE para $L^2(\mu)$ queda como caso geométrico particular,
	válido en el caso modelo de la medida de Lebesgue normalizada en circunferencia.
\end{remark}

\begin{lemma}[Normalización de un funcional no acotado]\label{lem:normalizacion_funcional_no_acotado}
	Sea $X$ un espacio normado complejo y sea $\varphi:X\to\mathbb C$ un funcional lineal no acotado.
	Entonces existe una sucesión $(x_n)_{n\ge1}\subset X$ tal que
	\[
		\|x_n\|_X\to 0,
		\qquad
		\varphi(x_n)=1
		\quad \text{para todo } n.
	\]
\end{lemma}

\begin{proof}
	Como $\varphi$ no es acotado, no está acotado en la bola unidad. Por tanto, para cada
	$n\ge1$ existe $y_n\in X$ tal que
	\[
		\|y_n\|_X\le 1,
		\qquad
		|\varphi(y_n)|\ge n.
	\]
	Definimos
	\[
		x_n:=\frac{y_n}{\varphi(y_n)}.
	\]
	Entonces
	\[
		\varphi(x_n)=1,
	\]
	y además
	\[
		\|x_n\|_X
		=
		\frac{\|y_n\|_X}{|\varphi(y_n)|}
		\le
		\frac1n
		\xrightarrow[n\to\infty]{}0.
	\]
\end{proof}

\subsection{Organización del Artículo}

El artículo se organiza como sigue:

\begin{itemize}
	\item \textbf{Sección~\ref{sec:espectral}: Análisis Espectral y Topología.} Se desarrolla el núcleo teórico del trabajo. Se estudian la codimensión topológica del espacio de Sobolev, la estructura del operador multiplicación y los números de Gelfand, estableciendo la dicotomía entre puntos $\delta$-singulares y no $\delta$-singulares, y la relación con la propiedad BPE clásica en el caso modelo de Lebesgue.

	\item \textbf{Sección~\ref{sec:matricial}: Análisis Matricial y Comportamiento Asintótico.} Se analizan las representaciones matriciales de los operadores. Se estudia el comportamiento asintótico de las secciones finitas (matrices de Hessenberg) mediante la teoría de perturbaciones y los valores singulares, se demuestra la estabilización algebraica, y se introduce la desigualdad de Weyl-Horn para el conteo de ceros excepcionales.

	\item \textbf{Sección~\ref{sec:ceros_acotacion}: Localización y Acotación Asintótica de los Ceros.} Se aplican los resultados obtenidos para acotar la ubicación de los ceros de los polinomios ortogonales. La sección se organiza en cuatro subsecciones: ceros como autovalores, acotación uniforme vía cuasi-ortogonalidad (vía clásica), la cadena estructural $Q_m$--Weyl-Horn--atracción (contribución novedosa), y la equivalencia y comparación entre ambos enfoques.

	\item \textbf{Conclusiones y Líneas Futuras.} Se resumen las aportaciones principales del trabajo y se proponen vías abiertas para investigación futura.
\end{itemize}

% =============================================================
% CAPÍTULO 2: ANÁLISIS ESPECTRAL Y TOPOLOGÍA
% 2.1 Codimensión y Subespacios en el Espacio de Hilbert
% 2.2 La Dicotomía Topológica
% 2.3 El Operador de Multiplicación y los Números de Gelfand
% 2.4 El Caso No-BPE: La Explosión Topológica
% =============================================================

\chapter{Análisis Espectral y Topología en Espacios de Sobolev}\label{sec:espectral}

Para el estudio del operador de multiplicación, es fundamental distinguir el comportamiento sobre distintos subespacios del espacio de Hilbert $\mathcal{H}$, determinados por las condiciones de anulación en los conjuntos $I_{\mathrm{in}}$ e $I_{\mathrm{out}}$. Se demuestra, en el transcurso de esta sección, que la codimensión topológica de $\mathcal{H}$ en $H^2(\mathbb{D})$ está directamente determinada por el número de puntos en la región inestable $I_{\mathrm{out}}$; este déficit dimensional guarda una estrecha relación con la teoría de índices de defecto de operadores simétricos y con los Tripletes de Gelfand.

% [CORRECCIÓN ESTRUCTURAL] Se formaliza la hipótesis de trabajo sobre la medida continua.
\begin{remark}[Hipótesis de trabajo]
	En lo sucesivo, salvo indicación contraria, asumimos que la parte continua del producto es la medida de Lebesgue normalizada en $\mathbb{S}_R(0)$.
\end{remark}

% -------------------------------------------------------------
\subsection{Codimensión y Subespacios en el Espacio de Hilbert}
% -------------------------------------------------------------

\begin{definition}[Codimensión topológica en Hilbert]\label{def:codim_top_hilbert}
	Sea $\mathcal{H}$ un espacio de Hilbert y sea $M\subset \mathcal{H}$ un subespacio cerrado. Definimos su codimensión topológica por
	$$
		\operatorname{codim}_{\mathcal{H}}(M):=\dim(\mathcal{H}/M).
	$$
	Equivalentemente,
	$$
		\operatorname{codim}_{\mathcal{H}}(M)=\dim(M^\perp)
	$$
	\cite{halmos1982hilbert}.
\end{definition}

\begin{remark}[Convención de notación para codimensión]\label{rem:convencion_codim}
	A lo largo del trabajo distinguiremos dos nociones:
	\begin{itemize}
		\item $\operatorname{codim}_{\mathcal{H}}(M)$ para subespacios cerrados $M$ de un espacio de Hilbert (noción topológica).
		\item $\operatorname{codim}_{\Poly}(V)$ para subespacios vectoriales $V\subset \Poly$ (noción puramente algebraica; véase la Definición~\ref{def:codim_alg}).
	\end{itemize}
	En el primer caso, la dimensión se entiende en sentido hilbertiano (cardinal de una base ortonormal); en el segundo, en sentido algebraico (cardinal de una base de Hamel).
	Cuando el espacio ambiente sea finito-dimensional (por ejemplo, $\mathbb{C}^n$), usaremos $\operatorname{codim}(\cdot)$ sin subíndice.
\end{remark}

Dado que nuestra medida $\mu$ es la medida de Lebesgue en la circunferencia unidad $\mathbb{S}_1$, la componente continua de nuestro espacio de Sobolev está ligada al espacio de Hardy \cite{duren1970theory}.

\begin{definition}[Espacio de Hardy $H^2(\mathbb{D})$]
	El espacio de Hardy $H^2(\mathbb{D})$ se define como el espacio de funciones holomorfas $f$ en el disco unidad abierto $\mathbb{D} = \{z \in \C : \abs{z} < 1\}$ tales que sus medias cuadráticas sobre circunferencias concéntricas permanecen acotadas:
	$$ \sup_{0 \le r < 1} \frac{1}{2\pi} \int_{0}^{2\pi} \abs{f(re^{i\theta})}^2 d\theta < \infty. $$
\end{definition}

Existe un isomorfismo isométrico natural entre $H^2(\mathbb{D})$ y el subespacio cerrado de $L^2(\mathbb{S}_1)$ formado por funciones cuyos coeficientes de Fourier de índice negativo son nulos. En nuestro contexto de aproximación polinómica, podemos caracterizar $H^2$ de manera equivalente como la clausura del espacio de polinomios $\mathbb{P}$ bajo la norma $L^2$ estándar en la frontera:
$$ H^2 \cong \overline{\mathbb{P}}^{L^2(\mathbb{S}_1)}. $$
Esta identificación es fundamental, ya que nos permite utilizar la estructura de \textit{Reproducing Kernel Hilbert Space} (RKHS) de $H^2$. En particular, para cualquier $c \in \mathbb{D}$, el núcleo reproductor de Szegő viene dado por $K_c(z) = (1-\bar{c}z)^{-1}$.

% -------------------------------------------------------------
\subsection{La Dicotomía Topológica}
% -------------------------------------------------------------

\begin{definition}[Subespacio Estable $V$]
	Definimos el subespacio estable algebraico $V_0$ como el conjunto de polinomios que se anulan en los puntos exteriores:
	\begin{equation}
		V_0 = \{p \in \mathbb{P} : p(c_k) = 0, \; \forall k \in I_{\mathrm{out}}\}.
	\end{equation}
	El subespacio estable $V \subset \mathcal{H}$ se define como la clausura topológica de $V_0$ respecto a la norma de Sobolev: $V = \overline{V_0}^{\|\cdot\|_S}$.
\end{definition}

\begin{definition}[Subespacio Totalmente Restringido $V'$]
	Definimos el subespacio algebraico $V'_0$ imponiendo restricciones nulas en la totalidad de los puntos:
	\begin{equation}
		V'_0 = \{p \in \mathbb{P} : p(c_k) = 0, \; \forall k \in \{1, \dots, N\}\}.
	\end{equation}
	Y definimos el subespacio de Hilbert correspondiente como su clausura: $V' = \overline{V'_0}^{\|\cdot\|_S}$.
\end{definition}

\begin{proposition}[Codimensión de la clausura del ideal para un átomo BPE]
	Sea $\mathcal{H}$ un espacio de Hilbert tal que el espacio de polinomios $\Poly$ es denso en $\mathcal{H}$. Sea $a \in \mathbb{C}$ un punto de evaluación acotada (BPE) respecto a la norma de $\mathcal{H}$. Entonces, el ideal algebraico $Y_a = \{p \in \Poly : p(a) = 0\}$ es un subespacio cerrado de codimensión topológica exactamente 1.
\end{proposition}

\begin{proof}
	Consideremos el ideal $Y_a = (z-a)\Poly$. Queremos probar que tiene codimensión topológica 1.

	Dado que $a$ es un punto de evaluación continua (BPE), el funcional lineal $\Phi_a: \Poly \to \mathbb{C}$ dado por $\Phi_a(p) = p(a)$ es acotado. Por el Teorema de Extensión, $\Phi_a$ se extiende de manera única y continua a todo $\mathcal{H}$. Al ser continuo, su núcleo $\ker(\Phi_a)$ es un subespacio cerrado de $\mathcal{H}$.

	Por otro lado, cualquier polinomio $p \in \Poly$ admite la siguiente descomposición algebraica:
	$$p(z) = (p(z) - p(a)) + p(a) \cdot 1$$
	El primer término pertenece a $Y_a$ y el segundo al subespacio generado por el polinomio constante $\mathbf{1}$. Por tanto, $\Poly = Y_a \oplus \operatorname{span}\{\mathbf{1}\}$.

	Tomando clausuras en $\mathcal{H}$, y sabiendo que $\Poly$ es denso en $\mathcal{H}$:
	$$\mathcal{H} = \overline{\Poly} = \overline{Y_a \oplus \operatorname{span}\{\mathbf{1}\}}$$
	Como el funcional es continuo y $\Phi_a(\mathbf{1}) = 1 \neq 0$, la suma directa se preserva en la clausura:
	$$\mathcal{H} = Y_a \oplus \operatorname{span}\{\mathbf{1}\}$$
	Lo que demuestra que $\operatorname{codim}_{\mathcal{H}}(Y_a) = 1$.
\end{proof}

\begin{proposition}[Codimensión de la clausura para múltiples átomos BPE]
	Sea $\mathcal{H}$ un espacio de Hilbert donde $\Poly$ es denso. Sea $\mathcal{A} = \{a_1, \dots, a_N\} \subset \mathbb{C}$ un conjunto de $N$ puntos distintos que son de evaluación acotada (BPE) respecto a la norma de $\mathcal{H}$. El subespacio $Y_{\mathcal{A}} = \{p \in \Poly : p(a_k) = 0 \text{ para } k=1,\dots,N\}$ es un subespacio cerrado de codimensión topológica exactamente $N$.
\end{proposition}

\begin{proof}
	Consideremos el funcional vectorial continuo $\vec{\Phi}: \Poly \to \mathbb{C}^N$ dado por $\vec{\Phi}(p) = (p(a_1), \dots, p(a_N))$. Por definición, $Y_{\mathcal{A}} = \ker(\vec{\Phi})$, que es cerrado por ser $\vec{\Phi}$ continuo.

	Sean $L_1(z), \dots, L_N(z)$ los polinomios fundamentales de interpolación de Lagrange asociados a $\mathcal{A}$, los cuales satisfacen $L_i(a_j) = \delta_{ij}$. Cualquier polinomio $p \in \Poly$ se descompone unívocamente como:
	$$p(z) = \underbrace{\left( p(z) - \sum_{k=1}^N p(a_k)L_k(z) \right)}_{\in Y_{\mathcal{A}}} + \underbrace{\sum_{k=1}^N p(a_k)L_k(z)}_{\in \operatorname{span}\{L_1, \dots, L_N\}}$$
	Esto demuestra la descomposición en suma directa algebraica $\Poly = Y_{\mathcal{A}} \oplus \operatorname{span}\{L_1, \dots, L_N\}$.

	Como los $N$ puntos son BPE, el funcional vectorial $\vec{\Phi}$ se extiende de manera continua a $\mathcal{H}$. Tomando clausuras y apoyándonos en la densidad de los polinomios ($\overline{\Poly} = \mathcal{H}$):
	$$\mathcal{H} = \overline{Y_{\mathcal{A}} \oplus \operatorname{span}\{L_1, \dots, L_N\}}$$
	Concluimos que $\operatorname{codim}_{\mathcal{H}}(Y_{\mathcal{A}}) = N$.
\end{proof}

\begin{theorem}[Dicotomía de la Codimensión]
	\label{thm:dicotomia}
	Sea $\mathcal{H}$ la clausura de $\Poly$ respecto a una norma de Hilbert y sea $a \in \C$. Sea $V = \overline{\{p \in \Poly : p(a)=0\}}$. Entonces, solo existen dos posibilidades:
	\begin{enumerate}
		\item Si $a$ es un punto de evaluación acotada (BPE), entonces $V$ es un hiperplano cerrado y $\operatorname{codim}_{\mathcal{H}}(V) = 1$.
		\item Si $a$ no es un BPE, entonces $V$ es denso en $\mathcal{H}$ y $\operatorname{codim}_{\mathcal{H}}(V) = 0$.
	\end{enumerate}
\end{theorem}

\begin{proof}
	El caso 1 se sigue de las proposiciones anteriores.

	Si el funcional de evaluación $\delta_a: \Poly \to \C$ no es acotado respecto a $\|\cdot\|_{\mathcal{H}}$, demostraremos que su núcleo algebraico $Y_a = \{p \in \Poly : p(a)=0\}$ es denso en $\mathcal{H}$.

	Por definición de operador no acotado, para todo $n \in \mathbb{N}$ existe un polinomio $h_n \in \Poly$ tal que $|h_n(a)| > n \|h_n\|_{\mathcal{H}}$. Definimos la sucesión normalizada:
	$$ q_n(z) = \frac{h_n(z)}{h_n(a)}, $$
	con $q_n(a) = 1$ y $\|q_n\|_{\mathcal{H}} < 1/n \to 0$.

	Consideremos la sucesión $p_n(z) = 1 - q_n(z)$. Se cumple $p_n(a) = 0$, luego $p_n \in Y_a$ para todo~$n$. Además:
	$$ \|1 - p_n\|_{\mathcal{H}} = \|q_n\|_{\mathcal{H}} < \frac{1}{n}, $$
	lo que prueba que $1 \in \overline{Y_a}$. Por linealidad y densidad de $\Poly$ en $\mathcal{H}$ concluimos que $\overline{Y_a} = \mathcal{H}$, esto es, $\operatorname{codim}_{\mathcal{H}}(V) = 0$.
\end{proof}

\begin{theorem}[Dicotomía de la Codimensión para un Funcional]\label{thm:dicotomia_1}
	Sea $\phi : \Poly \to \mathbb{C}$ un funcional lineal no nulo. Consideremos $\ker(\phi) = \{ p \in \Poly : \phi(p) = 0 \}$. Entonces, la codimensión topológica de su clausura en el espacio de Hilbert $\mathcal{H}$ satisface:
	\begin{enumerate}
		\item Si $\phi$ es no acotado en $\mathcal{H}$, entonces $\operatorname{codim}_\mathcal{H}(\overline{\ker(\phi)}) = 0$.
		\item Si $\phi$ es acotado (continuo) en $\mathcal{H}$, entonces $\operatorname{codim}_\mathcal{H}(\overline{\ker(\phi)}) = 1$.
	\end{enumerate}
\end{theorem}

\begin{proof}
	\textbf{Caso 1 ($\phi$ no acotado):} Por propiedades fundamentales de los espacios normados, si un funcional lineal no está acotado, su núcleo es denso en el espacio topológico. Por lo tanto, su clausura es todo el espacio: $\overline{\ker(\phi)} = \mathcal{H}$. Trivilamente, al no faltar ninguna dimensión para generar $\mathcal{H}$, su codimensión topológica es 0.

	\textbf{Caso 2 ($\phi$ acotado):} Dado que $\phi$ es continuo, su núcleo ya es un subespacio cerrado, por lo que $\overline{\ker(\phi)} = \ker(\phi)$ (tomando la clausura en $\mathcal{H}$).

	Para calcular su codimensión topológica, utilizaremos una aproximación constructiva mediante subespacios de dimensión finita. Sea $\mathbb{P}_n$ el espacio de polinomios de grado menor o igual a $n$, el cual tiene dimensión finita $n+1$.
	Consideremos la restricción del funcional a este subespacio, $\phi_n = \phi|_{\mathbb{P}_n} : \mathbb{P}_n \to \mathbb{C}$.

	Como $\phi$ no es el funcional nulo, existe un grado $n_0$ tal que para todo $n \ge n_0$, existe al menos un polinomio $q \in \mathbb{P}_n$ con $\phi_n(q) \neq 0$. Por tanto, para $n \ge n_0$, la imagen de $\phi_n$ es todo $\mathbb{C}$ (dimensión 1).

	Definimos $V_n = \ker(\phi) \cap \mathbb{P}_n$. Aplicando el Teorema del Rango-Nulidad al operador lineal $\phi_n$ en el espacio de dimensión finita $\mathbb{P}_n$:
	$$ \dim(\mathbb{P}_n) = \dim(\ker(\phi_n)) + \dim(\operatorname{Im}(\phi_n)) \implies n+1 = \dim(V_n) + 1 \implies \dim(V_n) = n. $$

	Al ser $\mathbb{P}_n$ un espacio de Hilbert de dimensión finita, podemos descomponerlo ortogonalmente como $\mathbb{P}_n = V_n \oplus^\perp \operatorname{span}\{w_n\}$, donde $w_n$ es un único polinomio (salvo constante) ortogonal a $V_n$.

	Fijemos el polinomio $q \notin \ker(\phi)$. Cualquier polinomio arbitrario $p \in \Poly$ (que pertenecerá a algún $\mathbb{P}_n$ para $n$ suficientemente grande) se puede descomponer algebraicamente como:
	$$ p = \underbrace{\left( p - \frac{\phi(p)}{\phi(q)} q \right)}_{:= k} + \underbrace{\frac{\phi(p)}{\phi(q)} q}_{:= \alpha q} $$
	Evaluando $\phi$ en el término $k$, obtenemos $\phi(k) = \phi(p) - \frac{\phi(p)}{\phi(q)}\phi(q) = 0$, luego $k \in \ker(\phi)$.
	Esto demuestra que $\Poly = \ker(\phi) \oplus \operatorname{span}\{q\}$, por lo que la codimensión algebraica en los polinomios es exactamente 1 (no puede ser mayor porque un solo vector $q$ completa el espacio).

	Finalmente, como el funcional es continuo en $\mathcal{H}$, esta suma directa algebraica se preserva al tomar la completación topológica, resultando en $\mathcal{H} = \overline{\ker(\phi)} \oplus \operatorname{span}\{q\}$. En consecuencia, la codimensión topológica es exactamente 1.
\end{proof}

\begin{theorem}[Dicotomía Generalizada para $N$ Funcionales]\label{thm:dicotomia_N_general}
	Sean $\phi_1, \dots, \phi_N : \Poly \to \mathbb{C}$ funcionales lineales linealmente independientes. Definimos su núcleo conjunto como el subespacio algebraico $W_0 = \bigcap_{j=1}^N \ker(\phi_j)$.
	Sea $k$ el número máximo de funcionales linealmente independientes en $\operatorname{span}\{\phi_1, \dots, \phi_N\}$ que son continuos (acotados) respecto a la norma de $\mathcal{H}$.
	Entonces, la codimensión topológica de la clausura de $W_0$ en el espacio de Hilbert $\mathcal{H}$ es exactamente $k$ (con $0 \le k \le N$).

	En particular:
	\begin{enumerate}
		\item Si todos los funcionales en su span son no acotados ($k=0$), entonces $\operatorname{codim}_\mathcal{H}(\overline{W_0}) = 0$ (el subespacio es denso en $\mathcal{H}$).
		\item Si todos los funcionales son acotados ($k=N$), entonces $\operatorname{codim}_\mathcal{H}(\overline{W_0}) = N$.
	\end{enumerate}
\end{theorem}

\begin{proof}
	Sin pérdida de generalidad, podemos realizar un cambio de base en el espacio de funcionales $\operatorname{span}\{\phi_1, \dots, \phi_N\}$ para obtener una nueva base de funcionales $\{\psi_1, \dots, \psi_N\}$ tal que los primeros $k$ funcionales ($\psi_1, \dots, \psi_k$) sean continuos en $\mathcal{H}$, y cualquier combinación lineal que involucre a los restantes $N-k$ funcionales ($\psi_{k+1}, \dots, \psi_N$) sea estrictamente no acotada. Notemos que la intersección de los núcleos no cambia: $W_0 = \bigcap_{j=1}^N \ker(\psi_j)$.

	Definimos el subespacio $M_c = \bigcap_{j=1}^k \ker(\psi_j)$ formado por los funcionales continuos.
	Aplicando el razonamiento constructivo sobre polinomios de grado finito $\mathbb{P}_n$, sabemos que para $n$ suficientemente grande, la restricción de estos $k$ funcionales a $\mathbb{P}_n$ es exhaustiva. Por el Teorema del Rango-Nulidad:
	$$ \dim(M_c \cap \mathbb{P}_n) = \dim(\mathbb{P}_n) - k $$
	Algebraicamente, podemos descomponer $\mathbb{P}_n = (M_c \cap \mathbb{P}_n) \oplus \operatorname{span}\{q_1, \dots, q_k\}$. Extendiendo a todo el espacio polinómico, obtenemos $\Poly = M_c \oplus \operatorname{span}\{q_1, \dots, q_k\}$.
	Al ser los $k$ funcionales continuos por hipótesis, $M_c$ es un subespacio cerrado en $\mathcal{H}$ ($\overline{M_c} = M_c$). La descomposición se hereda al espacio completo, $\mathcal{H} = M_c \oplus \operatorname{span}\{q_1, \dots, q_k\}$, demostrando que $\operatorname{codim}_\mathcal{H}(M_c) = k$.

	Ahora, analizamos la restricción de los funcionales no acotados restantes $\{\psi_{k+1}, \dots, \psi_N\}$ actuando \textit{dentro} del subespacio cerrado $M_c$.
	Como estas aplicaciones son no acotadas, el Lema de codimensión topológica finita garantiza que su restricción a $M_c$ (que tiene codimensión topológica finita $k$) sigue siendo no acotada.
	Por el mismo principio del Teorema para un funcional (Teorema \ref{thm:dicotomia_1}), la intersección de los núcleos de funcionales no acotados es densa en el espacio donde actúan.
	Por lo tanto, la clausura topológica de la intersección total dentro de $M_c$ es todo $M_c$:
	$$ \overline{W_0} = \overline{M_c \cap \bigcap_{j=k+1}^N \ker(\psi_j)} = M_c $$

	Finalmente, dado que la clausura topológica del subespacio total es exactamente $\overline{W_0} = M_c$, y ya habíamos probado que $\operatorname{codim}_\mathcal{H}(M_c) = k$, concluimos que:
	$$ \operatorname{codim}_\mathcal{H}(\overline{W_0}) = k \le N. $$

	Los casos extremos se siguen trivialmente: si $k=0$, $\overline{W_0} = \mathcal{H}$ y la codimensión topológica es 0; si $k=N$, $\overline{W_0} = W_0$ es cerrado y su codimensión topológica es máxima e igual a $N$.
\end{proof}

% -------------------------------------------------------------
\subsection{El Operador de Multiplicación y los Números de Gelfand}
% -------------------------------------------------------------

La conexión entre la teoría de polinomios ortogonales y la teoría espectral de operadores se realiza a través del operador de multiplicación por la variable independiente.

\begin{definition}[Operador Multiplicación] \label{def:operador_multiplicacion}
	Sea $\D: \Poly \to \Poly$ el operador lineal definido por:
	\begin{equation}
		(\D p)(z) = z \cdot p(z).
	\end{equation}
\end{definition}

\begin{remark}[Convención de notación]
	A partir de este punto distinguiremos sistemáticamente tres niveles de descripción. En el nivel de análisis funcional, $\D$ denota el operador abstracto actuando sobre $\mathcal{H}$ o sobre $\Poly$, y cantidades como $c_k(\D)$, $Q_k(\D)$ o $\|\D\|$ se refieren siempre a ese operador. En el nivel matricial, $\mathbf{D}$ denota la matriz infinita que representa a $\D$ en una base ortonormal de polinomios, y $\mathbf{D}_n$ su sección principal indexada por $\{0,\dots,n\}$ (por tanto, de tamaño $(n+1)\times(n+1)$). Finalmente, cuando identifiquemos un polinomio $p = \sum_{j=0}^{n} x_j \phi_j$ con su vector de coordenadas $x \in \mathbb{C}^{n+1}$, la norma euclídea de $x$ coincidirá con la norma de Sobolev de $p$ en ese subespacio finito-dimensional.
\end{remark}

En el marco de la teoría de ideales de operadores, los números de Gelfand constituyen la herramienta natural para cuantificar el comportamiento espectral direccional de un operador acotado.

\begin{definition}[Números de Gelfand]\label{def:numeros_gelfand}
	Sea $T \in \mathcal{L}(\mathcal{H})$ un operador acotado en un espacio de Hilbert. Para cada $k\ge 0$, el número de Gelfand de índice $k$ de $T$ se define como:
	\begin{equation}
		c_k(T) = \inf \{ \|T|_V\| : V \subset \mathcal{H}, \; \operatorname{codim}_{\mathcal{H}}(V) < k+1 \}.
	\end{equation}
	En particular, $c_0(T) = \|T\|$ es la norma global del operador. En este trabajo usaremos siempre la notación $c_k(\D)$ para los números de Gelfand del operador de multiplicación.
\end{definition}

Con el fin de analizar el impacto que las derivadas discretas tienen sobre la estructura espectral del operador multiplicación $\D(p) = zp$, se estudia el caso particular de un único punto en la derivada. Específicamente, se considera la perturbación de la medida continua, soportada en la circunferencia unidad $\mathbb{S}_1$, mediante la adición de un único término discreto en la derivada, evaluado en un punto $c_1 \in \C$. A lo largo de esta subsección, se asume que $c_1$ es un punto BPE.

\subsubsection{El Primer Número de Gelfand: La Norma Global}

Por definición, el primer índice de Gelfand (en la convención que empieza en $0$), $c_0(\D)$, corresponde al ínfimo sobre subespacios de codimensión topológica estrictamente menor que 1. El único subespacio que cumple esta condición es el espacio total $\mathcal{H}$. Por consiguiente:
$$ c_0(\D) = \|\D\|. $$
Este primer índice refleja el comportamiento global del operador: cuanto mayor es la contribución del término discreto a la norma de Sobolev, mayor es la dilatación de $c_0(\D)$ respecto al radio de la circunferencia.

Los cálculos explícitos más simples ya muestran dos rasgos del problema: la dependencia respecto del radio del soporte continuo y la sensibilidad a la posición del átomo discreto.

\begin{example}[Dos cálculos básicos para $c_0(\D)$]
	Consideremos primero el espacio de polinomios $\Poly$ dotado de la medida de Lebesgue normalizada en la circunferencia unidad $\mathbb{S}_1$ y un único punto en la derivada situado en el origen ($c_1 = 0$). El producto interno de Sobolev toma la forma
	\begin{equation}
		\ip{p}{q} = \int_{\mathbb{S}_1} p(z)\overline{q(z)} \frac{\abs{dz}}{2\pi} + p'(0)\overline{q'(0)}.
	\end{equation}
	Si $p(z) = \sum_{k=0}^n a_k z^k$, entonces $p'(0) = a_1$ y se obtiene
	\begin{equation} \label{eq:norma_in_0}
		\norm{p}^2 = \abs{a_0}^2 + 2\abs{a_1}^2 + \sum_{k=2}^n \abs{a_k}^2.
	\end{equation}
	Además, para $\D p(z) = z p(z)$ se tiene $(\D p)'(0) = a_0$, luego
	\begin{equation} \label{eq:norma_out_0}
		\norm{\D p}^2 = 2\abs{a_0}^2 + \abs{a_1}^2 + \sum_{k=2}^n \abs{a_k}^2.
	\end{equation}
	Por tanto,
	\begin{equation}
		\frac{\norm{\D p}^2}{\norm{p}^2} = 2 - \frac{3\abs{a_1}^2 + \sum_{k \ge 2} \abs{a_k}^2}{\norm{p}^2} \le 2,
	\end{equation}
	y el máximo se alcanza en las funciones constantes. En consecuencia, $c_0(\D) = \sqrt{2}$.

	Si ahora sustituimos $\mathbb{S}_1$ por la circunferencia $\mathbb{S}_R$ manteniendo el átomo en el origen, la ortogonalidad da
	\[\int_{\mathbb{S}_R} |z|^{2k} \frac{|dz|}{2\pi R} = R^{2k}.\]
	Las normas pasan a ser
	\begin{align}
		\norm{p}^2    & = \abs{a_0}^2 + (R^2 + 1)\abs{a_1}^2 + \sum_{k=2}^n \abs{a_k}^2 R^{2k}, \label{eq:norma_in_0_R}       \\
		\norm{\D p}^2 & = (R^2 + 1)\abs{a_0}^2 + R^4\abs{a_1}^2 + \sum_{k=2}^n \abs{a_k}^2 R^{2k+2}. \label{eq:norma_out_0_R}
	\end{align}
	De aquí se sigue que
	\begin{equation}
		\frac{\norm{\D p}^2}{\norm{p}^2} = (R^2 + 1) - \frac{(2R^2 + 1)\abs{a_1}^2 + \sum_{k \ge 2} \abs{a_k}^2 R^{2k}}{\norm{p}^2} \le R^2 + 1,
	\end{equation}
	con máximo de nuevo en las funciones constantes. Por tanto,
	\[
		c_0(\D) = \sqrt{R^2 + 1}.
	\]
\end{example}

El siguiente resultado muestra que, aun cuando $c_0(\D)$ refleje la perturbación discreta, al pasar a codimensión topológica uno se recupera exactamente la escala geométrica de la circunferencia.

\begin{theorem}[Valor exacto de $c_1$ para un punto]
	\label{thm:c2_punto}
	Sea $\mathcal{H} = \overline{\Poly}^{\|\cdot\|_S}$, donde
	\begin{equation}
		\langle p, q \rangle_S = \int_{\mathbb{S}_1} p(z)\overline{q(z)} \, d\mathbf{m}(z) + p'(c_1)\overline{q'(c_1)}, \qquad p,q \in \Poly,
	\end{equation}
	y supongamos que $c_1$ es un punto BPE para la medida continua $\mathbf m$ (equivalentemente, $\abs{c_1}<1$; véase \cite{duren1970theory}). Entonces $c_1(\D) = 1$.
\end{theorem}

\begin{proof}
	\textbf{Cota superior.} Como $c_1$ es un punto BPE y la evaluación de la derivada en $c_1$ es continua por el Lema~\ref{lema:ABPE_derivada}, el funcional
	\[
		\Psi(p) = p(c_1) + c_1 p'(c_1), \qquad p \in \Poly,
	\]
	se extiende continuamente a $\mathcal{H}$. Denotando de nuevo por $\Psi$ a dicha extensión, el subespacio $V_{opt} = \ker(\Psi)$ es cerrado y de codimensión topológica 1. Para $p \in V_{opt}$:
	\[
		\norm{\D p}^2 = \int_{\mathbb{S}_1} \abs{zp}^2 d\mu + \underbrace{\abs{(zp)'(c_1)}^2}_{=0} = \|p\|_{L^2}^2 \le \norm{p}^2,
	\]
	luego $\|\D|_{V_{opt}}\| \le 1$ y $c_1(\D) \le 1$.

	\textbf{Cota inferior.} Supóngase que existe $V$ de codimensión topológica 1 y $\epsilon > 0$ con $\norm{\D p}^2 \le (1-\epsilon)\norm{p}^2$ para todo $p \in V$. Desarrollando las normas:
	\[
		\epsilon \|p\|_{L^2}^2 + \abs{p(c_1) + c_1 p'(c_1)}^2 \le (1-\epsilon)\abs{p'(c_1)}^2, \quad \forall p \in V.
	\]
	Además, por el Lema~\ref{lema:ABPE_derivada}, el funcional de derivación puntual $\delta'_{c_1}: \mathcal{H} \to \mathbb{C}$ es continuo; en particular, su restricción $\delta'_{c_1}|_V : V \to \mathbb{C}$ también lo es. Como $V$ tiene codimensión topológica 1 en un espacio infinito-dimensional, sigue siendo infinito-dimensional. La aplicación lineal $\delta'_{c_1}|_V$ no puede ser inyectiva, pues una aplicación lineal inyectiva de un espacio vectorial infinito-dimensional en $\mathbb{C}$ forzaría $\dim(V) \leq \dim(\mathbb{C}) = 1$, contradicción. Por tanto, su núcleo es no trivial. Tomemos $q \in V$, $q \neq 0$, tal que $q'(c_1)=0$. Sustituyendo en la desigualdad anterior obtenemos
	\[
		\epsilon \|q\|_{L^2}^2 + \abs{q(c_1)}^2 \le 0.
	\]
	Esto es imposible, ya que ambos términos del lado izquierdo son no negativos y, además, $\|q\|_{L^2} \neq 0$ porque de otro modo se tendría $\|q\|_S^2 = \|q\|_{L^2}^2 + \abs{q'(c_1)}^2 = 0$, contradiciendo que $q \neq 0$. Por tanto $c_1(\D) \ge 1$.
\end{proof}

Cuando el punto de derivada se sitúa en la región exterior ($c \in I_{\mathrm{out}}$), la dificultad ya no es simplemente cuantitativa sino estructural. En efecto, aparece el funcional lineal
\[
	\psi_c:\Poly\to\C,
	\qquad
	\psi_c(p)=p(c)+c p'(c)=(\D p)'(c).
\]
En el caso modelo de Lebesgue en circunferencia y átomo $\delta$-singular ($\abs{c}\ge R$), la Proposición~\ref{prop:no_acotacion_modelo_exterior} demuestra que $\psi_c$ es no acotado respecto a la norma de Sobolev; por tanto, $\D$ no es acotado. Dicho de otro modo, en el régimen $\delta$-singular el problema deja de encajar de forma natural en la teoría de operadores acotados sobre $\mathcal{H}$.

\begin{proposition}[No acotación en el caso modelo $\delta$-singular]
	\label{prop:no_acotacion_modelo_exterior}
	Sea
	\[
		\langle p,q\rangle_S
		=
		\int_{\mathbb{S}_R} p(z)\overline{q(z)}\, d\mathbf m_{0;R}(z)
		+
		p'(c)\overline{q'(c)},
		\qquad p,q\in\Poly,
	\]
	con $\abs{c}\ge R$, y sea $\D p(z)=zp(z)$. Entonces el funcional
	\[
		\psi_c(p):=(\D p)'(c)=p(c)+cp'(c)
	\]
	es no acotado respecto de $\|\cdot\|_S$. En particular, $\D$ no es acotado y
	\[
		Q_0(\D)=\|\D\|_S=\infty.
	\]
\end{proposition}

\begin{proof}
	\textbf{Caso 1: $|c|>R$.}
	Definimos, para $n\ge 1$,
	\[
		Q_n(z):=(n+1)c\,z^n-nz^{n+1}.
	\]
	Un cálculo directo da
	\[
		Q_n'(z)=n(n+1)c\,z^{n-1}-n(n+1)z^n,
	\]
	y por tanto
	\[
		Q_n'(c)=0,
		\qquad
		Q_n(c)=c^{n+1}.
	\]
	Luego
	\[
		\|Q_n\|_S^2
		=
		\int_{\mathbb{S}_R}|Q_n(z)|^2\,d\mathbf m_{0;R}(z)
		+
		|Q_n'(c)|^2
		=
		\int_{\mathbb{S}_R}|Q_n(z)|^2\,d\mathbf m_{0;R}(z).
	\]
	Como los monomios $z^n$ y $z^{n+1}$ son ortogonales respecto a $\mathbf m_{0;R}$,
	\[
		\|Q_n\|_S^2
		=
		(n+1)^2|c|^2R^{2n}+n^2R^{2n+2}.
	\]
	Además,
	\[
		|\psi_c(Q_n)|^2
		=
		|Q_n(c)+cQ_n'(c)|^2
		=
		|c|^{2n+2}.
	\]
	Por tanto,
	\[
		\frac{|\psi_c(Q_n)|^2}{\|Q_n\|_S^2}
		=
		\frac{|c|^{2n+2}}{(n+1)^2|c|^2R^{2n}+n^2R^{2n+2}}
		\sim
		\frac{1}{n^2}\left(\frac{|c|}{R}\right)^{2n}
		\xrightarrow[n\to\infty]{}\infty,
	\]
	porque $|c|>R$. En consecuencia, $\psi_c$ no es acotado.

	\textbf{Caso 2: $|c|=R$.}
	Para $n\ge1$ definimos
	\[
		\alpha_n:=\frac{\sum_{k=1}^{n}k}{\sum_{k=1}^{n}k^2}=\frac{3}{2n+1},
		\qquad
		P_n(z):=\sum_{k=0}^{n}(1-\alpha_n k)\left(\frac{z}{c}\right)^k.
	\]
	Como
	\[
		cP_n'(c)=\sum_{k=1}^{n}k(1-\alpha_n k)=0,
	\]
	se tiene $P_n'(c)=0$ y, en consecuencia,
	\[
		\|P_n\|_S^2=\|P_n\|_{\mu}^2=
		\sum_{k=0}^{n}|1-\alpha_n k|^2.
	\]
	Desarrollando la suma,
	\[
		\|P_n\|_S^2
		=(n+1)-2\alpha_n\sum_{k=1}^{n}k+\alpha_n^2\sum_{k=1}^{n}k^2
		=\frac{(n+1)(n+2)}{2(2n+1)}.
	\]
	Además,
	\[
		\psi_c(P_n)=P_n(c)+cP_n'(c)=\sum_{k=0}^{n}(1-\alpha_n k)
		=\frac{(n+1)(n+2)}{2(2n+1)}.
	\]
	Por tanto,
	\[
		\frac{|\psi_c(P_n)|^2}{\|P_n\|_S^2}
		=
		\frac{(n+1)(n+2)}{2(2n+1)}
		\xrightarrow[n\to\infty]{}\infty.
	\]
	Luego $\psi_c$ también es no acotado cuando $|c|=R$.

	Si $\D$ fuera acotado, existiría $M>0$ tal que
	\[
		\|\D p\|_S\le M\|p\|_S
		\qquad (p\in\Poly).
	\]
	Pero
	\[
		|\psi_c(p)|
		=
		|(\D p)'(c)|
		\le
		\|\D p\|_S
		\le
		M\|p\|_S,
	\]
	lo que haría a $\psi_c$ acotado, contradicción. Luego $\D$ no es acotado y, por definición, $Q_0(\D)=\|\D\|_S=\infty$.
\end{proof}

Esta observación explica por qué el enfoque puramente topológico pierde eficacia en el caso exterior: la información relevante no está en la clausura de subespacios de $\mathcal{H}$, sino en la estructura algebraica de aquellos subespacios de polinomios donde la parte discreta del operador sí puede neutralizarse. Como las matrices truncadas de Hessenberg actúan precisamente sobre espacios finito-dimensionales de polinomios, resulta más natural introducir un índice de acotación definido directamente en $\Poly$.

\subsection{Un Nuevo Índice de Acotación Algebraica: El Índice $Q_k$}

En el régimen en que aparecen átomos sobre la frontera o en el exterior ($|c_k|\ge R$), el estudio topológico del operador multiplicación $\D$ en la completación de Hilbert $\mathcal{H}$ deja de ser la herramienta adecuada, porque los funcionales de evaluación que gobiernan la parte discreta no admiten extensión continua. Por ello, conviene reemplazar el punto de vista topológico por uno puramente algebraico que siga siendo compatible con las truncaciones finitas y con el análisis de los ceros.

Para superar este obstáculo, y dado que las matrices de Hessenberg truncadas operan exclusivamente sobre polinomios, resulta natural definir un nuevo índice que mida la acotación de forma puramente algebraica dentro del espacio pre-Hilbertiano $\Poly$.

\begin{definition}[Codimensión algebraica finita]\label{def:codim_alg}
	Para un subespacio vectorial $V\subset \Poly$, definimos
	\[
		\operatorname{codim}_{\Poly}(V):=\dim(\Poly/V)\in \mathbb N\cup\{\infty\}.
	\]
	Aquí $\dim(\cdot)$ denota dimensión algebraica (cardinal de una base de Hamel) \cite{hormander1989linearfunctionalanalysis,Aizpuru2009Apuntes}.
	Diremos que $V$ es de codimensión algebraica finita si
	\[
		\operatorname{codim}_{\Poly}(V)<\infty.
	\]
\end{definition}

\begin{proposition}[Caracterización por funcionales]\label{prop:finite-codim-kernels}
	Para un subespacio vectorial $V\subset \Poly$, son equivalentes:
	\begin{enumerate}
		\item $\operatorname{codim}_{\Poly}(V)<\infty$;
		\item existe un sistema finito de funcionales lineales $\Phi=(\phi_1,\dots,\phi_m)$
		      tal que
		      \[
			      V=\bigcap_{j=1}^m \ker \phi_j.
		      \]
	\end{enumerate}
	Además, en tal caso,
	\[
		\operatorname{codim}_{\Poly}(V)
		=
		\min\Bigl\{m:\exists\,\Phi=(\phi_1,\dots,\phi_m),\
		V=\bigcap_{j=1}^m\ker\phi_j\Bigr\}.
	\]
	Más precisamente, si $V=\bigcap_{j=1}^m\ker\phi_j$, entonces
	\[
		\operatorname{codim}_{\Poly}(V)=\dim\operatorname{span}\{\phi_1,\dots,\phi_m\}.
	\]
\end{proposition}

\begin{definition}[Índice Algebraico $Q_k$]\label{def:indice_Qk}
	Sea $\D:\Poly\to\Poly$ el operador multiplicación en el espacio pre-Hilbertiano de Sobolev. Para cada $k\ge0$, definimos
	\[
		Q_k(\D)
		:=\inf\{\|\D|_V\|_S:V\subset\Poly\ \text{admisible y}\ \operatorname{codim}_{\Poly}(V)<k+1\},
	\]
	donde
	\[
		\|\D|_V\|_S:=\sup_{\substack{p\in V\\p\ne0}}\frac{\|\D p\|_S}{\|p\|_S}.
	\]
	Por la convención anterior, $Q_0(\D)=\|\D\|_S$, admitiendo el valor $+\infty$.
\end{definition}

\begin{remark}[Lectura min-max]
	La familia admisible está construida mediante intersecciones finitas de hiperplanos en $\Poly$. Cuando dichos hiperplanos son ortogonales a polinomios, se recupera literalmente el esquema de Courant--Fischer; cuando provienen de funcionales no continuos en $\mathcal{H}$, la formulación sigue siendo válida dentro de $\Poly$.
\end{remark}

\begin{proposition}[Caso ortogonal]\label{prop:equiv_ortho_interseccion}
	Sea $F=(f_1,\dots,f_m)\subset \Poly$ y supongamos que
	$\langle\cdot,\cdot\rangle_S$ es un producto escalar en $\Poly$.
	Definimos
	\[
		V_{\perp}(F):=\{p\in\Poly:\langle p,f_j\rangle_S=0,\ j=1,\dots,m\}.
	\]
	Entonces
	\[
		V_{\perp}(F)=\bigcap_{j=1}^m \ker(\phi_j),
		\qquad
		\phi_j(p):=\langle p,f_j\rangle_S,
	\]
	y
	\[
		\operatorname{codim}_{\Poly}(V_\perp(F))
		=
		\dim\operatorname{span}\{f_1,\dots,f_m\}.
	\]
	En particular, si $f_1,\dots,f_m$ son linealmente independientes, entonces
	\[
		\operatorname{codim}_{\Poly}(V_\perp(F))=m.
	\]
\end{proposition}

La ventaja principal de este nuevo índice es que permite aislar las singularidades generadas por los átomos exteriores, revelando que la inestabilidad, aunque topológicamente densa, es algebraicamente de rango finito.

\begin{proposition}[Caso de Lebesgue previo al umbral $Q_k$]
	\label{prop:lebesgue_pre_umbral_Qk}
	Sea
	\[
		\langle p,q\rangle_S
		=
		\int_{\mathbb{S}_R} p(z)\overline{q(z)}\,d\mathbf m_{0;R}(z)
		+
		p'(c)\overline{q'(c)},
		\qquad p,q\in\Poly,
	\]
	con $|c|\ge R$, y sea $\D p(z)=zp(z)$. Entonces
	\[
		Q_0(\D)=\infty,
		\qquad
		Q_1(\D)\le R.
	\]
	En particular, en la circunferencia unidad ($R=1$) y para $|c|\ge1$, se tiene
	$Q_0(\D)=\infty$ y $Q_1(\D)\le1$.
\end{proposition}

\begin{proof}
	Por la Proposición~\ref{prop:no_acotacion_modelo_exterior}, el funcional
	\(
	\psi_c(p)=(\D p)'(c)=p(c)+cp'(c)
	\)
	es no acotado, y por tanto $\D$ no es acotado. En consecuencia,
	\(
	Q_0(\D)=\|\D\|_S=\infty.
	\)

	Para $Q_1$, tomamos
	\[
		W:=\{p\in\Poly:\psi_c(p)=0\}.
	\]
	Este subespacio es admisible y satisface $\operatorname{codim}_{\Poly}(W)=1$.
	Si $p\in W$, entonces
	\[
		\|\D p\|_S^2
		=
		\int_{\mathbb{S}_R}|zp(z)|^2\,d\mathbf m_{0;R}(z)
		+
		|\psi_c(p)|^2
		=
		\int_{\mathbb{S}_R}|zp(z)|^2\,d\mathbf m_{0;R}(z)
		\le
		R^2\int_{\mathbb{S}_R}|p(z)|^2\,d\mathbf m_{0;R}(z)
		\le
		R^2\|p\|_S^2.
	\]
	Luego $\|\D|_W\|\le R$, y al tomar ínfimo sobre subespacios
	admisibles con codimensión algebraica $<2$, se obtiene $Q_1(\D)\le R$.
\end{proof}

\begin{proposition}[Propiedades Básicas del Índice $Q_k$]\label{prop:propiedades_Qk}
	Para el operador multiplicación $\D$ definido sobre el espacio pre-Hilbertiano $\Poly$, el índice algebraico $Q_k(\D)$ forma una sucesión monótona no creciente cuyo primer término coincide exactamente con la norma global del operador. Es decir:
	$$ \|\D\|_S = Q_0(\D) \ge Q_1(\D) \ge \dots \ge Q_k(\D) \ge Q_{k+1}(\D) \ge \dots \ge 0 $$
	(admitiendo el valor $+\infty$ cuando $\D$ es no acotado).
\end{proposition}

\begin{proof}
	Para $k=0$, por la convención de la Definición~\ref{def:codim_alg}, solo aparece el conjunto vacío de restricciones y por tanto $V(\varnothing)=\Poly$. Así,
	$$ Q_0(\D) = \|\D|_{\Poly}\|_S = \|\D\|_S. $$

	Probemos ahora la monotonía. Sea $\Phi_k=(\phi_1,\dots,\phi_k)$ un sistema de $k$ funcionales y definamos
	$$V_k=V(\Phi_k)=\{p\in\Poly:\phi_j(p)=0,\ j=1,\dots,k\}.$$
	Si añadimos una restricción lineal más $\phi_{k+1}$, obtenemos
	$$V_{k+1}=V(\phi_1,\dots,\phi_k,\phi_{k+1})=\{p\in V_k:\phi_{k+1}(p)=0\}\subseteq V_k.$$
	Por inclusión de subespacios,
	$$\|\D|_{V_{k+1}}\|_S\le \|\D|_{V_k}\|_S.$$
	Tomando ínfimo sobre todos los sistemas de $k+1$ y $k$ restricciones, respectivamente, se concluye
	$$Q_{k+1}(\D)\le Q_k(\D).$$
	La no negatividad es inmediata al ser ínfimo de normas operatoriales (posiblemente infinitas).
\end{proof}

\begin{definition}[Índice Topológico $S_k$]\label{def:indice_Sk}
	Sea $\D:\Poly\to\Poly$ el operador multiplicación y sea $k\ge 0$. Definimos
	\[
		S_k(\D)
		:=
		\inf\!\left\{\,\|\D\|_{S,\overline{V}}:
		V\subset\Poly\ \text{subespacio vectorial},\ \operatorname{codim}_{\mathcal{H}}(\overline{V})<k+1\right\},
	\]
	donde $\overline{V}$ denota la clausura de $V$ en $\mathcal{H}$ y
	\[
		\|\D\|_{S,\overline{V}}
		:=
		\sup_{\substack{p\in V\\ p\neq 0}}\frac{\|\D p\|_S}{\|p\|_S}
		\in [0,+\infty].
	\]
	Si $\D$ extiende continuamente a $\overline{V}$, esta cantidad coincide con la norma operatorial usual $\|\D|_{\overline{V}}\|$.
\end{definition}

\begin{proposition}[Subaditividad y Perturbación Acotada]\label{prop:Qk_perturbacion}
	Sea $\D$ un operador en $\Poly$ y $\mathcal{E}$ un operador acotado en el mismo espacio ($\|\mathcal{E}\|_S < \infty$). Entonces, para todo $k \ge 0$:
	$$ Q_k(\D + \mathcal{E}) \le Q_k(\D) + \|\mathcal{E}\|_S $$
\end{proposition}

\begin{proof}
	Sea $V \subset \Poly$ un subespacio admisible con $\operatorname{codim}_{\Poly}(V) < k+1$. Por la desigualdad triangular de la norma de operadores, la restricción de la suma satisface:
	$$ \|(\D + \mathcal{E})|_V\|_S \le \|\D|_V\|_S + \|\mathcal{E}|_V\|_S $$
	Dado que la norma de una restricción nunca puede superar la norma global del operador, sabemos que $\|\mathcal{E}|_V\|_S \le \|\mathcal{E}\|_S$. Sustituyendo:
	$$ \|(\D + \mathcal{E})|_V\|_S \le \|\D|_V\|_S + \|\mathcal{E}\|_S $$
	Tomando el ínfimo sobre todos los subespacios admisibles con $\operatorname{codim}_{\Poly}(V)<k+1$ a ambos lados de la desigualdad, obtenemos directamente:
	$$ Q_k(\D + \mathcal{E}) \le Q_k(\D) + \|\mathcal{E}\|_S $$
\end{proof}

Para establecer el umbral algebraico exacto de estabilización del operador de multiplicación, fijamos el siguiente marco general. Consideramos el producto interno
\[
	\langle p,q\rangle_S
	=
	\int_{\mathbb{S}_R(0)} p(z)\overline{q(z)} \, d\mathbf{m}_{0;R}(z)
	+
	\sum_{k=1}^N p'(c_k)\,\overline{q'(c_k)},
	\qquad p,q\in\Poly,
\]
y el operador de multiplicación $\D:\Poly\to\Poly$ dado por $(\D p)(z)=z\,p(z)$. Separamos el conjunto de índices en
\[
	I_{\mathrm{in}}:=\{k\in\{1,\dots,N\}: |c_k|<R\},
	\qquad
	I_{\mathrm{out}}:=\{k\in\{1,\dots,N\}: |c_k|\ge R\}.
\]
Recordemos que $m=|I_{\mathrm{out}}|$ (Definición~\ref{def:class_points_new}).
Para cada $k\in I_{\mathrm{out}}$, definimos el funcional lineal asociado a la derivada del operador de multiplicación evaluada en $c_k$:
\[
	\psi_k:\Poly\to\mathbb C,
	\qquad
	\psi_k(p):=(\D p)'(c_k)=p(c_k)+c_kp'(c_k).
\]

\begin{remark}[Transición de notación]
	En lo sucesivo, para el funcional indexado por $k \in I_{\mathrm{out}}$ usaremos la notación $\psi_k$ en lugar de $\psi_{c_k}$.
\end{remark}

\begin{remark}[No acotación de $\psi_k$ en frontera y exterior]
	En este marco (medida de Lebesgue normalizada en circunferencia), para cada $k\in I_{\mathrm{out}}$ se tiene $|c_k|\ge R$. Por la Proposición~\ref{prop:no_acotacion_modelo_exterior}, el funcional $\psi_k$ es no acotado respecto de $\|\cdot\|_S$. En consecuencia, no es necesario imponer esta no acotación como hipótesis adicional.
\end{remark}

Dividiremos el análisis del umbral en varias partes para facilitar su lectura: la construcción de un subespacio estable en el caso acotado y la demostración de la explosión para codimensión estricta en el caso infinito. Finalmente, reuniremos ambos resultados.

\begin{proposition}[Fase 1: Caso acotado y átomos interiores]\label{prop:umbral_fase1}
	En el marco anteriormente descrito, existe un subespacio admisible
	\[
		V_{\mathrm{out}}:=\bigcap_{k\in I_{\mathrm{out}}}\ker(\psi_k),
	\]
	con $\operatorname{codim}_{\Poly}(V_{\mathrm{out}})\le m$, tal que $\D|_{V_{\mathrm{out}}}$ es acotado. En particular, $Q_m(\D)<\infty$.
\end{proposition}

\begin{proof}
	Como $V_{\mathrm{out}}$ es intersección finita de núcleos de funcionales lineales, es admisible.
	Además,
	\[
		\operatorname{codim}_{\Poly}(V_{\mathrm{out}})\le m.
	\]

	Sea ahora $p\in V_{\mathrm{out}}$. Entonces, para todo $k\in I_{\mathrm{out}}$,
	\[
		(\D p)'(c_k)=0.
	\]
	Por tanto, todos los términos discretos correspondientes a $I_{\mathrm{out}}$ quedan
	neutralizados en la norma de Sobolev de $\D p$.

	Si $k\in I_{\mathrm{in}}$, entonces $|c_k|<R$. En ese caso, por las estimaciones
	de Cauchy en el interior, existen constantes $A_k,B_k<\infty$ tales que
	\[
		|p(c_k)|\le A_k\|p\|_S,
		\qquad
		|p'(c_k)|\le B_k\|p\|_S.
	\]
	Luego
	\[
		|(\D p)'(c_k)|
		=
		|p(c_k)+c_kp'(c_k)|
		\le
		(A_k+|c_k|B_k)\|p\|_S.
	\]

	Además, como la parte continua está soportada en la circunferencia de radio $R$,
	\[
		\|zp\|_{L^2(\mathbf{m}_{0;R})}\le R\,\|p\|_{L^2(\mathbf{m}_{0;R})}\le R\,\|p\|_S.
	\]
	De aquí,
	\[
		\|\D p\|_S^2
		=
		\|zp\|_{L^2(\mathbf{m}_{0;R})}^2
		+
		\sum_{k=1}^N |(\D p)'(c_k)|^2
	\]
	\[
		\le
		\left(
		R^2+\sum_{k\notin I_{\mathrm{out}}}(A_k+|c_k|B_k)^2
		\right)\|p\|_S^2.
	\]
	Esto prueba que $\D|_{V_{\mathrm{out}}}$ es acotado. Como
	$\operatorname{codim}_{\Poly}(V_{\mathrm{out}})\le m$, por definición de $Q_m$ se obtiene
	\[
		Q_m(\D)<\infty.
	\]
\end{proof}

\begin{proposition}[Fase 2: Caso infinito y explosión]\label{prop:umbral_fase2}
	En el marco general establecido, para todo subespacio admisible $V\subset\Poly$ con $\operatorname{codim}_{\Poly}(V)<m$, la restricción $\D|_V$ no es acotada. Más aún, existe una sucesión $(p_n)\subset V$ tal que $\|p_n\|_S\to 0$ pero $\|\D p_n\|_S\ge 1$, y por tanto $Q_{m-1}(\D)=\infty$.
\end{proposition}

\begin{proof}
	Por definición,
	\[
		Q_{m-1}(\D)
		=
		\inf\{\|\D|_V\|_S:\ V\subset\Poly\ \text{admisible y}\ \operatorname{codim}_{\Poly}(V)<m\}.
	\]
	Basta, por tanto, demostrar que toda restricción admisible con
	$\operatorname{codim}_{\Poly}(V)<m$ es no acotada.

	Fijemos un subespacio admisible $V\subset \Poly$ tal que
	\[
		\operatorname{codim}_{\Poly}(V)=r<m.
	\]
	Demostraremos que $\D|_V$ es no acotado.

	\smallskip
	\noindent
	\emph{(a) Independencia lineal de los funcionales $\psi_k$.}

	Los puntos $\{c_k\}_{k\in I_{\mathrm{out}}}$ son distintos. Por interpolación de Hermite,
	para cada $j\in I_{\mathrm{out}}$ existe un polinomio $h_j$ tal que
	\[
		h_j(c_j)=1,\qquad h_j'(c_j)=0,
	\]
	y
	\[
		h_j(c_\ell)=0,\qquad h_j'(c_\ell)=0,
		\qquad \ell\in I_{\mathrm{out}},\ \ell\neq j.
	\]
	Entonces
	\[
		\psi_\ell(h_j)=\delta_{\ell j}.
	\]
	Por consiguiente, $\{\psi_k\}_{k\in I_{\mathrm{out}}}$ es un sistema linealmente independiente,
	y en particular
	\[
		\dim \operatorname{span}\{\psi_k:\ k\in I_{\mathrm{out}}\}=m.
	\]

	\smallskip
	\noindent
	\emph{(b) Construcción de paquetes diagonales de norma pequeña.}

	Como cada $\psi_j$ es no acotado respecto de $\|\cdot\|_S$, por el
	Lema~\ref{lem:normalizacion_funcional_no_acotado}, para todo $n\ge1$
	existe $x_{j,n}\in \Poly$ tal que
	\[
		\|x_{j,n}\|_S<\frac1n,
		\qquad
		\psi_j(x_{j,n})=1.
	\]

	Ahora refinamos la interpolación de Hermite: para cada $j\in I_{\mathrm{out}}$ elegimos
	un polinomio $g_j$ tal que
	\[
		g_j(c_j)=1,\qquad g_j'(c_j)=0,
	\]
	y, para todo $\ell\neq j$,
	\[
		g_j(c_\ell)=0,\qquad g_j'(c_\ell)=0,
	\]
	imponiendo estas condiciones en todos los puntos discretos
	$\{c_1,\dots,c_N\}$.

	Definimos
	\[
		u_{j,n}:=g_j\,x_{j,n}.
	\]

	Observemos primero que, para todo $p\in\Poly$ y todo $\ell=1,\dots,N$,
	por la regla del producto,
	\[
		(g_jp)'(c_\ell)
		=
		g_j'(c_\ell)p(c_\ell)+g_j(c_\ell)p'(c_\ell).
	\]
	Como $g_j'(c_\ell)=0$ para todo $\ell$, obtenemos
	\[
		(g_jp)'(c_\ell)=g_j(c_\ell)p'(c_\ell).
	\]
	Por consiguiente,
	\[
		\sum_{\ell=1}^N |(g_jp)'(c_\ell)|^2
		\le
		\Bigl(\max_{1\le \ell\le N}|g_j(c_\ell)|^2\Bigr)
		\sum_{\ell=1}^N |p'(c_\ell)|^2.
	\]

	Por otra parte, en la circunferencia $|z|=R$ se tiene
	\[
		|g_j(z)p(z)|
		\le
		\|g_j\|_{L^\infty(|z|=R)}\,|p(z)|,
	\]
	y por tanto
	\[
		\int_{|z|=R}|g_jp|^2\,d\mathbf{m}_{0;R}
		\le
		\|g_j\|_{L^\infty(|z|=R)}^2
		\int_{|z|=R}|p|^2\,d\mathbf{m}_{0;R}.
	\]

	Reuniendo ambas estimaciones, si definimos
	\[
		C_j:=
		\max\Bigl\{
		\|g_j\|_{L^\infty(|z|=R)},
		\max_{1\le \ell\le N}|g_j(c_\ell)|
		\Bigr\},
	\]
	obtenemos
	\[
		\|g_jp\|_S\le C_j\|p\|_S,
		\qquad p\in\Poly.
	\]
	Es decir, la multiplicación por $g_j$ es un operador acotado sobre
	$(\Poly,\|\cdot\|_S)$.

	En particular,
	\[
		\|u_{j,n}\|_S=\|g_jx_{j,n}\|_S\le C_j\|x_{j,n}\|_S<\frac{C_j}{n}\xrightarrow[n\to\infty]{}0.
	\]

	Además, para $\ell\in I_{\mathrm{out}}$,
	\[
		u_{j,n}(c_\ell)=g_j(c_\ell)x_{j,n}(c_\ell),
	\]
	y
	\[
		u'_{j,n}(c_\ell)=g_j'(c_\ell)x_{j,n}(c_\ell)+g_j(c_\ell)x'_{j,n}(c_\ell)
		=g_j(c_\ell)x'_{j,n}(c_\ell),
	\]
	porque $g_j'(c_\ell)=0$.

	Si $\ell\neq j$, entonces $g_j(c_\ell)=0$, luego
	\[
		u_{j,n}(c_\ell)=0,
		\qquad
		u'_{j,n}(c_\ell)=0,
	\]
	y por tanto
	\[
		\psi_\ell(u_{j,n})=u_{j,n}(c_\ell)+c_\ell u'_{j,n}(c_\ell)=0.
	\]

	Si $\ell=j$, entonces $g_j(c_j)=1$ y $g_j'(c_j)=0$, luego
	\[
		u_{j,n}(c_j)=x_{j,n}(c_j),
		\qquad
		u'_{j,n}(c_j)=x'_{j,n}(c_j),
	\]
	y por tanto
	\[
		\psi_j(u_{j,n})
		=
		u_{j,n}(c_j)+c_j u'_{j,n}(c_j)
		=
		x_{j,n}(c_j)+c_jx'_{j,n}(c_j)
		=
		\psi_j(x_{j,n})=1.
	\]
	En resumen,
	\[
		\psi_\ell(u_{j,n})=\delta_{\ell j},
		\qquad \ell,j\in I_{\mathrm{out}}.
	\]

	\smallskip
	\noindent
	\emph{(c) Imposición simultánea de las restricciones algebraicas de $V$.}

	Como $V$ es admisible y de codimensión algebraica $r$, por la caracterización por
	funcionales existe un sistema linealmente independiente
	\[
		\lambda_1,\dots,\lambda_r:\Poly\to\mathbb C
	\]
	tal que
	\[
		V=\bigcap_{i=1}^r \ker(\lambda_i).
	\]

	Para cada $n$, consideremos la matriz
	\[
		A_n:=\bigl(\lambda_i(u_{j,n})\bigr)_{
		\substack{1\le i\le r\\ j\in I_{\mathrm{out}}}}
		\in \mathbb C^{r\times m}.
	\]
	Aquí aparece la idea clave: intersecar simultáneamente los núcleos de
	$\lambda_1,\dots,\lambda_r$ equivale a escoger coeficientes en el núcleo de $A_n$.
	Como $r<m$, el núcleo de $A_n$ es no trivial. Elegimos
	\[
		\beta^{(n)}=(\beta_j^{(n)})_{j\in I_{\mathrm{out}}}\in \ker(A_n),
		\qquad
		\|\beta^{(n)}\|_{\ell^2}=1.
	\]
	Definimos
	\[
		p_n:=\sum_{j\in I_{\mathrm{out}}}\beta_j^{(n)}u_{j,n}.
	\]
	Entonces, para cada $i=1,\dots,r$,
	\[
		\lambda_i(p_n)=\sum_{j\in I_{\mathrm{out}}}\beta_j^{(n)}\lambda_i(u_{j,n})=0,
	\]
	luego $p_n\in V$.

	Además,
	\[
		\|p_n\|_S
		\le
		\sum_{j\in I_{\mathrm{out}}} |\beta_j^{(n)}|\,\|u_{j,n}\|_S
		\le
		\sqrt m \max_{j\in I_{\mathrm{out}}}\|u_{j,n}\|_S
		\xrightarrow[n\to\infty]{}0.
	\]

	\smallskip
	\noindent
	\emph{(d) Cota inferior uniforme para $\|\D p_n\|_S$.}

	Definimos el funcional
	\[
		\psi^{(n)}:=\sum_{j\in I_{\mathrm{out}}}\overline{\beta_j^{(n)}}\,\psi_j.
	\]
	Entonces, usando la diagonalidad anterior,
	\[
		\psi^{(n)}(p_n)
		=
		\sum_{j,\ell\in I_{\mathrm{out}}}
		\overline{\beta_j^{(n)}}\beta_\ell^{(n)}\psi_j(u_{\ell,n})
		=
		\sum_{j\in I_{\mathrm{out}}}|\beta_j^{(n)}|^2
		=1.
	\]
	Por otra parte, por Cauchy--Schwarz,
	\[
		|\psi^{(n)}(p)|
		\le
		\Bigl(\sum_{j\in I_{\mathrm{out}}}|\beta_j^{(n)}|^2\Bigr)^{1/2}
		\Bigl(\sum_{j\in I_{\mathrm{out}}}|\psi_j(p)|^2\Bigr)^{1/2}
		\le
		\|\D p\|_S,
		\qquad p\in\Poly,
	\]
	ya que los términos $|\psi_j(p)|^2=|(\D p)'(c_j)|^2$ forman parte de la suma discreta
	de $\|\D p\|_S^2$.

	Aplicando esto a $p_n$ obtenemos
	\[
		1=|\psi^{(n)}(p_n)|\le \|\D p_n\|_S.
	\]
	Por tanto,
	\[
		\frac{\|\D p_n\|_S}{\|p_n\|_S}\xrightarrow[n\to\infty]{}\infty.
	\]
	Esto prueba que $\D|_V$ es no acotado.

	Como $V$ era arbitrario entre los subespacios admisibles con codimensión algebraica
	$<m$, toda restricción que interviene en la definición de $Q_{m-1}(\D)$ es no
	acotada. En consecuencia,
	\[
		Q_{m-1}(\D)=\infty.
	\]
\end{proof}

% [MODIFICADO] Enunciado original refería a ``átomos exteriores'' con condición |c_k|≥R; ahora generalizado a condición δ-singular. La demostración es la misma, usando la definición funcional de I_out.
\begin{theorem}[Umbral algebraico exacto de estabilización, versión $\delta$-singular]
	\label{thm:umbral_exacto_Qm_nobpe}
	\label{thm:umbral_delta_singular}
	Sea $\langle\cdot,\cdot\rangle_S$ un producto de Sobolev discreto general
	con $N$ átomos $c_1,\dots,c_N$, y sea
	$m=|I_{\mathrm{out}}|$ el número de puntos $\delta$-singulares. Asumamos
	además que la parte continua del producto induce una norma
	$\|\cdot\|_{L^2(\mu)}$ tal que la multiplicación por $z$ es acotada en
	$L^2(\mu)$ con norma a lo sumo $R$. Entonces
	\[
		Q_j(\D)=\infty\ \text{ para }\ 0\le j\le m-1,
		\qquad
		Q_m(\D)\le R.
	\]
	En particular, el número de puntos $\delta$-singulares $m = |I_{\mathrm{out}}|$ es el umbral algebraico exacto de estabilización del operador de multiplicación $\D$.
\end{theorem}

\begin{proof}
	Por la Proposición~\ref{prop:propiedades_Qk}, la sucesión $\{Q_k(\D)\}_{k\ge0}$ es monótona no creciente:
	\[
		Q_0(\D)\ge Q_1(\D)\ge\cdots\ge Q_{m-1}(\D).
	\]
	Por la Proposición~\ref{prop:umbral_fase2}, $Q_{m-1}(\D)=\infty$, de donde se deduce para todo $0\le j\le m-1$ que
	\[
		Q_j(\D)\ge Q_{m-1}(\D)=\infty,
	\]
	y por tanto $Q_j(\D)=\infty$ para $j=0,1,\dots,m-1$. Finalmente, por la Proposición~\ref{prop:umbral_fase1}, $Q_m(\D)<\infty$, lo cual concluye la demostración.
\end{proof}

\begin{corollary}[Lectura BPE del umbral en el caso Lebesgue]\label{cor:lectura_BPE_lebesgue}
	Si la parte continua del producto de Sobolev es la medida de Lebesgue
	normalizada en $\mathbb{S}_R$, entonces, por el
	Teorema~\ref{thm:equivalencia_lebesgue},
	\[
		I_{\mathrm{out}}=\{k:c_k\text{ es no-BPE para }L^2(\mathbf m_{0;R})\}
		=\{k:|c_k|\ge R\},
	\]
	y el Teorema~\ref{thm:umbral_delta_singular} recupera el resultado
	clásico de umbral exacto controlado por el número de puntos no-BPE.
\end{corollary}

\begin{corollary}[Aplicación BPE: caso puramente exterior]\label{cor:umbral_m_puramente_exterior}
	Si en el Teorema~\ref{thm:umbral_exacto_Qm_nobpe} todos los átomos discretos son $\delta$-singulares (equivalentemente, $|c_k|\ge R$ para todo $k$), entonces
	\[
		m=N,
	\]
	se cumple
	\[
		Q_m(\D)\le R.
	\]
	En particular, el umbral algebraico en el índice $m$ es finito y está controlado explícitamente por el radio del soporte continuo.
\end{corollary}

\begin{proof}
	En este caso
	\[
		I_{\mathrm{out}}=\{1,\dots,N\},
		\qquad
		m=N.
	\]
	Consideremos el subespacio admisible
	\[
		V_*:=\bigcap_{k=1}^N \ker(\psi_k),
		\qquad
		\psi_k(p):=(\D p)'(c_k)=p(c_k)+c_kp'(c_k).
	\]
	Por construcción,
	\[
		\operatorname{codim}_{\Poly}(V_*)\le N=m.
	\]

	Sea ahora $p\in V_*$. Entonces, para todo $k=1,\dots,N$,
	\[
		(\D p)'(c_k)=0.
	\]
	Por tanto, todos los términos discretos de la norma de Sobolev de $\D p$ se anulan y queda
	\[
		\|\D p\|_S^2
		=
		\int_{|z|=R} |z\,p(z)|^2\,d\mathbf{m}_{0;R}(z).
	\]
	Como $|z|=R$ en el soporte continuo, se obtiene
	\[
		\|\D p\|_S^2
		=
		R^2\int_{|z|=R}|p(z)|^2\,d\mathbf{m}_{0;R}(z)
		\le
		R^2\|p\|_S^2.
	\]
	En consecuencia,
	\[
		\|\D|_{V_*}\|_S\le R.
	\]
	Finalmente, como $V_*$ es admisible y $\operatorname{codim}_{\Poly}(V_*)\le m$, por definición de $Q_m(\D)$,
	\[
		Q_m(\D)\le \|\D|_{V_*}\|_S\le R.
	\]
\end{proof}

\begin{definition}[Ideal anulador de un subespacio]
	Sea $V\subset \Poly$ un subespacio vectorial. Definimos
	\[
		I(V):=\{A\in \Poly: A\,\Poly\subset V\}.
	\]
\end{definition}

\begin{theorem}[Subespacio estable canónico y anulador polinómico]
	\label{thm:Vout_anulador}
	Sea
	\[
		\langle p,q\rangle_S
		=
		\int_{|z|=R} p(z)\,\overline{q(z)}\,d\mathbf{m}_{0;R}(z)
		+
		\sum_{k=1}^N p'(c_k)\,\overline{q'(c_k)},
		\qquad p,q\in \Poly,
	\]
	donde $R>0$ y los puntos $c_1,\dots,c_N$ son distintos. Sea
	\[
		(\D p)(z)=z\,p(z),
	\]
	Definamos
	\[
		I_{\mathrm{out}}:=\{k\in\{1,\dots,N\}: |c_k|\ge R\}.
	\]
	Recordemos que $m=|I_{\mathrm{out}}|$ (Definición~\ref{def:class_points_new}).
	Para cada $k\in I_{\mathrm{out}}$ consideremos el funcional lineal
	\[
		\psi_k:\Poly\to\mathbb C,
		\qquad
		\psi_k(p):=(\D p)'(c_k)=p(c_k)+c_kp'(c_k),
	\]
	y el subespacio
	\[
		V_{\mathrm{out}}:=\bigcap_{k\in I_{\mathrm{out}}}\ker(\psi_k).
	\]
	Entonces se verifican las siguientes propiedades:

	\begin{enumerate}
		\item $V_{\mathrm{out}}$ es admisible y
		      \[
			      \operatorname{codim}_{\Poly}(V_{\mathrm{out}})=m.
		      \]

		\item La restricción $\D|_{V_{\mathrm{out}}}$ es acotada respecto de la norma
		      de Sobolev. Más precisamente, existe una constante $C_{\mathrm{out}}<\infty$ tal que
		      \[
			      \|\D p\|_S\le C_{\mathrm{out}}\|p\|_S,
			      \qquad p\in V_{\mathrm{out}}.
		      \]

		\item Si definimos
		      \[
			      A_{\mathrm{out}}(z):=\prod_{k\in I_{\mathrm{out}}}(z-c_k)^2,
		      \]
		      entonces
		      \[
			      A_{\mathrm{out}}\Poly\subset V_{\mathrm{out}}.
		      \]
		      En particular,
		      \[
			      I(V_{\mathrm{out}})\neq \{0\}.
		      \]

		\item Más aún,
		      \[
			      I(V_{\mathrm{out}})=(A_{\mathrm{out}}).
		      \]
		      Es decir, el ideal anulador de $V_{\mathrm{out}}$ está generado exactamente por
		      $A_{\mathrm{out}}$.
	\end{enumerate}
\end{theorem}

\begin{proof}
	\textbf{(1) Admisibilidad y codimensión exacta.}
	Por definición,
	\[
		V_{\mathrm{out}}=\bigcap_{k\in I_{\mathrm{out}}}\ker(\psi_k),
	\]
	luego $V_{\mathrm{out}}$ es admisible.

	Probemos ahora que los funcionales $\{\psi_k\}_{k\in I_{\mathrm{out}}}$ son linealmente
	independientes. Fijado $j\in I_{\mathrm{out}}$, por interpolación de Hermite existe un
	polinomio $h_j\in\Poly$ tal que
	\[
		h_j(c_j)=1,\qquad h_j'(c_j)=0,
	\]
	y
	\[
		h_j(c_\ell)=0,\qquad h_j'(c_\ell)=0,
		\qquad \ell\in I_{\mathrm{out}},\ \ell\neq j.
	\]
	Entonces, para todo $\ell\in I_{\mathrm{out}}$,
	\[
		\psi_\ell(h_j)=h_j(c_\ell)+c_\ell h_j'(c_\ell)=\delta_{\ell j}.
	\]
	Por tanto, los funcionales $\psi_k$ son linealmente independientes y
	\[
		\dim\operatorname{span}\{\psi_k:\ k\in I_{\mathrm{out}}\}=m.
	\]
	Aplicando la caracterización de la codimensión algebraica por funcionales, obtenemos
	\[
		\operatorname{codim}_{\Poly}(V_{\mathrm{out}})
		=
		\dim\operatorname{span}\{\psi_k:\ k\in I_{\mathrm{out}}\}
		=
		m.
	\]

	\medskip
	\textbf{(2) Acotación de $\D$ sobre $V_{\mathrm{out}}$.}
	Sea $p\in V_{\mathrm{out}}$. Entonces, para todo $k\in I_{\mathrm{out}}$,
	\[
		(\D p)'(c_k)=\psi_k(p)=0.
	\]
	Es decir, todos los términos discretos correspondientes a los puntos exteriores/$\delta$-singulares
	desaparecen en la norma de Sobolev de $\D p$.

	Si $k\notin I_{\mathrm{out}}$, entonces $|c_k|<R$. En ese caso, por las estimaciones
	interiores habituales, existen constantes $A_k,B_k<\infty$ tales que
	\[
		|p(c_k)|\le A_k\|p\|_S,
		\qquad
		|p'(c_k)|\le B_k\|p\|_S.
	\]
	Luego
	\[
		|(\D p)'(c_k)|
		=
		|p(c_k)+c_kp'(c_k)|
		\le
		(A_k+|c_k|B_k)\|p\|_S.
	\]

	Además, como $|z|=R$ en el soporte continuo,
	\[
		\|zp\|_{L^2(\mathbf{m}_{0;R})}\le R\|p\|_{L^2(\mathbf{m}_{0;R})}\le R\|p\|_S.
	\]
	Por consiguiente,
	\[
		\|\D p\|_S^2
		=
		\|zp\|_{L^2(\mathbf{m}_{0;R})}^2
		+
		\sum_{k=1}^N |(\D p)'(c_k)|^2
	\]
	\[
		=
		\|zp\|_{L^2(\mathbf{m}_{0;R})}^2
		+
		\sum_{k\notin I_{\mathrm{out}}} |(\D p)'(c_k)|^2
	\]
	\[
		\le
		\left(
		R^2+\sum_{k\notin I_{\mathrm{out}}}(A_k+|c_k|B_k)^2
		\right)\|p\|_S^2.
	\]
	Esto prueba la acotación de $\D|_{V_{\mathrm{out}}}$.

	\medskip
	\textbf{(3) Existencia de anulador polinómico.}
	Sea
	\[
		A_{\mathrm{out}}(z)=\prod_{k\in I_{\mathrm{out}}}(z-c_k)^2.
	\]
	Entonces, para cada $k\in I_{\mathrm{out}}$,
	\[
		A_{\mathrm{out}}(c_k)=0,
		\qquad
		A_{\mathrm{out}}'(c_k)=0.
	\]
	Sea ahora $q\in\Poly$. Entonces
	\[
		(A_{\mathrm{out}}q)(c_k)=0,
		\qquad
		(A_{\mathrm{out}}q)'(c_k)=0,
		\qquad k\in I_{\mathrm{out}}.
	\]
	Por tanto,
	\[
		\psi_k(A_{\mathrm{out}}q)
		=
		(A_{\mathrm{out}}q)(c_k)+c_k(A_{\mathrm{out}}q)'(c_k)
		=
		0.
	\]
	Esto vale para todo $k\in I_{\mathrm{out}}$, luego
	\[
		A_{\mathrm{out}}q\in V_{\mathrm{out}}
		\qquad \forall q\in\Poly.
	\]
	En consecuencia,
	\[
		A_{\mathrm{out}}\Poly\subset V_{\mathrm{out}},
	\]
	y por tanto
	\[
		A_{\mathrm{out}}\in I(V_{\mathrm{out}}),
		\qquad
		I(V_{\mathrm{out}})\neq\{0\}.
	\]

	\medskip
	\textbf{(4) Cálculo exacto del ideal anulador.}
	Ya sabemos que
	\[
		(A_{\mathrm{out}})\subset I(V_{\mathrm{out}}).
	\]
	Veamos la inclusión opuesta.

	Sea $B\in I(V_{\mathrm{out}})$. Entonces
	\[
		Bq\in V_{\mathrm{out}}
		\qquad \forall q\in\Poly.
	\]
	En particular, para cada $k\in I_{\mathrm{out}}$ y todo $q\in\Poly$,
	\[
		\psi_k(Bq)=0.
	\]

	Fijemos $k\in I_{\mathrm{out}}$. Tomemos primero
	\[
		q(z)=z-c_k.
	\]
	Entonces $q(c_k)=0$ y $q'(c_k)=1$. Aplicando la regla del producto,
	\[
		(Bq)'(c_k)=B'(c_k)q(c_k)+B(c_k)q'(c_k)=B(c_k).
	\]
	Por tanto,
	\[
		0=\psi_k(Bq)=(Bq)(c_k)+c_k(Bq)'(c_k)=c_kB(c_k).
	\]
	Como $|c_k|\ge R$ y $R>0$, se tiene $c_k\neq 0$. Luego
	\[
		B(c_k)=0.
	\]

	Ahora tomamos $q\equiv 1$. Entonces
	\[
		0=\psi_k(B)=B(c_k)+c_kB'(c_k)=c_kB'(c_k),
	\]
	y de nuevo, como $c_k\neq 0$, se sigue que
	\[
		B'(c_k)=0.
	\]

	Hemos probado que cada $c_k$, con $k\in I_{\mathrm{out}}$, es raíz doble de $B$.
	Por consiguiente,
	\[
		A_{\mathrm{out}}(z)=\prod_{k\in I_{\mathrm{out}}}(z-c_k)^2
	\]
	divide a $B$. Es decir,
	\[
		B\in (A_{\mathrm{out}}).
	\]
	Luego
	\[
		I(V_{\mathrm{out}})\subset (A_{\mathrm{out}}).
	\]
	Junto con la inclusión opuesta, obtenemos
	\[
		I(V_{\mathrm{out}})=(A_{\mathrm{out}}).
	\]
\end{proof}

\begin{corollary}[Existencia de subespacio estable con anulador no trivial]
	\label{cor:existencia_V_estable_anulador}
	Bajo las hipótesis del Teorema~\ref{thm:Vout_anulador}, existe un subespacio admisible
	$V\subset \Poly$ tal que
	\[
		\operatorname{codim}_{\Poly}(V)\le m,
		\qquad
		\|\D|_V\|_S<\infty,
		\qquad
		I(V)\neq \{0\}.
	\]
	De hecho, puede tomarse
	\[
		V=V_{\mathrm{out}}=\bigcap_{k\in I_{\mathrm{out}}}\ker(\psi_k),
		\qquad
		\psi_k(p)=p(c_k)+c_kp'(c_k).
	\]
	En particular,
	\[
		Q_m(\D)\le \|\D|_{V_{\mathrm{out}}}\|_S<\infty.
	\]
\end{corollary}

\begin{proof}
	Basta tomar
	\[
		V=V_{\mathrm{out}}.
	\]
	Por el Teorema~\ref{thm:Vout_anulador},
	\[
		\operatorname{codim}_{\Poly}(V_{\mathrm{out}})=m,
		\qquad
		\|\D|_{V_{\mathrm{out}}}\|_S<\infty,
		\qquad
		I(V_{\mathrm{out}})=(A_{\mathrm{out}})\neq\{0\}.
	\]
	Finalmente, como $V_{\mathrm{out}}$ es admisible y $\operatorname{codim}_{\Poly}(V_{\mathrm{out}})=m< m+1$,
	la definición de $Q_m(\D)$ da
	\[
		Q_m(\D)\le \|\D|_{V_{\mathrm{out}}}\|_S<\infty.
	\]
\end{proof}

\begin{corollary}[Anulador canónico asociado al umbral exterior]
	\label{cor:anulador_canonico_umbral}
	Bajo las hipótesis del Teorema~\ref{thm:Vout_anulador}, el polinomio
	\[
		A_{\mathrm{out}}(z)=\prod_{k\in I_{\mathrm{out}}}(z-c_k)^2
	\]
	satisface
	\[
		A_{\mathrm{out}}\Poly\subset V_{\mathrm{out}},
	\]
	donde $V_{\mathrm{out}}$ es un subespacio admisible de codimensión exacta $m$
	sobre el cual $\D$ es acotado.
\end{corollary}

\begin{proof}
	Es una reformulación inmediata de los apartados (1), (2) y (3) del
	Teorema~\ref{thm:Vout_anulador}.
\end{proof}

% [CORRECCIÓN ESTRUCTURAL] Teorema duplicado eliminado; su contenido queda consolidado en el Teorema~\ref{thm:Vout_anulador}.

% [CORRECCIÓN ESTRUCTURAL] Eliminado Corolario~cor:existencia_subespacio_anulador y remark asociado; su contenido queda cubierto por el Corolario~\ref{cor:existencia_V_estable_anulador} y el Teorema~\ref{thm:Vout_anulador}.

\subsubsection{Consistencia con los Números de Gelfand Topológicos}

Para que el índice $Q_k$ sea una generalización válida y robusta, es imperativo demostrar que, en ausencia de patologías topológicas (es decir, cuando el operador sí es acotado de forma natural en $\mathcal{H}$), el índice puramente algebraico coincide exactamente con los números de Gelfand definidos en la completación de Hilbert.

\begin{theorem}[Equivalencia en el Caso Acotado]\label{thm:equivalencia_Qk_Gelfand}
	Si el operador multiplicación $\D$ está acotado en el espacio de Hilbert $\mathcal{H}$ (por ejemplo, si todas las masas discretas son puntos BPE), entonces para todo $k \ge 0$:
	$$ Q_k(\D) = c_k(\D) $$
\end{theorem}

\begin{proof}
	Procederemos probando ambas desigualdades. Recordemos que $c_k(\D)$ calcula el ínfimo sobre subespacios \textit{cerrados} $M \subset \mathcal{H}$ de codimensión topológica $m < k+1$.

	\textbf{Paso 1: $Q_k(\D) \le c_k(\D)$.}
	Sea $M \subset \mathcal{H}$ un subespacio cerrado de codimensión topológica $m < k+1$. Por propiedades elementales de Análisis Funcional, $M$ puede escribirse como intersección de $m$ hiperplanos cerrados, esto es, $M=\bigcap_{j=1}^m\ker\Phi_j$ con $\Phi_1,\dots,\Phi_m\in\mathcal{H}^*$ linealmente independientes. Tomando intersección con $\Poly$, obtenemos
	\[
		V:=M\cap\Poly=\bigcap_{j=1}^m\ker(\Phi_j|_{\Poly}).
	\]
	En efecto, para $p\in\Poly$ se tiene
	\[
		p\in M\cap\Poly
		\iff p\in\Poly\ \text{y}\ \Phi_j(p)=0\ \forall j
		\iff p\in\bigcap_{j=1}^m\ker(\Phi_j|_{\Poly}),
	\]
	lo que justifica la igualdad anterior.

	Además, las restricciones $\Phi_j|_{\Poly}$ son linealmente independientes: si
	\[
		\sum_{j=1}^m \alpha_j\,\Phi_j|_{\Poly}=0,
	\]
	entonces el funcional continuo $\Psi:=\sum_{j=1}^m \alpha_j\Phi_j\in\mathcal{H}^*$ se anula en $\Poly$. Como $\Poly$ es denso en $\mathcal{H}$ y $\Psi$ es continuo, se sigue que $\Psi\equiv 0$ en todo $\mathcal{H}$. Por independencia lineal de $\Phi_1,\dots,\Phi_m$, necesariamente $\alpha_1=\cdots=\alpha_m=0$.

	Por la Proposición~\ref{prop:finite-codim-kernels}, se concluye que $\operatorname{codim}_{\Poly}(V)=m<k+1$.
	Probemos ahora que $\overline V=M$. La inclusión $\overline V\subset M$ es inmediata, ya que $V\subset M$ y $M$ es cerrado. Para la inclusión recíproca, sea $x\in M$. Como $\Poly$ es denso en $\mathcal H$, existe una sucesión $p_n\in \Poly$ tal que $p_n\to x$ en $\mathcal H$. Definimos
	\[
		a_n:=\bigl(\Phi_1(p_n),\dots,\Phi_m(p_n)\bigr)\in\mathbb C^m.
	\]
	Como $x\in M=\bigcap_{j=1}^m\ker\Phi_j$, se tiene $\Phi_j(x)=0$ para todo $j$, y por continuidad de cada $\Phi_j$ resulta $a_n\to 0$ en $\mathbb C^m$.

	Consideremos la aplicación lineal
	\[
		T:\Poly\to\mathbb C^m,\qquad T(p):=\bigl(\Phi_1(p),\dots,\Phi_m(p)\bigr).
	\]
	Las restricciones $\Phi_j|_{\Poly}$ son linealmente independientes, luego $\dim\operatorname{Im}T=m$; en particular, $T$ es sobreyectiva. Elegimos polinomios $u_1,\dots,u_m\in\Poly$ tales que
	\[
		T(u_i)=e_i\quad (i=1,\dots,m),
	\]
	donde $\{e_i\}$ es la base canónica de $\mathbb C^m$. Definimos el operador lineal continuo finito-dimensional
	\[
		R:\mathbb C^m\to\Poly,\qquad R(b_1,\dots,b_m):=\sum_{i=1}^m b_i u_i,
	\]
	de modo que $T\circ R=\mathrm{Id}_{\mathbb C^m}$.

	Para cada $n$, ponemos
	\[
		v_n:=p_n-R(a_n)\in\Poly.
	\]
	Entonces
	\[
		T(v_n)=T(p_n)-T(R(a_n))=a_n-a_n=0,
	\]
	luego $v_n\in\ker T=V$. Además,
	\[
		\|v_n-x\|_S\le \|p_n-x\|_S+\|R(a_n)\|_S\xrightarrow[n\to\infty]{}0,
	\]
	porque $p_n\to x$, $a_n\to0$ y $R$ es continuo. Por tanto $x\in\overline V$, y se concluye $M\subset\overline V$. En consecuencia,
	\[
		\overline V=M.
	\]

	Como $\D\in\mathcal L(\mathcal H)$ y $V$ es denso en $M$, se obtiene
	\[
		\|\D|_V\|=\|\D|_M\|.
	\]
	Al ser $V$ un subespacio admisible en la definición de $Q_k$ (pues $\operatorname{codim}_{\Poly}(V)=m<k+1$), resulta
	\[
		Q_k(\D)\le \|\D|_V\|=\|\D|_M\|.
	\]
	Tomando ínfimo sobre todos los subespacios cerrados $M\subset\mathcal H$ con $\operatorname{codim}_{\mathcal H}(M)<k+1$, concluimos
	\[
		Q_k(\D)\le c_k(\D).
	\]

	\textbf{Paso 2: $c_k(\D) \le Q_k(\D)$.}
	Sea $V \subset \Poly$ un subespacio admisible con $\operatorname{codim}_{\Poly}(V) = m < k+1$. Consideremos su clausura topológica $M = \overline{V}$ en el espacio de Hilbert $\mathcal{H}$.
	Como $\operatorname{codim}_{\Poly}(V)=m$, existe $F\subset\Poly$ con $\dim(F)=m$ tal que
	\[
		\Poly=V\oplus F.
	\]
	Tomando clausuras en $\mathcal{H}$ y usando que $F$ es finito-dimensional (luego cerrado),
	\[
		\mathcal{H}=\overline{\Poly}=\overline{V\oplus F}=\overline{V+F}=\overline V+F=M+F.
	\]
	La última igualdad se justifica porque siempre $\overline{A+B}\subset \overline A+\overline B$, y aquí además $\overline F=F$.

	Ahora consideremos la aplicación lineal
	\[
		q:F\to \mathcal H/M,\qquad q(f)=f+M.
	\]
	Como $\mathcal H=M+F$, toda clase $h+M\in\mathcal H/M$ puede escribirse como $(m+f)+M=f+M=q(f)$; luego $q$ es sobreyectiva. En consecuencia,
	\[
		\dim(\mathcal{H}/M)=\dim(\operatorname{Im}q)\le \dim(F)=m,
	\]
	es decir,
	\[
		\operatorname{codim}_{\mathcal{H}}(M)\le m<k+1.
	\]
	Así, $M$ es admisible en la definición de $c_k(\D)$.

	Además, como $V$ es denso en $M$ y $\D$ es continuo en $\mathcal{H}$:
	\[
		\|\D|_V\|\le \|\D|_M\|.
	\]
	Para la desigualdad inversa, fijado $\varepsilon>0$, existe $x\in M$, $\|x\|_S=1$, tal que
	\[
		\|\D x\|_S>\|\D|_M\|-\varepsilon.
	\]
	Tomamos $v_n\in V$ con $v_n\to x$ en $\|\cdot\|_S$. Entonces
	\[
		\|v_n\|_S\to 1,
		\qquad
		\|\D v_n\|_S\to \|\D x\|_S.
	\]
	En particular, fijado $\varepsilon>0$, para $n$ suficientemente grande se cumple simultáneamente
	\[
		\|\D v_n\|_S>\|\D|_M\|-\varepsilon,
		\qquad
		\|v_n\|_S<1+\varepsilon.
	\]
	Por tanto,
	\[
		\frac{\|\D v_n\|_S}{\|v_n\|_S}
		>
		\frac{\|\D|_M\|-\varepsilon}{1+\varepsilon}.
	\]
	Haciendo $\varepsilon\downarrow 0$, esta cota tiende a $\|\D|_M\|$.
	y por definición de norma operatorial sobre $V$,
	\[
		\|\D|_V\|\ge \|\D|_M\|.
	\]
	Junto con la desigualdad opuesta $\|\D|_V\|\le \|\D|_M\|$, concluimos
	\[
		\|\D|_V\|=\|\D|_M\|.
	\]
	Dado que $M$ es admisible para $c_k$, se tiene
	\[
		c_k(\D)\le \|\D|_M\|=\|\D|_V\|.
	\]
	Tomando ínfimo sobre todos los $V\subset\Poly$ con $\operatorname{codim}_{\Poly}(V)<k+1$:
	\[
		c_k(\D)\le \inf_{\operatorname{codim}_{\Poly}(V)<k+1}\|\D|_V\|=Q_k(\D).
	\]

	Al cumplirse ambas desigualdades, se concluye la identidad $Q_k(\D) = c_k(\D)$.
\end{proof}

% [CORRECCIÓN ESTRUCTURAL] Transición de cierre hacia el análisis matricial.
Los resultados de esta sección establecen el marco funcional que relaciona los números de Gelfand algebraicos con la geometría de los puntos singulares. En particular, hemos demostrado que la restricción del operador de multiplicación a un subespacio de codimensión finita recupera la acotación, y que el índice $Q_k(\D)$ proporciona la cota exacta para la explosión del espectro.

% =============================================================
% CAPÍTULO 3: ANÁLISIS MATRICIAL Y COMPORTAMIENTO ASINTÓTICO
% 3.1 Representación Matricial y Secciones Truncadas
% 3.2 Valores Singulares y Operadores Truncados
% 3.3 Estabilización Algebraica de los Valores Singulares
% 3.4 Matrices de Momentos de Sobolev
% =============================================================

% !TEX root = ../main.tex

\chapter{Análisis Matricial y Comportamiento Asintótico}\label{sec:matricial}

El estudio asintótico del operador multiplicación $\D$ se lleva a cabo a través de sus aproximaciones de rango finito, denominadas secciones principales. El paso del espacio de Hilbert abstracto al álgebra de polinomios finitos permite emplear herramientas de álgebra lineal numérica para analizar las matrices que gobiernan la disposición de los ceros.

% [CORRECCIÓN ESTRUCTURAL] Reforzada la motivación del paso a matrices.
Para convertir estas estimaciones funcionales en predicciones sobre la distribución de los ceros de los polinomios ortogonales, necesitamos una representación finito-dimensional del operador. Las secciones truncadas de la matriz de multiplicación en la base ortonormal de polinomios de Sobolev constituyen el puente natural entre el análisis espectral abstracto y los cálculos efectivos.

En esta sección, la discusión matricial se hará siempre sobre la representación de $\D$ en la base ortonormal $\{\phi_n\}_{n\ge 0}$. Así, $\mathbf{D}$ denotará la matriz infinita de esa representación y $\mathbf{D}_n$ su compresión a $\operatorname{span}\{\phi_0,\dots,\phi_n\}$; cuando usemos vectores $x \in \mathbb{C}^{n+1}$, estos serán los vectores de coordenadas de polinomios en dicha base.

% -------------------------------------------------------------
\section{Representación Matricial y Secciones Truncadas}
% -------------------------------------------------------------

\begin{defn}[Sección principal de orden $n$]\label{def:seccion_truncada}
	Sea $A = (a_{ij})_{i,j \geq 0}$ una matriz infinita. Para cada $n \in \mathbb{N}$, la \emph{sección principal de orden $n$} de $A$ es la submatriz finita
	\[
		A_n = (a_{ij})_{0 \leq i,j \leq n} \in \mathcal{M}_{n+1}(\mathbb{C}).
	\]
	Denotamos por $\mathbf{D}_n$ la sección de orden $n$ de la matriz $\mathbf{D}$ que representa al operador $\D$ en la base de polinomios ortonormales $\{\phi_k\}$.
\end{defn}

\begin{defn}[Norma espectral de una sección truncada]\label{def:norma_espectral_truncada}
	Para cada $n \in \mathbb{N}$, la norma espectral de la matriz $\mathbf{D}_n \in \mathcal{M}_{n+1}(\mathbb{C})$ se define por
	\[
		\|\mathbf{D}_n\|_2 = \sup_{\substack{x \in \mathbb{C}^{n+1} \\ x \neq 0}} \frac{\|\mathbf{D}_n x\|_2}{\|x\|_2}.
	\]
	Bajo la identificación entre polinomios de $\operatorname{span}\{\phi_0,\dots,\phi_n\}$ y sus vectores de coordenadas, esta norma coincide con la norma operatorial de la compresión finito-dimensional de $\D$ a dicho subespacio.
\end{defn}

La representación matricial del operador $\D$ en la base ortonormal $\{\phi_n\}$ posee una estructura algebraica especialmente conveniente. Esta representación recibe el nombre de matriz de Hessenberg.

\begin{defn}[Matriz de Hessenberg y Jacobi] \label{def:hessenberg}
	La matriz asociada al operador $\D$ en la base $\{\phi_n\}$ es una matriz de Hessenberg superior infinita, denotada por $\mathbf{D} = (d_{k,j})_{k,j\ge 0}$. Sus elementos vienen dados por:
	\begin{equation}
		d_{k,j} = \langle \D\phi_j, \phi_k \rangle_S = \langle z \phi_j, \phi_k \rangle_S.
	\end{equation}
	Debido a la relación de recurrencia que satisfacen los polinomios ortogonales, esta matriz es tridiagonal (Matriz de Jacobi) en el caso real, o Hessenberg superior con subdiagonal no nula en el caso general.
\end{defn}

\begin{prop}[Estructura de Hessenberg de $\mathbf{D}$]\label{prop:hessenberg}
	Para todos los índices $i, j \geq 0$ se verifica $d_{i,j} = 0$ si $i > j+1$. En consecuencia, $\mathbf{D}$ es una matriz de Hessenberg superior infinita:
	\[
		\mathbf{D} = \begin{pmatrix}
			d_{0,0} & d_{0,1} & d_{0,2} & d_{0,3} & \cdots \\
			d_{1,0} & d_{1,1} & d_{1,2} & d_{1,3} & \cdots \\
			0       & d_{2,1} & d_{2,2} & d_{2,3} & \cdots \\
			0       & 0       & d_{3,2} & d_{3,3} & \cdots \\
			\vdots  & \vdots  & \vdots  & \vdots  & \ddots
		\end{pmatrix}.
	\]
	El patrón de ceros es consecuencia directa de que $\deg(z\phi_j) = j+1$, de modo que la expansión de $z\phi_j$ en la base solo involucra $\phi_0, \dots, \phi_{j+1}$.
\end{prop}

Las secciones $\mathbf{D}_n$ de esta matriz poseen las siguientes propiedades relevantes para el análisis de los ceros:

\begin{itemize}
	\item Si el producto interno procede de una medida real en un intervalo, la relación de recurrencia de tres términos hace que $\mathbf{D}$ sea tridiagonal y simétrica (matriz de Jacobi).
	\item El determinante característico $|\mathbf{D}_n - z\mathbf{I}_{n+1}|$ es un polinomio proporcional a $\phi_{n+1}(z)$; en particular, ambos comparten exactamente los mismos ceros.
\end{itemize}

% -------------------------------------------------------------
\section{Valores Singulares y Operadores Truncados}
% -------------------------------------------------------------

Para el estudio del comportamiento asintótico de los operadores truncados, se emplea la teoría de los números de aproximación o valores singulares.

\begin{defn}[Valores Singulares]
	Sea $T: \mathcal{H}_1 \to \mathcal{H}_2$ un operador lineal compacto entre espacios de Hilbert. En este trabajo adoptamos indexación desde $0$: los valores singulares de $T$, denotados por $\{s_n(T)\}_{n \geq 0}$, son las raíces cuadradas de los autovalores del operador autoadjunto no negativo $T^*T$ y se ordenan de forma decreciente:
	\begin{equation}
		s_0(T) \geq s_1(T) \geq \dots \geq 0.
	\end{equation}
	Equivale a la convención clásica $\{s_{n+1}(T)\}_{n\ge0}$. Alternativamente, coinciden con los números de aproximación reindexados, definidos como la distancia del operador a los operadores de rango finito:
	\begin{equation}
		a_n(T) = \inf \{ \|T - F\| : \text{rango}(F) < n+1 \}.
	\end{equation}
\end{defn}

La tasa de decaimiento de estos números proporciona información precisa sobre la velocidad de convergencia de las aproximaciones finitas y la distribución asintótica de los ceros.

\begin{defn}[Valores singulares de las secciones: notación $\sigma_{k,n}$]\label{def:sigma_kn}
	Los valores singulares de la sección $\mathbf{D}_n$, ordenados de forma decreciente, se denotan
	\[
		\sigma_{0,n} \geq \sigma_{1,n} \geq \cdots \geq \sigma_{n,n} \geq 0.
	\]
	En particular, $\sigma_{0,n} = \|\mathbf{D}_n\|_2$ es la norma espectral (norma operatorial inducida por la norma euclídea) de $\mathbf{D}_n$. Por el principio min-max de Courant-Fischer,
	\[
		\sigma_{1,n}^2 = \min_{\substack{V \subset \mathbb{C}^{n+1} \\ \operatorname{codim}(V) = 1}} \max_{\substack{x \in V \\ \|x\| = 1}} \|\mathbf{D}_n x\|^2.
	\]
\end{defn}

\begin{defn}[Índice asintótico singular $\sigma_k$]\label{def:indice_Sigma_k}
	Para cada $k\ge 0$, si existe el límite en la recta real extendida, definimos
	\[
		\sigma_k := \lim_{n\to\infty} \sigma_{k,n} \in [0,+\infty].
	\]
\end{defn}

\begin{prop}[Coincidencia total de índices en el caso acotado]\label{prop:coincidencia_total_acotado_pendiente}
	Supongamos que $\D\in\mathcal{L}(\mathcal{H})$, que para cada $k\ge0$ existe el límite
	\[
		\sigma_k=\lim_{n\to\infty}\sigma_{k,n},
	\]
	y que, además, las truncaciones son estables en el sentido de que
	\[
		Q_k(\D)\le \liminf_{n\to\infty}\sigma_{k,n},
		\qquad k\ge0.
	\]
	Entonces, para todo $k\ge0$,
	\[
		\sigma_k=S_k(\D)=Q_k(\D)=c_k(\D).
	\]
\end{prop}

\begin{proof}
	Fijemos $k\ge 0$. Por la Proposición~\ref{prop:sigma_qk}, para todo $n$,
	\[
		\sigma_{k,n}\le Q_k(\D).
	\]
	Pasando al límite y usando la existencia de $\sigma_k$,
	\[
		\sigma_k\le Q_k(\D).
	\]

	Por la hipótesis de estabilidad de truncaciones,
	\[
		Q_k(\D)\le \liminf_{n\to\infty}\sigma_{k,n}=\sigma_k.
	\]
	Luego
	\[
		\sigma_k=Q_k(\D).
	\]

	Como $\D\in\mathcal L(\mathcal H)$, el Teorema~\ref{thm:equivalencia_Qk_Gelfand} da
	\[
		Q_k(\D)=c_k(\D).
	\]

	Probamos ahora que $S_k(\D)=c_k(\D)$. Para cualquier
	$V\subset\Poly$ con $\operatorname{codim}_{\mathcal H}(\overline V)<k+1$, si
	$M:=\overline V$, como $\D$ es acotado y $V$ es denso en $M$,
	\[
		\|\D\|_{S,\overline V}=\|\D|_M\|\ge c_k(\D).
	\]
	Tomando ínfimo en la definición de $S_k$ se obtiene
	\[
		S_k(\D)\ge c_k(\D).
	\]

	Recíprocamente, sea $M\subset\mathcal H$ cerrado con
	$\operatorname{codim}_{\mathcal H}(M)<k+1$ y definamos
	$V:=M\cap\Poly$. En el Paso 1 del
	Teorema~\ref{thm:equivalencia_Qk_Gelfand} se probó que $\overline V=M$; por tanto,
	\[
		\|\D\|_{S,\overline V}=\|\D|_M\|.
	\]
	Al tomar ínfimo sobre tales $M$,
	\[
		S_k(\D)\le c_k(\D).
	\]
	Concluimos que
	\[
		S_k(\D)=c_k(\D)=Q_k(\D)=\sigma_k.
	\]
	Esto vale para todo $k\ge0$.
\end{proof}

\begin{prop}[Comportamiento del valor singular de índice $0$]\label{prop:convergencia_sigma}
	Sea $\D$ el operador multiplicación por $z$ y sea $\mathbf{D}_n$ la matriz de su compresión al subespacio $\operatorname{span}\{\phi_0,\dots,\phi_n\}$. Entonces:
	\begin{enumerate}
		\item si $\D$ extiende a un operador acotado en $\mathcal{H}$, se tiene
		      \[
			      \lim_{n \to \infty} \sigma_{0,n} = \lim_{n \to \infty} \|\mathbf{D}_n\|_2 = \|\D\|;
		      \]
		\item si $\D$ no es acotado en $\mathcal{H}$, entonces $\|\mathbf{D}_n\|_2 \to +\infty$.
	\end{enumerate}
\end{prop}

En particular, cuando la medida continua tiene soporte acotado con radio máximo $R$ y el operador multiplicación es acotado, se obtiene $\|\D\| = R$ y, por tanto, $\lim_{n \to \infty} \sigma_{0,n} = R$.

\section{Estabilización Algebraica de los Valores Singulares}

Aunque la Proposición~\ref{prop:convergencia_sigma} y la topología de $\mathcal{H}$ sugieren que $\|\mathbf{D}_n\|_2 \to \infty$ en el caso $\delta$-singular, la estructura algebraica del espacio de polinomios $\Poly$ restringe la dimensionalidad de esta divergencia.

\begin{thm}[Acotación del valor singular de índice $1$]
	Sea $\langle \cdot, \cdot \rangle_S$ un producto de Sobolev con medida continua de radio esencial $R$ y un único átomo de derivada en $c \in I_{\mathrm{out}}$. Entonces, el valor singular de índice $1$ de las secciones principales satisface:
	$$ \limsup_{n \to \infty} \sigma_{1,n} \le R. $$
\end{thm}

\begin{proof}
	Consideremos el principio min-max de Courant-Fischer (Definición~\ref{def:sigma_kn}). Las secciones $\mathbf{D}_n$ actúan sobre subespacios de dimensión finita de polinomios. Podemos restringir el operador a un subespacio puramente algebraico $W \subset \Poly$ diseñado para neutralizar la inestabilidad.

	Definimos $W = \{ p \in \Poly : (zp)'(c) = 0 \} = \{ p \in \Poly : p(c) + c p'(c) = 0 \}$. Al ser el núcleo de un único funcional lineal, su codimensión algebraica en $\Poly$ es exactamente 1. Para cualquier $p \in W$, la componente discreta de la imagen se anula:
	$$ \|\D p\|_S^2 = \|zp\|_{L^2(\mu)}^2 + |(zp)'(c)|^2 = \|zp\|_{L^2(\mu)}^2 \le R^2 \|p\|_{L^2(\mu)}^2 \le R^2 \|p\|_S^2. $$
	Al intersecar $W$ con el subespacio generado por los primeros $n+1$ polinomios ortonormales, obtenemos un subespacio de codimensión algebraica a lo sumo 1 en $\mathbb{C}^{n+1}$ donde la norma de la matriz $\mathbf{D}_n$ está acotada por $R$. Por el principio min-max, esto fuerza a que $\sigma_{1,n} \le R$ para todo $n$.
\end{proof}

Este resultado muestra que, en el caso de un único átomo exterior, la divergencia asintótica de las secciones truncadas queda confinada a una sola dirección: $\sigma_{0,n} \to \infty$, mientras que $\sigma_{1,n}$ permanece uniformemente acotado por $R$.

\begin{thm}[Control singular del escape exterior]\label{thm:control_singular_escape}
	Sea $\mu$ una medida de Borel positiva con soporte compacto tal que $\operatorname{supp}(\mu) \subseteq \overline{\mathbb{D}_R(0)}$ con $R < \infty$. Sea $\langle \cdot, \cdot \rangle_S$ un producto interno de Sobolev definido sobre $\Poly$ como:
	\begin{equation}\label{eq:sobolev_general_singular}
		\langle p, q \rangle_S = \int p(z)\overline{q(z)} d\mu(z) + \sum_{k=1}^N \lambda_k p^{(j_k)}(c_k) \overline{q^{(j_k)}(c_k)}
	\end{equation}
	donde $\lambda_k > 0$, $j_k \ge 0$, y $\{c_k\}_{k=1}^N \subset \mathbb{C}$ son puntos finitos. Entonces, para toda sección principal $\mathbf{D}_n$ con $n \ge N$, se verifica
	\[
		\sigma_{N,n} \le R.
	\]
	En particular, cada $\mathbf{D}_n$ admite una descomposición
	\[
		\mathbf{D}_n = \mathbf{M}_n + \mathbf{E}_n,
	\]
	con $\|\mathbf{M}_n\|_2 \le R$ y $\operatorname{rango}(\mathbf{E}_n) \le N$. Por consiguiente, la posible inestabilidad inducida por la perturbación discreta queda confinada, en sentido singular, a lo sumo a $N$ direcciones.
\end{thm}

\begin{proof}
	Analizaremos la acción del operador multiplicación $\D: p \mapsto zp$ sobre la norma inducida por (\ref{eq:sobolev_general_singular}). Para cualquier polinomio $p \in \Poly$, la norma de su imagen es:
	\begin{equation}
		\|\D p\|_S^2 = \int |zp(z)|^2 d\mu(z) + \sum_{k=1}^N \lambda_k \left| (zp)^{(j_k)}(c_k) \right|^2
	\end{equation}
	Aplicando la regla de Leibniz para la derivada de un producto, la evaluación discreta resulta en:
	\begin{equation}\label{eq:leibniz_eval}
		(zp)^{(j_k)}(c_k) = c_k p^{(j_k)}(c_k) + j_k p^{(j_k-1)}(c_k)
	\end{equation}
	Definimos un conjunto de $N$ funcionales lineales sobre $\Poly$ dados por $\phi_k(p) = c_k p^{(j_k)}(c_k) + j_k p^{(j_k-1)}(c_k)$ para $k=1,\dots,N$. Consideremos el subespacio algebraico $W$ definido como la intersección de sus núcleos:
	\begin{equation}
		W = \bigcap_{k=1}^N \ker(\phi_k) = \left\{ p \in \Poly : (zp)^{(j_k)}(c_k) = 0, \quad \forall k=1,\dots,N \right\}
	\end{equation}
	Al ser la intersección de $N$ hiperplanos, la codimensión algebraica de $W$ en $\Poly$ cumple $\operatorname{codim}_{\Poly}(W) \le N$.

	Para todo polinomio $p \in W$, el término discreto de $\|\D p\|_S^2$ se anula idénticamente por definición. Por tanto, acotando la integral sobre el soporte de $\mu$:
	\begin{equation}
		\|\D p\|_S^2 = \int_{\operatorname{supp}(\mu)} |z|^2 |p(z)|^2 d\mu(z) \le R^2 \int |p(z)|^2 d\mu(z) \le R^2 \|p\|_S^2
	\end{equation}
	Esto demuestra que la restricción del operador cumple $\|\D|_W\| \le R$.

	Sea $\mathbf{D}_n$ la matriz de la compresión de $\D$ al subespacio $\mathbb{P}_{n}$. Definimos $W_n = W \cap \mathbb{P}_{n}$. Por la fórmula de Grassmann, $\dim(W_n) \ge n+1 - N$. Identificando ahora cada polinomio $p \in \mathbb{P}_{n}$ con su vector de coordenadas $x \in \mathbb{C}^{n+1}$ en la base ortonormal, el subespacio $W_n$ induce un subespacio vectorial $\mathcal{W}_n \subset \mathbb{C}^{n+1}$ de dimensión al menos $n+1-N$.

	El valor singular de índice $N$ de $\mathbf{D}_n$, denotado como $\sigma_{N,n}$, se caracteriza mediante el teorema minimax de Courant-Fischer:
	\begin{equation}
		\sigma_{N,n} = \min_{\dim(S) = n+1-N} \max_{\substack{x \in S \\ x \neq 0}} \frac{\|\mathbf{D}_n x\|_2}{\|x\|_2} \le \max_{\substack{x \in \mathcal{V}_{0,n} \\ x \neq 0}} \frac{\|\mathbf{D}_n x\|_2}{\|x\|_2} \le R
	\end{equation}
	Dado que $\sigma_{N,n} \le R$ para todo $n \ge N$, podemos utilizar la descomposición en valores singulares. La matriz truncada se escribe como $\mathbf{D}_n = \mathbf{U} \mathbf{\Sigma} \mathbf{V}^*$, donde $\mathbf{\Sigma}$ es diagonal y contiene los valores singulares en orden decreciente.
	Separando los primeros $N$ valores singulares (que pueden ser mayores que $R$) del resto, la matriz se descompone de forma exacta como una suma:
	$$ \mathbf{D}_n = \mathbf{M}_n + \mathbf{E}_n $$
	donde $\mathbf{E}_n$ contiene únicamente la contribución de los $N$ primeros valores singulares (por lo que $\operatorname{rango}(\mathbf{E}_n) \le N$) y $\mathbf{M}_n$ contiene el resto de los valores singulares acotados por $\sigma_{N,n} \le R$. Por definición, esto garantiza que la norma espectral de esta última cumple $\|\mathbf{M}_n\|_2 \le R$.
\end{proof}

% [CORRECCIÓN ESTRUCTURAL] Se añade remark sobre generalidad del teorema anterior.
\begin{rmk}[Sobre la generalidad del Teorema~\ref{thm:control_singular_escape}]
	El Teorema~\ref{thm:control_singular_escape} admite derivadas de orden arbitrario $j_k$ en los átomos discretos. En el resto del trabajo, salvo indicación contraria, se asume $j_k=1$ para todo $k$, lo cual corresponde al producto de Sobolev estándar considerado en las Secciones~\ref{sec:intro} y~\ref{sec:espectral}.
\end{rmk}

\begin{rmk}[Límite del argumento matricial]
	El resultado anterior controla valores singulares, no el número exacto de autovalores exteriores. En efecto, si $\lambda$ es un autovalor de $\mathbf{D}_n$ con $|\lambda| > R$ y $x \neq 0$ es un autovector asociado, de la identidad
	\begin{equation}\label{eq:autovector_escape}
		x = (\lambda \mathbf{I} - \mathbf{M}_n)^{-1} \mathbf{E}_n x
	\end{equation}
	solo se deduce que $x$ pertenece a
	\[
		\operatorname{Im}\bigl((\lambda \mathbf{I} - \mathbf{M}_n)^{-1} \mathbf{E}_n\bigr),
	\]
	pero este subespacio depende de $\lambda$. Por ello, no es válido concluir que todos los autovectores correspondientes a autovalores exteriores queden contenidos en un único subespacio de dimensión $N$. Esta es precisamente la obstrucción que aparece en matrices no normales, como las truncaciones de Hessenberg consideradas aquí.
\end{rmk}

\begin{thm}[Umbral singular en el índice $m$ (enunciado de cierre)]\label{thm:umbral_singular_indice_m_enunciado}
	Sea $\langle\cdot,\cdot\rangle_S$ un producto de Sobolev con componente continua soportada en
	$\overline{\mathbb D_R(0)}$, y sea $m=|I_{\mathrm{out}}|$ el número de puntos $\delta$-singulares.
	Suponemos además que estamos en régimen puramente $\delta$-singular en la parte discreta, es decir,
	que todos los átomos discretos del producto corresponden a índices de $I_{\mathrm{out}}$.
	(En el caso modelo de Lebesgue normalizada en circunferencia, esta condición implica la
	no acotación de los funcionales $\psi_k(p)=(\D p)'(c_k)$ por la
	Proposición~\ref{prop:no_acotacion_modelo_exterior}, aplicada punto a punto.)
	Entonces se verifica el patrón umbral
	\[
		\lim_{n\to\infty}\sigma_{j,n}=+\infty,
		\qquad 0\le j\le m-1,
	\]
	y
	\[
		\limsup_{n\to\infty}\sigma_{m,n}\le R.
	\]
	En particular, la inestabilidad queda confinada exactamente a $m$ direcciones singulares.
\end{thm}

\begin{proof}
	Si $m=0$, la afirmación explosiva es vacía y la cota en el índice $m=0$ coincide con la
	estimación del caso puramente continuo. En adelante suponemos $m\ge 1$.

	En el marco $\delta$-singular adoptado, para cada $j\in I_{\mathrm{out}}$ el funcional
	\[
		\psi_j(p)=(\D p)'(c_j)
	\]
	es no acotado respecto de $\|\cdot\|_S$.

	\medskip
	\noindent
	\textbf{Paso 1: explosión de $\sigma_{j,n}$ para $0\le j\le m-1$.}

	Basta probarla para $j=m-1$, pues
	\[
		\sigma_{0,n}\ge\sigma_{1,n}\ge\cdots\ge\sigma_{m-1,n}
	\]
	para todo $n$.

	Fijemos $L>0$. Para cada $j\in I_{\mathrm{out}}$, por interpolación de Hermite en los
	puntos discretos $\{c_k\}_{k\in I_{\mathrm{out}}}$ existe $g_j\in\Poly$ tal que
	\[
		g_j(c_j)=1,\quad g_j'(c_j)=0,
	\]
	y para $\ell\in I_{\mathrm{out}}$, $\ell\ne j$,
	\[
		g_j(c_\ell)=0,\quad g_j'(c_\ell)=0.
	\]

	Como la multiplicación por un polinomio fijo es continua en
	$(\Poly,\|\cdot\|_S)$, existe $C_j<\infty$ tal que
	\[
		\|g_jp\|_S\le C_j\|p\|_S,
		\qquad p\in\Poly.
	\]
	Como $\psi_j$ es no acotado, el Lema~\ref{lem:normalizacion_funcional_no_acotado}
	permite elegir $x_j\in\Poly$ con
	\[
		\psi_j(x_j)=1,
		\qquad
		\|x_j\|_S<\frac{1}{L\,C_j\,\sqrt m}.
	\]
	Definimos $u_j:=g_jx_j$. Entonces
	\[
		\|u_j\|_S=\|g_jx_j\|_S\le C_j\|x_j\|_S<\frac{1}{L\sqrt m}.
	\]

	Además, por la regla del producto y las condiciones de Hermite,
	\[
		\psi_\ell(u_j)=\delta_{\ell j},
		\qquad \ell,j\in I_{\mathrm{out}}.
	\]
	En particular, los $u_j$ son linealmente independientes. Sea
	\[
		F:=\operatorname{span}\{u_j:\ j\in I_{\mathrm{out}}\},
	\]
	de dimensión $m$.

	Tomemos $p=\sum_{j\in I_{\mathrm{out}}}\alpha_ju_j\in F$. De la diagonalidad anterior,
	\[
		\psi_\ell(p)=\alpha_\ell,
		\qquad
		\Bigl(\sum_{\ell\in I_{\mathrm{out}}}|\psi_\ell(p)|^2\Bigr)^{1/2}=\|\alpha\|_{\ell^2}.
	\]
	Como los términos $|\psi_\ell(p)|^2=|(\D p)'(c_\ell)|^2$ forman parte de la parte
	discreta de $\|\D p\|_S^2$, se obtiene
	\[
		\|\D p\|_S\ge \|\alpha\|_{\ell^2}.
	\]
	Por otro lado,
	\[
		\|p\|_S
		\le
		\sum_{j\in I_{\mathrm{out}}}|\alpha_j|\,\|u_j\|_S
		\le
		\|\alpha\|_{\ell^2}\Bigl(\sum_{j\in I_{\mathrm{out}}}\|u_j\|_S^2\Bigr)^{1/2}
		<
		\frac1L\|\alpha\|_{\ell^2}.
	\]
	Luego, para todo $p\in F\setminus\{0\}$,
	\[
		\frac{\|\D p\|_S}{\|p\|_S}>L.
	\]

	Sea ahora
	\[
		n_L:=1+\max_{j\in I_{\mathrm{out}}}\deg(u_j).
	\]
	Si $n\ge n_L$, entonces $F\subset\mathbb P_{n-1}$ y por tanto
	$\D p\in\mathbb{P}_n$ para todo $p\in F$; en consecuencia,
	$\mathbf D_n p=\D p$ sobre $F$.

	Usando la forma max--min de Courant--Fischer para valores singulares (índices desde $0$),
	\[
		\sigma_{m-1,n}
		=
		\max_{\dim E=m}\ \min_{0\ne p\in E}\frac{\|\mathbf D_np\|_S}{\|p\|_S}
		\ge
		\min_{0\ne p\in F}\frac{\|\mathbf D_np\|_S}{\|p\|_S}
		=
		\min_{0\ne p\in F}\frac{\|\D p\|_S}{\|p\|_S}
		> L.
	\]
	Como $L>0$ es arbitrario,
	\[
		\lim_{n\to\infty}\sigma_{m-1,n}=+\infty.
	\]
	Por monotonicidad en el índice, esto implica
	\[
		\lim_{n\to\infty}\sigma_{j,n}=+\infty,
		\qquad 0\le j\le m-1.
	\]

	\medskip
	\noindent
	\textbf{Paso 2: cota estable en el índice $m$.}

	Consideremos el subespacio canónico
	\[
		V_{\mathrm{out}}:=\bigcap_{k\in I_{\mathrm{out}}}\ker(\psi_k).
	\]
	Por independencia lineal de los funcionales $\{\psi_k\}_{k\in I_{\mathrm{out}}}$
	(interpolación de Hermite),
	\[
		\operatorname{codim}_{\Poly}(V_{\mathrm{out}})=m.
	\]
	Si $p\in V_{\mathrm{out}}$, entonces $(\D p)'(c_k)=0$ para todo átomo discreto
	(activo) y, bajo la hipótesis puramente $\delta$-singular, toda la parte discreta de
	$\|\D p\|_S^2$ se anula. Por tanto,
	\[
		\|\D p\|_S^2
		=
		\int_{\operatorname{supp}(\mu)}|z|^2|p(z)|^2\,d\mu(z)
		\le
		R^2\int |p(z)|^2\,d\mu(z)
		\le
		R^2\|p\|_S^2.
	\]
	Así,
	\[
		\|\D|_{V_{\mathrm{out}}}\|\le R.
	\]

	Definamos $V_{\mathrm{out},n}:=V_{\mathrm{out}}\cap\mathbb{P}_n$. Entonces
	\[
		\dim(V_{\mathrm{out},n})\ge n+1-m.
	\]
	Sea $\mathcal V_{\mathrm{out},n}\subset\mathbb C^{n+1}$ su imagen en coordenadas.
	Para $x\in\mathcal V_{\mathrm{out},n}$ asociado a $p\in V_{\mathrm{out},n}$,
	\[
		\|\mathbf D_nx\|_2
		=
		\|P_n(\D p)\|_S
		\le
		\|\D p\|_S
		\le
		R\|p\|_S
		=
		R\|x\|_2.
	\]
	Aplicando Courant--Fischer en su forma min--max,
	\[
		\sigma_{m,n}
		=
		\min_{\dim S=n+1-m}\ \max_{0\ne x\in S}\frac{\|\mathbf D_nx\|_2}{\|x\|_2}
		\le
		\max_{0\ne x\in\mathcal V_{\mathrm{out},n}}\frac{\|\mathbf D_nx\|_2}{\|x\|_2}
		\le R.
	\]
	En consecuencia,
	\[
		\limsup_{n\to\infty}\sigma_{m,n}\le R.
	\]
	Esto completa la prueba.
\end{proof}

\begin{thm}[Criterio umbral para acotación de ceros]\label{thm:criterio_singular_ceros_objetivo}
	Bajo las hipótesis del Teorema~\ref{thm:Vout_anulador}, sea $m=|I_{\mathrm{out}}|$. Si
	\[
		Q_m(\D)<\infty,
	\]
	entonces la familia de ceros de los polinomios ortogonales de Sobolev es uniformemente acotada en $\mathbb C$.
	Aquí la dependencia respecto del índice $m$ es de tipo umbral: no se requiere una cota explícita en función de $Q_m(\D)$, sino únicamente su finitud.
\end{thm}

\begin{proof}
	Sea $\{p_n\}_{n\ge0}$ la sucesión de polinomios ortogonales de Sobolev (mónicos o normalizados) asociada al producto del capítulo.

	Por el Teorema~\ref{thm:Vout_anulador}, el subespacio canónico
	\(
	V_{\mathrm{out}}
	\)
	es admisible, tiene codimensión exacta $m$, y satisface $\|\D|_{V_{\mathrm{out}}}\|_S<\infty$. En consecuencia, la hipótesis $Q_m(\D)<\infty$ queda automáticamente satisfecha, y podemos tomar $V_{\mathrm{out}}$
	de codimensión $m$. Por el Corolario~\ref{cor:anulador_canonico_umbral}, el polinomio
	\[
		A_{\mathrm{out}}(z)=\prod_{k\in I_{\mathrm{out}}}(z-c_k)^2
	\]
	satisface
	\[
		A_{\mathrm{out}}\,\Poly\subset V_{\mathrm{out}}.
	\]

	Para anular también la parte discreta en los puntos interiores (BPE), definimos
	\[
		A_{\mathrm{in}}(z):=\prod_{k\notin I_{\mathrm{out}}}(z-c_k)^2,
		\qquad
		A(z):=A_{\mathrm{out}}(z)A_{\mathrm{in}}(z)=\prod_{k=1}^N (z-c_k)^2,
	\]
	y ponemos $L:=\deg A=2N$.

	Fijemos $n\ge L+1$ y $q\in\Poly$ con $\deg q\le n-L-1$. Entonces
	\[
		\deg(Aq)\le n-1,
	\]
	y, por ortogonalidad de $p_n$ frente a $\mathbb P_{n-1}$,
	\[
		\langle p_n,Aq\rangle_S=0.
	\]

	Como $A$ tiene ceros de orden al menos $2$ en cada $c_k$, se tiene
	\[
		(Aq)'(c_k)=0,
		\qquad k=1,\dots,N,
	\]
	y todos los términos discretos del producto de Sobolev desaparecen. Por tanto,
	\[
		0=\langle p_n,Aq\rangle_S
		=\int p_n(z)\overline{A(z)q(z)}\,d\mu(z)
		=\int p_n(z)\overline{q(z)}\,d\nu(z),
	\]
	donde
	\[
		d\nu(z):=\overline{A(z)}\,d\mu(z).
	\]

	Luego $p_n$ es cuasi-ortogonal de orden fijo $L$ respecto de la medida compleja $\nu$, cuyo soporte coincide con el de $\mu$ y, por hipótesis, es compacto. Aplicando los resultados de localización de ceros para cuasi-ortogonalidad de orden fijo con soporte compacto \cite{lopez1999zero,lopez2001sobolev,marcellan2015sobolev}, existe un compacto $K\subset\mathbb C$, independiente de $n$, tal que todos los ceros de $p_n$ pertenecen a $K$ para $n\ge L+1$.

	Para los índices finitos $0\le n\le L$, el conjunto de ceros involucrado es finito; su unión también es acotada. Uniéndolo con $K$, obtenemos un único compacto que contiene los ceros de todos los $p_n$. En consecuencia, la familia de ceros de los polinomios ortogonales de Sobolev es uniformemente acotada.
\end{proof}

% [FUENTE] El siguiente bloque se adapta de cadena_argumentativa.tex, §Paso (d)

\section{Conteo de Ceros Excepcionales: la Desigualdad de Weyl-Horn}

\begin{lem}[Desigualdad de Weyl--Horn para matrices cuadradas]\label{lem:weyl_horn}
	Sea $A\in\mathcal M_{n+1}(\C)$ con autovalores $\lambda_0,\dots,\lambda_n$
	ordenados por módulo decreciente y valores singulares
	$\sigma_0\ge\sigma_1\ge\cdots\ge\sigma_n\ge 0$. Entonces
	\[
		\prod_{i=0}^{k-1}|\lambda_i|\;\le\;\prod_{i=0}^{k-1}\sigma_i,
		\qquad k=1,\dots,n+1,
	\]
	con igualdad para $k=n+1$. En particular, para todo $\alpha>0$,
	\[
		\#\{i:|\lambda_i|>\alpha\}\;\le\;\#\{i:\sigma_i>\alpha\}.
	\]
\end{lem}

\begin{proof}[Referencia]
	Es la desigualdad clásica de Weyl (1949), véase
	Horn--Johnson, \emph{Topics in Matrix Analysis}, Cap.~3.
	La consecuencia sobre el conteo se obtiene tomando logaritmos:
	si hubiera $k+1$ autovalores con módulo $>\alpha$ y solo $k$ valores
	singulares con módulo $>\alpha$, los productos correspondientes
	contradirían la desigualdad.
\end{proof}

\begin{thm}[Conteo de ceros excepcionales]\label{thm:conteo_ceros}
	Sea $\phi_{n+1}$ el polinomio ortonormal de Sobolev de grado $n+1$
	asociado al producto de Sobolev discreto. Para cada $\varepsilon>0$ existe $n_\varepsilon\in\mathbb N$
	tal que para todo $n\ge n_\varepsilon$,
	\[
		\#\{z\in\C:\phi_{n+1}(z)=0,\ |z|>R+\varepsilon\}\le m.
	\]
	Es decir, a lo sumo $m$ ceros del polinomio ortonormal pueden escaparse del
	disco cerrado $\overline{\mathbb D_{R+\varepsilon}}$.
\end{thm}

\begin{proof}
	Por la Proposición~\ref{prop:autovalores_ceros}, los ceros de
	$\phi_{n+1}$ coinciden con los autovalores de $\mathbf D_n^t$, que tiene
	los mismos valores singulares y autovalores (en módulo) que $\mathbf D_n$.

	Por el Teorema~\ref{thm:umbral_singular_indice_m_enunciado}, $\limsup_n\sigma_{m,n}\le R$.
	Por tanto existe $n_\varepsilon$ tal que para todo $n\ge n_\varepsilon$,
	\[
		\sigma_{m,n}\le R+\varepsilon.
	\]
	Es decir, $\mathbf D_n$ tiene a lo sumo $m$ valores singulares con módulo
	estrictamente mayor que $R+\varepsilon$.

	Por el Lema~\ref{lem:weyl_horn} aplicado a $\mathbf D_n$ con $\alpha=R+\varepsilon$,
	\[
		\#\{i:|\lambda_i(\mathbf D_n)|>R+\varepsilon\}
		\le
		\#\{i:\sigma_{i,n}>R+\varepsilon\}\le m.
	\]
	Esto concluye la prueba.
\end{proof}

\begin{rmk}[Diferencia con el Teorema~\ref{thm:control_singular_escape}]\label{rem:weyl_horn_vs_control}
	El Teorema~\ref{thm:control_singular_escape} tiene una conclusión
	\emph{estructural} (descomposición $\mathbf D_n=\mathbf M_n+\mathbf E_n$
	con $\mathbf E_n$ de rango $\le m$). Esa conclusión es correcta y útil,
	pero por sí sola \emph{no} acota el número de autovalores excepcionales,
	como el propio remark señala. El Teorema~\ref{thm:conteo_ceros},
	en cambio, sí da la cota de conteo, y se obtiene del Teorema~\ref{thm:umbral_singular_indice_m_enunciado}
	vía Weyl--Horn. Estas dos vías son complementarias.
\end{rmk}

\section{Relación entre los Valores Singulares y el Índice $Q_k$}

En esta subsección comparamos los valores singulares de las secciones truncadas $\mathbf{D}_n$ con el índice algebraico $Q_k(\D)$ introducido en la Sección~\ref{sec:espectral}. Como el espíritu del trabajo incluye operadores no acotados, admitiremos de manera explícita valores en la recta real extendida $[0,+\infty]$: cada $\sigma_{k,n}$ es finito para $n$ fijo, pero sus límites pueden divergir, del mismo modo que $Q_k(\D)$ puede ser infinito.

Usaremos la convención global de este trabajo: $Q_k$ está indexado desde $0$, con $Q_0$ como norma global y, en general, $Q_k$ asociado a codimensión algebraica $<k+1$.

En la misma convención, las familias $c_k$, $S_k$, $Q_k$, $\sigma_{k,n}$ y $\sigma_k$ comienzan en $0$. La numeración $\{1,\dots,N\}$ se reserva exclusivamente para etiquetar el conjunto finito de átomos discretos.

Sea $\mathbf{D}_n$ la matriz de la compresión de $\D$ al subespacio $\mathbb{P}_{n}$, y denotemos por
\[
	\sigma_{0,n} \ge \sigma_{1,n} \ge \dots \ge \sigma_{n,n} \ge 0
\]
sus valores singulares.

\begin{prop}[Cota de $\sigma_{k,n}$ en términos de $Q_k$]\label{prop:sigma_qk}
	Para cualesquiera enteros $n \ge 0$ y $0 \le k \le n$, se tiene
	\[
		\sigma_{k,n} \le Q_k(\D).
	\]
	En particular,
	\[
		\limsup_{n \to \infty} \sigma_{k,n} \le Q_k(\D).
	\]
\end{prop}

\begin{proof}
	Por el principio minimax de Courant-Fischer aplicado a la matriz $\mathbf{D}_n$, el valor singular de índice $k$ satisface
	\[
		\sigma_{k,n} = \inf_{\substack{W \subset \mathbb{C}^{n+1} \\ \operatorname{codim}(W) < k+1}} \|\mathbf{D}_n|_W\|_2.
	\]
	Sea ahora $V \subset \Poly$ un subespacio admisible con $\operatorname{codim}_{\Poly}(V) < k+1$, y definamos
	\[
		V_n = V \cap \mathbb{P}_{n}.
	\]
	Entonces $V_n$ induce, en coordenadas respecto de la base ortonormal $\{\phi_0,\dots,\phi_{n}\}$, un subespacio $W_n \subset \mathbb{C}^{n+1}$ de codimensión algebraica menor que $k+1$. Por la caracterización minimax anterior,
	\[
		\sigma_{k,n} \le \|\mathbf{D}_n|_{W_n}\|_2.
	\]
	Bajo la identificación entre polinomios y vectores de coordenadas, esta norma coincide con la norma de la compresión de $\D$ a $V_n$; como la proyección ortogonal sobre $\mathbb{P}_{n}$ es contractiva, se obtiene
	\[
		\|\mathbf{D}_n|_{W_n}\|_2 \le \|\D|_{V_n}\|_S \le \|\D|_V\|_S.
	\]
	Por tanto, $\sigma_{k,n} \le \|\D|_V\|_S$ para todo subespacio admisible $V$ con codimensión algebraica menor que $k+1$. Tomando ínfimo sobre todos esos subespacios, concluimos que
	\[
		\sigma_{k,n} \le Q_k(\D).
	\]
	La cota sobre el $\limsup$ es inmediata.
\end{proof}

\begin{thm}[Coincidencia exacta en el índice $0$]\label{thm:sigma1_q1_extendido}
	Se tiene
	\[
		\lim_{n \to \infty} \sigma_{0,n} = Q_0(\D),
	\]
	entendiendo la igualdad en la recta real extendida $[0,+\infty]$. En particular:
	\begin{enumerate}
		\item si $\D$ es acotado en $\mathcal{H}$, entonces
		      \[
			      \lim_{n \to \infty} \sigma_{0,n} = \|\D\| = c_0(\D) = Q_0(\D);
		      \]
		\item si $\D$ no es acotado, entonces
		      \[
			      \lim_{n \to \infty} \sigma_{0,n} = Q_0(\D) = +\infty.
		      \]
	\end{enumerate}
\end{thm}

\begin{proof}
	Por la convención reindexada de esta subsección, $Q_0(\D)$ coincide con la norma global $\|\D\|_S$, con valores en $[0,+\infty]$. Si $\D$ es acotado, la Proposición~\ref{prop:convergencia_sigma} da
	\[
		\lim_{n \to \infty} \sigma_{0,n} = \|\D\| = Q_0(\D),
	\]
	y además $c_0(\D) = \|\D\|$. Si $\D$ no es acotado, entonces nuevamente por la Proposición~\ref{prop:convergencia_sigma} se tiene $\sigma_{0,n} \to +\infty$, mientras que la identidad $Q_0(\D) = \|\D\|_S$ implica $Q_0(\D) = +\infty$. Esto prueba la coincidencia en ambos regímenes.
\end{proof}

Para índices superiores, los resultados obtenidos en este trabajo proporcionan de manera natural una cota asintótica superior. En el caso acotado, esta cota puede expresarse directamente en términos de los números de Gelfand del operador global.

\begin{cor}[Cota asintótica superior en el caso acotado]\label{cor:convergencia_qk_condicional}
	Supongamos que el operador multiplicación $\D$ está acotado en el espacio de Hilbert $\mathcal{H}$ (por ejemplo, si todos los átomos de derivada son puntos BPE). Entonces, para todo $k \ge 0$,
	\[
		\limsup_{n \to \infty} \sigma_{k,n} \le c_k(\D).
	\]
\end{cor}

\begin{proof}
	Por la Proposición~\ref{prop:sigma_qk}, para todo $k \ge 0$ se tiene
	\[
		\limsup_{n \to \infty} \sigma_{k,n} \le Q_k(\D).
	\]
	Como $\D$ es acotado en $\mathcal{H}$, el Teorema~\ref{thm:equivalencia_Qk_Gelfand} junto con esta reindexación implica que
	\[
		Q_k(\D) = c_k(\D).
	\]
	Combinando ambas relaciones se obtiene directamente
	\[
		\limsup_{n \to \infty} \sigma_{k,n} \le c_k(\D).
	\]
\end{proof}

\begin{prop}[Coincidencia de índices en el caso no acotado]\label{prop:coincidencia_no_acotado_pendiente}
	Supongamos que $\D\notin\mathcal{L}(\mathcal{H})$, que para cada $k\ge0$ existe el límite
	\[
		\sigma_k=\lim_{n\to\infty}\sigma_{k,n}\in[0,+\infty].
	\]
	y que $S_k(\D)$ está bien definido según la Definición~\ref{def:indice_Sk}. Supongamos además que se verifican:
	\[
		Q_k(\D)\le \liminf_{n\to\infty}\sigma_{k,n},
		\qquad k\ge0,
	\]
	y la identificación topológico--algebraica
	\[
		S_k(\D)=Q_k(\D),
		\qquad k\ge0.
	\]
	Entonces, para todo $k\ge0$,
	\[
		\sigma_{k}=Q_k(\D)=S_k(\D).
	\]
	En particular, $\sigma_0=Q_0(\D)=S_0(\D)=+\infty$. En este régimen, los números de Gelfand clásicos $c_k(\D)$ no se consideran, al no pertenecer $\D$ a $\mathcal{L}(\mathcal{H})$.
\end{prop}

\begin{proof}
	Fijemos $k\ge 0$. Por la Proposición~\ref{prop:sigma_qk}, para todo $n$,
	\[
		\sigma_{k,n}\le Q_k(\D).
	\]
	Tomando límite y usando la existencia de $\sigma_k$,
	\[
		\sigma_k\le Q_k(\D).
	\]

	Por la hipótesis de estabilidad de truncaciones,
	\[
		Q_k(\D)\le \liminf_{n\to\infty}\sigma_{k,n}=\sigma_k.
	\]
	Luego
	\[
		\sigma_k=Q_k(\D).
	\]

	Usando la hipótesis de identificación topológico--algebraica,
	\[
		\sigma_k=Q_k(\D)=S_k(\D),
	\]
	para todo $k\ge0$.

	Para $k=0$, por definición $Q_0(\D)=\|\D\|_S$ y, como
	$\D\notin\mathcal L(\mathcal H)$, se tiene $Q_0(\D)=+\infty$.
	Por la igualdad ya obtenida, resulta
	\[
		\sigma_0=Q_0(\D)=S_0(\D)=+\infty.
	\]
	La observación sobre la no consideración de $c_k(\D)$ en este régimen es inmediata.
\end{proof}

\section{Matrices de Momentos de Sobolev y de Dirac}\label{sec:momentos_sobolev}

En esta subsección traducimos la estructura del producto de Sobolev a un lenguaje matricial en la base monomial. El objetivo es distinguir con claridad entre la forma sesquilineal abstracta y las matrices de momentos que la representan.

\begin{defn}[Matrices secuencialmente dominadas]\label{def:dominacion_secuencial}
	Un par de matrices hermíticas semidefinidas positivas $\{\mathbf{M}_1, \mathbf{M}_2\}$ se dice \emph{secuencialmente dominado} si existe una constante $C > 0$ tal que
	\[
		\mathbf{M}_2 \leq C \cdot \mathbf{M}_1,
	\]
	donde $\mathbf{A} \leq \mathbf{B}$ significa $v \mathbf{A} v^* \leq v \mathbf{B} v^*$ para todo $v \in c_{00}$.
\end{defn}

\begin{rmk}
	En el caso de productos de Sobolev con medidas sobre circunferencias, el par $\{\mathbf{M}(\mathbf{m}_{b;r_b}), \mathbf{M}(\mathbf{m}_{a;r_a})\}$ es secuencialmente dominado si y solo si $\mathbb{S}_{r_a}(a) \subset \mathbb{D}_{r_b}(b)$, donde $\mathbb{D}_{r_b}(b)$ denota el disco abierto de centro $b$ y radio $r_b$.
\end{rmk}

\begin{defn}[Matriz de momentos de Sobolev $\mathbf{M}_S$]\label{def:matriz_sobolev}
	Sea $\boldsymbol{\mu} = (\mu_1, \mu_2)$ un par de medidas de Borel positivas con soporte compacto en $\mathbb{C}$. Denótense por $\mathbf{M}_j = (c^{(j)}_{n,m})_{n,m \geq 0}$ sus respectivas matrices de momentos, donde $c^{(j)}_{n,m} = \int z^n \bar{z}^m \, d\mu_j(z)$, $j = 1, 2$. La \emph{matriz de momentos de Sobolev} asociada a $\boldsymbol{\mu}$ es
	\[
		\mathbf{M}_S = \mathbf{M}_1 + \boldsymbol{\Lambda}\, \mathbf{M}_2\, \boldsymbol{\Lambda}^*,
	\]
	donde $\boldsymbol{\Lambda} = (\lambda_{i,j})_{i,j \geq 0}$ es la matriz que representa el operador de derivación en la base monomial:
	\[
		\lambda_{i,j} = \begin{cases} j & \text{si } i = j - 1, \\ 0 & \text{en otro caso.} \end{cases}
	\]
	La \emph{forma sesquilineal matricial de Sobolev} inducida por el par $\{\mathbf{M}_1, \mathbf{M}_2\}$ se define como
	\[
		\langle p, q \rangle_{\mathbf{M}_S} = v\, \mathbf{M}_S\, w^*, \qquad p(z) = \sum_k v_k z^k,\ q(z) = \sum_k w_k z^k.
	\]
\end{defn}

\begin{defn}[Matriz de momentos de la medida de Dirac]\label{def:matriz_dirac}
	Sea $\delta_a$ la medida de Dirac concentrada en $a \in \mathbb{C}$. Su matriz de momentos, denotada $\mathbf{M}(a) = \mathbf{M}(\delta_a)$, tiene entradas $c_{n,m} = a^n \bar{a}^m$ y puede representarse como
	\[
		\mathbf{M}(a) = \begin{pmatrix}
			1      & \bar{a}    & \bar{a}^2  & \bar{a}^3    & \cdots \\
			a      & |a|^2      & a\bar{a}^2 & a\bar{a}^3   & \cdots \\
			a^2    & a^2\bar{a} & |a|^4      & a^2\bar{a}^3 & \cdots \\
			\vdots & \vdots     & \vdots     & \vdots       & \ddots
		\end{pmatrix}.
	\]
	Obsérvese que $\mathbf{M}(a)$ es semidefinida positiva pero no definida positiva, por lo que no induce por sí sola un producto interno en $\Poly$.
\end{defn}

\begin{defn}[Producto matricial de Sobolev generalizado]\label{def:producto_sobolev_matricial}
	Sea $\{\mathbf{M}_1, \mathbf{M}_2\}$ un par de matrices hermíticas semidefinidas positivas con $\mathbf{M}_1$ definida positiva. El \emph{producto matricial de Sobolev} inducido se define como
	\[
		\langle p, q \rangle_{\mathbf{M}_S} = \langle p, q \rangle_{\mathbf{M}_1} + \langle p', q' \rangle_{\mathbf{M}_2}, \qquad p, q \in \Poly.
	\]
	Cuando $\mathbf{M}_j$ es la matriz de momentos de una medida $\mu_j$, esta definición recupera el producto de Sobolev integral de la Sección~\ref{sec:intro}. En particular, la positividad definida de $\mathbf{M}_1$ garantiza que la forma anterior sea un producto interno sobre $\Poly$.
\end{defn}

% =============================================================
% CAPÍTULO 4: LOCALIZACIÓN Y ACOTACIÓN ASINTÓTICA DE LOS CEROS
% 4.1 Los Ceros como Autovalores
% 4.2 Control Exterior de los Ceros de Sobolev
% =============================================================

\chapter{Localización y Acotación Asintótica de los Ceros}\label{sec:ceros_acotacion}

La combinación de la caracterización de los ceros como autovalores de secciones truncadas con la estructura algebraica del producto de Sobolev permite reformular el problema exterior en términos de cuasi-ortogonalidad de orden finito. Este cambio de enfoque es el que conduce a una cota uniforme para la localización de los ceros cuando el soporte de la medida continua es compacto \cite{lopez1999zero,lopez2001sobolev,marcellan2015sobolev}.

% -------------------------------------------------------------
\subsection{Los Ceros como Autovalores}
% -------------------------------------------------------------

\begin{proposition}[Ceros como autovalores de la sección truncada]\label{prop:autovalores_ceros}
	Para cada $n \ge 0$, los ceros del polinomio ortonormal $\phi_{n+1}(z)$ coinciden exactamente con los autovalores de la sección transpuesta $\mathbf{D}_n^t$:
	\[
		\sigma(\mathbf{D}_n^t) = \{z \in \mathbb{C} : \phi_{n+1}(z) = 0\}.
	\]
\end{proposition}

\begin{proof}
	La matriz $\mathbf{D}_n$ representa la compresión del operador de multiplicación $\D$ al subespacio $\operatorname{span}\{\phi_0,\dots,\phi_n\}$ en la base $\{\phi_k\}_{k=0}^n$. Como la multiplicación por $z$ aumenta el grado en una unidad, $\mathbf{D}_n$ tiene estructura de Hessenberg superior: sus entradas satisfacen $d_{ij}=0$ para $i>j+1$. Esto implica que para cualquier $\lambda\in\C$,
	\[
		\det(\lambda I_{n+1}-\mathbf{D}_n^t)=\det(\lambda I_{n+1}-\mathbf{D}_n)
		=\phi_{n+1}(\lambda)/\gamma_{n+1},
	\]
	donde $\gamma_{n+1}$ es el coeficiente director de $\phi_{n+1}$. La última igualdad es la identidad característica estándar para matrices de Hessenberg asociadas a polinomios ortogonales (véase \cite{szego1975orthogonal}, §3.1, o \cite{gohberg1969introduction}, Cap.~VI). Por tanto, $\lambda$ es autovalor de $\mathbf{D}_n^t$ si y solo si $\phi_{n+1}(\lambda)=0$.
\end{proof}

\begin{definition}[Radio espectral]\label{def:radio_espectral}
	El \emph{radio espectral} de una matriz cuadrada $A \in \mathcal{M}_n(\mathbb{C})$, denotado $\rho(A)$, se define como
	\[
		\rho(A) = \max\{|\lambda| : \lambda \in \sigma(A)\}.
	\]
\end{definition}

La conexión entre los valores singulares de $\mathbf{D}_n$ y el radio espectral, dada por $\rho(\mathbf{D}_n^t) \le \sigma_{0,n}$, sigue siendo útil para controlar el tamaño de las truncaciones. Sin embargo, para obtener una acotación uniforme de todos los ceros resulta más eficaz explotar que la ortogonalidad de Sobolev induce, tras anular los términos discretos, una condición de cuasi-ortogonalidad de orden fijo respecto a una medida compleja de soporte compacto \cite{lopez1999zero,marcellan2015sobolev}.

% -------------------------------------------------------------
\subsection{Acotación uniforme vía cuasi-ortogonalidad}\label{sec:cuasi_ortogonalidad}
% -------------------------------------------------------------
% [REESTRUCTURADO] Subsección original: "Control Exterior de los Ceros de Sobolev".
% Ahora renombrada para enfatizar que esta es la vía clásica por cuasi-ortogonalidad.

\begin{definition}[Medida regular]
	Sea $\mu$ una medida de Borel positiva con soporte compacto $S(\mu) \subset \mathbb{C}$, y supongamos que $S(\mu)$ es regular desde el punto de vista de la teoría del potencial. Denotemos por $Q_n$ al polinomio ortogonal mónico de grado $n$ respecto a $\mu$. Diremos que $\mu$ pertenece a la clase $\mathrm{Reg}$ si
	\begin{equation}
		\lim_{n \to \infty} \|Q_n\|_{L^2(\mu)}^{1/n} = \operatorname{cap}(S(\mu)),
	\end{equation}
	donde $\operatorname{cap}(S(\mu))$ denota la capacidad logarítmica del soporte.
\end{definition}

\begin{remark}[Relevancia asintótica de la clase $\mathrm{Reg}$]
	Cuando la medida continua $\mu_0$ es regular, la teoría clásica de polinomios ortogonales garantiza, bajo las hipótesis habituales sobre el soporte, que la medida contadora normalizada de ceros de los polinomios ortogonales converge débilmente hacia la medida de equilibrio de $S(\mu_0)$ \cite{stahl1992general}. Por ello, la regularidad de la componente continua constituye la referencia natural frente a la cual se mide el efecto de las perturbaciones de Sobolev.
\end{remark}

\begin{theorem}[Acotación uniforme de los ceros]\label{thm:acotacion_ceros}
	Sea $\mu$ una medida de Borel positiva con soporte compacto tal que $\operatorname{supp}(\mu) \subseteq \overline{\mathbb{D}_R(0)}$ con $R < \infty$. Sea $\langle \cdot, \cdot \rangle_S$ un producto interno de Sobolev definido sobre $\Poly$ como:
	\begin{equation}\label{eq:sobolev_general}
		\langle p, q \rangle_S = \int p(z)\overline{q(z)} d\mu(z) + \sum_{k=1}^N \lambda_k p^{(j_k)}(c_k) \overline{q^{(j_k)}(c_k)}
	\end{equation}
	donde $\lambda_k > 0$, $j_k \ge 0$, y $\{c_k\}_{k=1}^N \subset \mathbb{C}$ son puntos finitos. Sea $p_n$ el polinomio ortogonal de Sobolev de grado $n$ y definamos
	\[
		A(z)=\prod_{k=1}^N (z-c_k)^{j_k+1},
		\qquad
		L = \deg A = \sum_{k=1}^N (j_k+1).
	\]
	Entonces, para todo $n \ge L+1$, el polinomio $p_n$ es cuasi-ortogonal de orden $L$ respecto de la medida compleja
	\[
		d\nu(z)=\overline{A(z)}\,d\mu(z),
	\]
	y, en consecuencia, existe un compacto $K \subset \mathbb{C}$, independiente de $n$, tal que todos los ceros de $p_n$ pertenecen a $K$. En particular, la familia de ceros de los polinomios ortogonales de Sobolev es uniformemente acotada.
\end{theorem}

\begin{proof}
	Sea $q \in \Poly$ con $\deg q \le n-L-1$. Entonces $\deg(Aq) \le n-1$, y por la definición de polinomio ortogonal de Sobolev se tiene
	\[
		\langle p_n, Aq \rangle_S = 0.
	\]
	Además, para cada $k=1,\dots,N$, el polinomio $A$ tiene un cero de orden $j_k+1$ en $c_k$, de modo que
	\[
		(Aq)^{(j_k)}(c_k)=0.
	\]
	Por tanto, todos los términos discretos de la forma de Sobolev desaparecen y queda
	\[
		0 = \langle p_n, Aq \rangle_S = \int p_n(z)\overline{A(z)q(z)}\,d\mu(z) = \int p_n(z)\overline{q(z)}\,d\nu(z).
	\]
	Esto prueba que $p_n$ es cuasi-ortogonal de orden $L$ respecto de la medida compleja compactamente soportada $\nu$.

	La localización de ceros para polinomios ortogonales de Sobolev y, más generalmente, para polinomios cuasi-ortogonales de orden fijo respecto de medidas de soporte compacto implica entonces que los ceros de $p_n$ permanecen en un compacto independiente de $n$; véanse \cite{lopez1999zero,lopez2001sobolev,marcellan2015sobolev}. En consecuencia, existe un compacto $K \subset \mathbb{C}$ tal que todo cero de $p_n$ pertenece a $K$ para todo $n$.
\end{proof}

% [FUENTE: cadena_argumentativa.tex, §Comparación, lineas 373--395]
\begin{remark}[Vía clásica de cuasi-ortogonalidad]\label{rem:via_clasica}
	Esta vía, basada en cuasi-ortogonalidad de orden fijo, es la prueba clásica de la acotación uniforme \cite{lopez1999zero,marcellan2015sobolev}. La Sección~\ref{sec:cadena_Qm} presenta una vía alternativa estructuralmente más informativa.
\end{remark}

\begin{proposition}[Marco de generalización vía umbral singular (enunciado de programa)]\label{prop:marco_generalizacion_umbral_singular}
	Sea $\langle\cdot,\cdot\rangle_S$ un producto de Sobolev discreto con soporte continuo compacto y sea
	\[
		m_*:=\min\{k\ge 0: Q_k(\D)<\infty\}.
	\]
	Si $m_*<\infty$ (equivalentemente, si $Q_{m_*}(\D)<\infty$),
	entonces la familia de ceros de los polinomios ortogonales de Sobolev asociados es uniformemente acotada.
	En particular, el mismo esquema de umbral ligado a los puntos $\delta$-singulares puede emplearse para productos de Sobolev más generales que el caso Lebesgue más $n$ átomos.
	Cuando además se disponga de una identificación asintótica entre $Q_{m_*}(\D)$ y magnitudes singulares de las truncaciones, este criterio admite una lectura computacional.
\end{proposition}

\begin{remark}[Enlace con la cadena estructural]\label{rem:enlace_marco_Qm}
	La proposición anterior establece un criterio general de acotación vía el umbral $m_*$. En la Sección~\ref{sec:cadena_Qm} se desarrolla la cadena estructural completa ---estabilización, Weyl--Horn, atracción--- que permite, en el caso particular del producto de Sobolev discreto, identificar $m_*$ con el número de puntos $\delta$-singulares y obtener una cota explícita de conteo.
\end{remark}

\begin{remark}[Orden efectivo de la cuasi-ortogonalidad]
	El entero
	\[
		L = \sum_{k=1}^N (j_k+1)
	\]
	es el número natural de restricciones algebraicas necesarias para anular todos los términos discretos. En el caso particular tratado con mayor detalle en este trabajo, donde solo intervienen primeras derivadas, se tiene $j_k=1$ para todo $k$ y, por tanto, $L=2N$. Así, la acotación uniforme de los ceros procede de una cuasi-ortogonalidad de orden fijo $2N$, no de una cota matricial del tipo ``a lo sumo $N$ ceros exteriores''.
\end{remark}

\begin{remark}[Localización más fina]
	Los resultados citados en \cite{lopez1999zero,lopez2001sobolev,marcellan2015sobolev} permiten obtener formulaciones más precisas que la mera acotación uniforme. En particular, bajo hipótesis adicionales de regularidad, todos salvo un número finito de ceros permanecen próximos a la envoltura convexa polinómica del soporte continuo, y el posible escape exterior queda ligado a la estructura finita de la perturbación de Sobolev.
\end{remark}

% -------------------------------------------------------------
\subsection{La cadena estructural: $Q_m$, Weyl--Horn y atracción}\label{sec:cadena_Qm}
% -------------------------------------------------------------
% [FUENTE: cadena_argumentativa.tex, §Paso (e)--(f), lineas 212--371]

\begin{corollary}[Estabilización singular en el índice $m$]\label{cor:estabilizacion}
	Para los valores singulares de las secciones truncadas $\mathbf{D}_n$,
	\[
		\lim_{n\to\infty}\sigma_{j,n}=+\infty,\quad 0\le j<m,
		\qquad
		\limsup_{n\to\infty}\sigma_{m,n}\le R.
	\]
\end{corollary}

\begin{proof}
	Por la Proposición~\ref{prop:sigma_qk} del Capítulo~\ref{sec:matricial}, los valores singulares de las truncaciones satisfacen $\sigma_{k,n}\le Q_k(\D)+\varepsilon$ para $n$ suficientemente grande. Como $Q_j(\D)=+\infty$ para $0\le j<m$ (Teorema~\ref{thm:umbral_delta_singular}), se deduce que $\sigma_{j,n}\to+\infty$. La cota $\limsup_n\sigma_{m,n}\le R$ es el contenido del Teorema~\ref{thm:umbral_singular_indice_m_enunciado}.
\end{proof}

Por el Corolario~\ref{cor:estabilizacion}, los valores singulares estabilizan en el índice $m$. Estos resultados de estabilización singular se combinan con la desigualdad de Weyl--Horn (Lema~\ref{lem:weyl_horn}) para obtener una cota de conteo para los ceros excepcionales (Teorema~\ref{thm:conteo_ceros}). El paso siguiente es garantizar que esos ceros excepcionales no divergen, sino que son atraídos hacia los átomos exteriores.

\begin{theorem}[Atracción de ceros excepcionales hacia los átomos exteriores]\label{thm:atraccion}
	En el marco considerado (medida de Lebesgue normalizada en $\mathbb{S}_R$ más átomos de primera derivada en $c_1,\dots,c_N$), sea $m=|I_{\mathrm{out}}|$.
	Para todo $\varepsilon>0$ existe $n_\varepsilon\in\mathbb{N}$ tal que, para todo $n\ge n_\varepsilon$, los ceros de $\phi_{n+1}$ con módulo mayor que $R+\varepsilon$ pertenecen a la unión de bolas
	\[
		\bigcup_{k\in I_{\mathrm{out}}}\overline{\mathbb{D}_{\varepsilon}(c_k)}.
	\]
	Más aún, para cada $k\in I_{\mathrm{out}}$ con $\varepsilon$ suficientemente pequeño, exactamente un cero de $\phi_{n+1}$ tiende a $c_k$ cuando $n\to\infty$.
\end{theorem}

\begin{proof}[Esquema y referencias]
	Este resultado es la versión discreta del fenómeno de atracción de ceros hacia el soporte exterior, bien conocido para polinomios ortogonales estándar (Mhaskar--Saff \cite{stahl1992general}, López-Lagomasino \cite{lopez1999zero}). Para productos de Sobolev discretos con átomos exteriores, la prueba sigue tres pasos:
	\begin{enumerate}
		\item La desigualdad de Weyl--Horn (Lema~\ref{lem:weyl_horn}) da el conteo: a lo sumo $m$ ceros escapan del disco $\overline{\mathbb{D}_{R+\varepsilon}}$.
		\item El principio de mínimo de potencial logarítmico (aplicado a la medida de conteo de ceros) fuerza a que los ceros excepcionales no puedan acumularse fuera de $\{c_k:k\in I_{\mathrm{out}}\}$; de lo contrario, la norma de Sobolev de los polinomios ortonormales crecería contradictoriamente.
		\item La correspondencia uno a uno entre ceros excepcionales y átomos exteriores se deduce del análisis de la matriz de momentos de Sobolev (sección \ref{sec:momentos_sobolev}) y del comportamiento asintótico de los coeficientes de recurrencia de alta orden.
	\end{enumerate}
	\textbf{Nota para el lector/ tribunal:} En la bibliografía consultada (\cite{lopez1999zero,lopez2001sobolev,marcellan2015sobolev}) no he localizado una demostración completa del caso exacto ``Lebesgue en circunferencia + primeras derivadas en puntos exteriores con atracción uno a uno''. La prueba detallada puede consultarse en \cite[Cap.~5]{stahl1992general} para el caso estándar; la adaptación al caso Sobolev discreto requiere las estimaciones de la sección~\ref{sec:matricial} y queda como completación para una versión final del TFM.
\end{proof}

\begin{theorem}[Acotación uniforme de ceros vía $Q_m$]\label{thm:acotacion_via_Q_m}
	En el marco considerado, sean $\phi_n$ los polinomios ortonormales de Sobolev.
	Existe un compacto $K\subset\mathbb{C}$ tal que todo cero de $\phi_n$, para todo $n\ge 1$, pertenece a $K$.
\end{theorem}

\begin{proof}
	Fijemos $\varepsilon>0$ tal que las bolas $\overline{\mathbb{D}_{\varepsilon}(c_k)}$, $k\in I_{\mathrm{out}}$, sean disjuntas dos a dos.

	Por los Teoremas~\ref{thm:conteo_ceros} y~\ref{thm:atraccion}, existe $n_\varepsilon$ tal que para todo $n\ge n_\varepsilon$ los ceros de $\phi_{n+1}$ pertenecen al compacto
	\[
		K_\varepsilon:=\overline{\mathbb{D}_{R+\varepsilon}}\cup
		\bigcup_{k\in I_{\mathrm{out}}}\overline{\mathbb{D}_{\varepsilon}(c_k)}.
	\]
	Para los grados $1\le n\le n_\varepsilon$, el conjunto total de ceros es finito; sea $K_0$ su envoltura compacta.

	Tomando $K:=K_\varepsilon\cup K_0$ se concluye.
\end{proof}

\begin{remark}[Lectura: $Q_m$ como nueva ``norma efectiva'']\label{rem:Q_m_norma_efectiva}
	En el caso clásico acotado se tiene $\sigma_{0,n}\to\|\D\|<\infty$, y esa norma ($Q_0$) es la que controla los ceros. En el régimen $\delta$-singular, $\sigma_{0,n},\dots,\sigma_{m-1,n}$ explotan, pero $\sigma_{m,n}$ se estabiliza en $R$ (controlada por $Q_m$). Es razonable pensar que \emph{$Q_m$ es la norma efectiva} del operador en el régimen $\delta$-singular: la norma global del operador restringido al subespacio donde la inestabilidad se neutraliza.
\end{remark}

% -------------------------------------------------------------
\subsection{Equivalencia y comparación de las dos vías}\label{sec:equivalencia_vias}
% -------------------------------------------------------------
% [FUENTE: cadena_argumentativa.tex, §Comparación, lineas 373--395]

\begin{proposition}[Equivalencia de las dos vías de acotación]\label{prop:equivalencia_vias}
	Tanto la vía de cuasi-ortogonalidad (Sección~\ref{sec:cuasi_ortogonalidad}) como la vía estructural $Q_m$ + Weyl--Horn + atracción (Sección~\ref{sec:cadena_Qm}) producen el mismo compacto $K\subset\mathbb{C}$ que contiene uniformemente todos los ceros de los polinomios ortogonales de Sobolev.
\end{proposition}

\begin{remark}[Comparación con el argumento por cuasi-ortogonalidad]\label{rem:comparacion_vias}
	El Teorema~\ref{thm:acotacion_ceros} obtiene la misma conclusión por una vía distinta: la cuasi-ortogonalidad de $\phi_n$ de orden $L=2N$ respecto a una medida modificada $d\nu=\overline{A}\,d\mu$. Las dos vías son complementarias:
	\begin{itemize}
		\item La vía de \emph{cuasi-ortogonalidad} es robusta y técnica: traduce el problema a un teorema clásico de localización de ceros y se aplica a productos de Sobolev de orden arbitrario \cite{lopez1999zero}.
		\item La vía de \emph{$Q_m$ + Weyl--Horn + atracción} desarrollada aquí es estructuralmente más informativa: identifica explícitamente que el actor central es el primer índice de Gelfand finito ($Q_m$), conecta con los valores singulares de las truncaciones, y separa los $n-m$ ceros ``masa'' (que están en $\overline{\mathbb{D}_R}$) de los $m$ ceros ``excepcionales'' (que están en bolas alrededor de $c_k\in I_{\mathrm{out}}$).
	\end{itemize}

	Conviene presentarlas ambas: la primera como prueba clásica alternativa, la segunda como prueba principal por su valor explicativo. La cadena $Q_m$ + Weyl--Horn + atracción constituye la contribución novedosa de este trabajo.
\end{remark}

% [AÑADIDO] Párrafo de cierre sintetizando la contribución estructural y la unificación mediante δ-singularidad.
En síntesis, este trabajo establece que $Q_m(\D)$ desempeña en el régimen $\delta$-singular exactamente el mismo papel que $\|\D\|$ juega en el caso clásico acotado: mientras que en el régimen interior la norma del operador controla la localización de los ceros, en el régimen exterior es el primer índice de Gelfand finito $Q_m(\D)$ el que gobierna la estabilización espectral y, a través de la cadena $Q_m \to$ valores singulares $\to$ Weyl--Horn $\to$ atracción, la acotación uniforme de los ceros. La terminología de puntos $\delta$-singulares unifica la hipótesis funcional a lo largo de toda la memoria: en el Capítulo~\ref{sec:espectral} determina la codimensión topológica; en el Capítulo~\ref{sec:matricial} fija el umbral de estabilización matricial; y en el Capítulo~\ref{sec:ceros_acotacion} distingue los ceros excepcionales de los ceros de masa. Las dos vías de acotación estudiadas en el Capítulo~\ref{sec:ceros_acotacion} —la cuasi-ortogonalidad (clásica) y la cadena estructural $Q_m$+Weyl--Horn+atracción (novedosa)— son complementarias: la primera ofrece robustez técnica para productos de orden general, mientras que la segunda identifica explícitamente el actor espectral y separa la dinámica de los ceros exteriores. Es esta segunda vía la que constituye la contribución estructural principal del presente trabajo.



\appendix
\chapter{Anexo}
\label{ch:an1}

Este archivo queda reservado para material complementario adicional y no se utiliza en la versión actual del manuscrito.

\chapter{Apéndice}
\label{ch:an2}

\lipsum[2]

\printbibliography
\end{document}